% Generated by the blueprint tool (blueprint-pdf/1) from a blueprint snapshot; edit the
% Markdown sources, not this file.  Compile with `latexmk -pdf` or by running
% `pdflatex` twice (the second run settles the contents and the page references).
% Images, if any, are read from assets/ next to this file.
\documentclass[11pt]{article}
\usepackage[T1]{fontenc}
\usepackage[utf8]{inputenc}
\usepackage{lmodern}
\usepackage{textcomp}
\usepackage[a4paper,margin=2.5cm,headheight=14pt]{geometry}
\usepackage{amsmath,amssymb,amsthm}
\usepackage{amscd}  % KaTeX's CD diagrams; part of amsmath, so always there
\usepackage{mathtools}
\usepackage{bm}
\usepackage{graphicx}
\usepackage{xcolor}
\usepackage{array,booktabs,tabularx}
\usepackage{fancyvrb}
\usepackage[normalem]{ulem}
\usepackage{microtype}
% Optional, and only with their fonts: a small TeX installation has the
% packages without the fonts, which fails at the first formula.
\IfFileExists{mathrsfs.sty}{\IfFileExists{rsfs10.tfm}{\usepackage{mathrsfs}}{}}{}
\IfFileExists{stmaryrd.sty}{\IfFileExists{stmary10.tfm}{\usepackage{stmaryrd}}{}}{}
\IfFileExists{cancel.sty}{\usepackage{cancel}}{}
\usepackage[hyphens]{url}
\usepackage[pdfusetitle,bookmarksnumbered=false]{hyperref}
\usepackage{bookmark}
\definecolor{bplink}{HTML}{1F4E9C}
\definecolor{bpmuted}{HTML}{5F676E}
\hypersetup{colorlinks=true,linkcolor=bplink,urlcolor=bplink,citecolor=bplink}
\setlength{\parindent}{0pt}
\setlength{\parskip}{0.55em plus 0.2em minus 0.1em}
% Lists as tight as the page's.
\makeatletter
\def\@listi{\leftmargin\leftmargini\parsep 0pt\topsep 0.3em plus 0.1em\itemsep 0.1em plus 0.1em}
\let\@listI\@listi
\def\@listii{\leftmargin\leftmarginii\labelwidth\leftmarginii\advance\labelwidth-\labelsep\parsep 0pt\topsep 0.1em\itemsep 0.1em}
% An ordered list that does not start at 1.
\newcommand{\bpstartat}[1]{\setcounter{\@enumctr}{\numexpr#1-1\relax}}
\makeatother
\setcounter{tocdepth}{3}
\emergencystretch=3em
\allowdisplaybreaks

% A numbered statement, set as a paper sets one: "Theorem 2.1 (Title)."
% #1 is the lead ("Theorem~2.1"), #2 what follows it in roman.
\newenvironment{bpstatement}[2]{\par\addvspace{0.6\baselineskip}\noindent\textbf{#1}#2\textbf{.}\enspace\ignorespaces}
  {\par\addvspace{0.4\baselineskip}}
\newenvironment{bpstatementit}[2]{\begin{bpstatement}{#1}{#2}\itshape\ignorespaces}{\end{bpstatement}}
% A step: a relation with prose of its own, read under the object it starts from.
\newenvironment{bpstep}[1]{\par\addvspace{0.3\baselineskip}\begin{list}{}{\leftmargin=1.6em\rightmargin=0pt\topsep=0pt}\item\relax
  {\small\color{bpmuted}#1}\par\nopagebreak}{\end{list}\par}
% Notes about where an object is written out.
\newenvironment{bpnote}{\par\small\color{bpmuted}\itshape}{\par}
\newcommand{\bpomitted}[1]{\par{\small\color{bpmuted}\itshape #1\par}}
\newcommand{\bpbodyheading}[1]{\par\addvspace{0.4\baselineskip}\noindent\textbf{#1}\par\nopagebreak}
\newcommand{\bprule}{\par\medskip\noindent\hfil\rule{0.3\linewidth}{0.4pt}\hfil\par\medskip}
\newcommand{\bpboundary}[1]{\par\nopagebreak{\small\color{bpmuted}#1\par}\nopagebreak}
% A reference to something this export leaves out, and one that resolves nowhere.
\newcommand{\bpout}[1]{#1\textsuperscript{\hyperlink{bp-outside}{\ensuremath{\ast}}}}
\newcommand{\bpbroken}[1]{\textcolor{red!70!black}{[#1?]}}
\newcommand{\bpslug}[1]{\texttt{#1}}
% A bibliography entry: label, anchor, entry; the label hangs in the margin.
\newcommand{\bpbibitem}[3]{\par\setbox0\hbox{[#1]\enspace}\ifdim\wd0<5em\setbox0\hbox to 5em{[#1]\hfil}\fi
  \hangindent=\wd0 \hangafter=1 \noindent\box0\phantomsection\label{#2}#3\par\smallskip}
% A character the export knows no way to set.
\newcommand{\bpglyph}[1]{\fbox{\scriptsize #1}}
% Images, at their own size unless they are wider than the line.
\newsavebox{\bpimagebox}
\newcommand{\bpimage}[1]{\sbox{\bpimagebox}{\includegraphics{#1}}%
  \ifdim\wd\bpimagebox>\linewidth\resizebox{\linewidth}{!}{\usebox{\bpimagebox}}\else\usebox{\bpimagebox}\fi}
\newcommand{\bpmissingimage}[1]{\fbox{\parbox{0.8\linewidth}{\centering\small\color{bpmuted}Image not included: #1}}}
\DefineVerbatimEnvironment{bpverbatim}{Verbatim}{fontsize=\small,xleftmargin=1em}
\renewcommand{\qedsymbol}{\ensuremath{\blacksquare}}
\hypersetup{pdftitle={Weil Conjectures — 2 The pro-étale site},pdfsubject={Chapter 2 The pro-étale site; all levels},pdfcreator={blueprint}}
% KaTeX commands the website knows
\providecommand{\R}{\mathbb{R}}\providecommand{\N}{\mathbb{N}}\providecommand{\Z}{\mathbb{Z}}
\providecommand{\Q}{\mathbb{Q}}\providecommand{\C}{\mathbb{C}}
\providecommand{\Reals}{\mathbb{R}}\providecommand{\reals}{\mathbb{R}}\providecommand{\Complex}{\mathbb{C}}
\providecommand{\cnums}{\mathbb{C}}\providecommand{\natnums}{\mathbb{N}}\providecommand{\Bbb}[1]{\mathbb{#1}}
\providecommand{\bold}[1]{\mathbf{#1}}\providecommand{\frak}[1]{\mathfrak{#1}}\providecommand{\mathbbm}[1]{\mathbb{#1}}
\providecommand{\lang}{\langle}\providecommand{\rang}{\rangle}\providecommand{\bra}[1]{\langle #1\rvert}
\providecommand{\ket}[1]{\lvert #1\rangle}\providecommand{\braket}[1]{\langle #1\rangle}
\providecommand{\set}[1]{\{#1\}}\providecommand{\Set}[1]{\left\{#1\right\}}
\providecommand{\infin}{\infty}\providecommand{\isin}{\in}\providecommand{\plusmn}{\pm}\providecommand{\sub}{\subset}
\providecommand{\sube}{\subseteq}\providecommand{\supe}{\supseteq}\providecommand{\exist}{\exists}
\providecommand{\bull}{\bullet}\providecommand{\larr}{\leftarrow}\providecommand{\rarr}{\rightarrow}
\providecommand{\lrarr}{\leftrightarrow}\providecommand{\Larr}{\Leftarrow}\providecommand{\Rarr}{\Rightarrow}
\providecommand{\Lrarr}{\Leftrightarrow}\providecommand{\uarr}{\uparrow}\providecommand{\darr}{\downarrow}
\providecommand{\harr}{\leftrightarrow}\providecommand{\lArr}{\Leftarrow}\providecommand{\rArr}{\Rightarrow}
\providecommand{\hArr}{\Leftrightarrow}\providecommand{\alef}{\aleph}\providecommand{\alefsym}{\aleph}
\providecommand{\real}{\Re}\providecommand{\image}{\Im}\providecommand{\weierp}{\wp}
\providecommand{\Alpha}{\mathrm{A}}\providecommand{\Beta}{\mathrm{B}}\providecommand{\Epsilon}{\mathrm{E}}
\providecommand{\Zeta}{\mathrm{Z}}\providecommand{\Eta}{\mathrm{H}}\providecommand{\Iota}{\mathrm{I}}
\providecommand{\Kappa}{\mathrm{K}}\providecommand{\Mu}{\mathrm{M}}\providecommand{\Nu}{\mathrm{N}}
\providecommand{\Omicron}{\mathrm{O}}\providecommand{\Rho}{\mathrm{P}}\providecommand{\Tau}{\mathrm{T}}
\providecommand{\Chi}{\mathrm{X}}\providecommand{\omicron}{o}\providecommand{\gt}{>}\providecommand{\lt}{<}
\providecommand{\And}{\&}\providecommand{\hdashline}{\hline}\providecommand{\operatornamewithlimits}[1]{\operatorname*{#1}}
\providecommand{\llbracket}{[\![}\providecommand{\rrbracket}{]\!]}\providecommand{\cancel}[1]{#1}
\providecommand{\mathscr}{\mathcal}
% the characters beyond ASCII this document uses
\DeclareUnicodeCharacter{2160}{\textup{I}}
\DeclareUnicodeCharacter{2161}{\textup{II}}
% the project's macros ([katex.macros] in blueprint.toml)
\providecommand{\A}{}\renewcommand{\A}{\mathbf A}
\providecommand{\AEn}{}\renewcommand{\AEn}{A_{E,n}}
\providecommand{\Aaf}{}\renewcommand{\Aaf}{\mathbf A}
\providecommand{\Ab}{}\renewcommand{\Ab}{\mathsf{Ab}}
\providecommand{\AffSch}{}\renewcommand{\AffSch}{\operatorname{AffSch}}
\providecommand{\Aone}{}\renewcommand{\Aone}{\mathbf A^1}
\providecommand{\Bl}{}\renewcommand{\Bl}{\operatorname{Bl}}
\providecommand{\Br}{}\renewcommand{\Br}{\mathrm{Br}}
\providecommand{\C}{}\renewcommand{\C}{\mathbb C}
\providecommand{\Cech}{}\renewcommand{\Cech}{\check{H}}
\providecommand{\Coext}{}\renewcommand{\Coext}{\operatorname{Coext}}
\providecommand{\Cone}{}\renewcommand{\Cone}{\operatorname{Cone}}
\providecommand{\Cons}{}\renewcommand{\Cons}{\operatorname{Cons}}
\providecommand{\Cref}{}\renewcommand{\Cref}[1]{\text{#1}}
\providecommand{\D}{}\renewcommand{\D}{D}
\providecommand{\DLAFT}{}\renewcommand{\DLAFT}{\operatorname{FT}}
\providecommand{\DLAFr}{}\renewcommand{\DLAFr}{\operatorname{Fr}}
\providecommand{\DLAQ}{}\renewcommand{\DLAQ}{{\overline{\mathbf Q}_\ell}}
\providecommand{\DLRQ}{}\renewcommand{\DLRQ}{{\overline{\mathbf Q}_\ell}}
\providecommand{\Dbc}{}\renewcommand{\Dbc}{D^b_c}
\providecommand{\Dc}{}\renewcommand{\Dc}{D_{\mathrm{cons}}}
\providecommand{\Dctf}{}\renewcommand{\Dctf}{D^b_{ctf}}
\providecommand{\Dic}{}\renewcommand{\Dic}{D_{\IndCons}}
\providecommand{\Dplus}{}\renewcommand{\Dplus}{D^{+}}
\providecommand{\End}{}\renewcommand{\End}{\operatorname{End}}
\providecommand{\Et}{}\renewcommand{\Et}{\mathrm{\acute{e}t}}
\providecommand{\Ex}{}\renewcommand{\Ex}{\mathrm{Ex}}
\providecommand{\Ext}{}\renewcommand{\Ext}{\operatorname{Ext}}
\providecommand{\F}{}\renewcommand{\F}{\mathbb F}
\providecommand{\FT}{}\renewcommand{\FT}{\operatorname{FT}}
\providecommand{\Fbar}{}\renewcommand{\Fbar}{{\overline{\mathbf F}_q}}
\providecommand{\Fell}{}\renewcommand{\Fell}{\mathbf F_\ell}
\providecommand{\Fib}{}\renewcommand{\Fib}{\operatorname{Fib}}
\providecommand{\Fix}{}\renewcommand{\Fix}{\operatorname{Fix}}
\providecommand{\Fp}{}\renewcommand{\Fp}{{\mathbf F_{p}}}
\providecommand{\Fq}{}\renewcommand{\Fq}{{\mathbf F_q}}
\providecommand{\Fqn}{}\renewcommand{\Fqn}{{\mathbf F_{q^{n}}}}
\providecommand{\Frob}{}\renewcommand{\Frob}{\operatorname{Frob}}
\providecommand{\GL}{}\renewcommand{\GL}{\operatorname{GL}}
\providecommand{\Gal}{}\renewcommand{\Gal}{\operatorname{Gal}}
\providecommand{\Hens}{}\renewcommand{\Hens}{\operatorname{Hens}}
\providecommand{\Hilb}{}\renewcommand{\Hilb}{\operatorname{Hilb}}
\providecommand{\Hom}{}\renewcommand{\Hom}{\operatorname{Hom}}
\providecommand{\IndCons}{}\renewcommand{\IndCons}{\operatorname{IndCons}}
\providecommand{\Kar}{}\renewcommand{\Kar}{\operatorname{Kar}}
\providecommand{\LT}{}\renewcommand{\LT}{\operatorname{LT}}
\providecommand{\Lie}{}\renewcommand{\Lie}{\operatorname{Lie}}
\providecommand{\Loc}{}\renewcommand{\Loc}{\operatorname{Loc}}
\providecommand{\Mod}{}\renewcommand{\Mod}{\mathsf{Mod}}
\providecommand{\N}{}\renewcommand{\N}{\mathbb{N}}
\providecommand{\Needspace}{}\renewcommand{\Needspace}{\par}
\providecommand{\O}{}\renewcommand{\O}{{\mathcal{O}}}
\providecommand{\Ob}{}\renewcommand{\Ob}{\mathrm{Ob}}
\providecommand{\Oe}{}\renewcommand{\Oe}{\mathcal O_E}
\providecommand{\PP}{}\renewcommand{\PP}{\mathbf{P}}
\providecommand{\Perf}{}\renewcommand{\Perf}{\operatorname{Perf}}
\providecommand{\Pic}{}\renewcommand{\Pic}{\mathrm{Pic}}
\providecommand{\Pm}{}\renewcommand{\Pm}{\mathbf P^{m}}
\providecommand{\Pone}{}\renewcommand{\Pone}{\mathbf P^{1}}
\providecommand{\Pqn}{}\renewcommand{\Pqn}{\mathbf P^{n}}
\providecommand{\ProFin}{}\renewcommand{\ProFin}{\operatorname{ProFin}}
\providecommand{\Proet}{}\renewcommand{\Proet}{\mathrm{pro\acute{e}t}}
\providecommand{\Q}{}\renewcommand{\Q}{\mathbb Q}
\providecommand{\Qb}{}\renewcommand{\Qb}{\overline{\mathbf Q}_{\ell}}
\providecommand{\Qbar}{}\renewcommand{\Qbar}{\overline{\mathbf Q}}
\providecommand{\QbarX}{}\renewcommand{\QbarX}{\overline{\mathbf Q}_{\ell,X}}
\providecommand{\Qbarl}{}\renewcommand{\Qbarl}{{\overline{\mathbf Q}_{\ell}}}
\providecommand{\Qell}{}\renewcommand{\Qell}{{\mathbf{Q}_{\ell}}}
\providecommand{\Ql}{}\renewcommand{\Ql}{\mathbf Q_{\ell}}
\providecommand{\Qlb}{}\renewcommand{\Qlb}{{\overline{\mathbf Q}_\ell}}
\providecommand{\Qlbar}{}\renewcommand{\Qlbar}{\overline{\mathbf Q}_\ell}
\providecommand{\R}{}\renewcommand{\R}{\mathbb{R}}
\providecommand{\RG}{}\renewcommand{\RG}{\operatorname{\mathrm{R}\Gamma}}
\providecommand{\RGamma}{}\renewcommand{\RGamma}{R\Gamma}
\providecommand{\RGammac}{}\renewcommand{\RGammac}{R\Gamma_c}
\providecommand{\RHom}{}\renewcommand{\RHom}{R\mathcal{H}om}
\providecommand{\RHomX}{}\renewcommand{\RHomX}{R\mathcal H om_X}
\providecommand{\Rat}{}\renewcommand{\Rat}{\operatorname{IsRat}}
\providecommand{\Res}{}\renewcommand{\Res}{\operatorname{Res}}
\providecommand{\Rlim}{}\renewcommand{\Rlim}{R\!\varprojlim}
\providecommand{\SIm}{}\renewcommand{\SIm}{\operatorname{SIm}}
\providecommand{\Sch}{}\renewcommand{\Sch}{\mathsf{Sch}}
\providecommand{\Sh}{}\renewcommand{\Sh}{\operatorname{Sh}}
\providecommand{\Span}{}\renewcommand{\Span}{\operatorname{span}}
\providecommand{\Spec}{}\renewcommand{\Spec}{\operatorname{Spec}}
\providecommand{\Tr}{}\renewcommand{\Tr}{\operatorname{Tr}}
\providecommand{\Trsm}{}\renewcommand{\Trsm}{\operatorname{Tr}^{\mathrm{sm}}}
\providecommand{\Wt}{}\renewcommand{\Wt}{\operatorname{Wt}}
\providecommand{\Z}{}\renewcommand{\Z}{\mathbb{Z}}
\providecommand{\Zar}{}\renewcommand{\Zar}{\mathrm{zar}}
\providecommand{\Zell}{}\renewcommand{\Zell}{{\mathbf{Z}_{\ell}}}
\providecommand{\Zl}{}\renewcommand{\Zl}{{\mathbf Z_\ell}}
\providecommand{\affet}{}\renewcommand{\affet}[1]{{#1}^{\mathrm{aff}}_{\mathrm{ét}}}
\providecommand{\affproet}{}\renewcommand{\affproet}[1]{{#1}^{\mathrm{aff}}_{\mathrm{proét}}}
\providecommand{\boxt}{}\renewcommand{\boxt}{\boxtimes_A^{\mathrm L}}
\providecommand{\cF}{}\renewcommand{\cF}{\mathcal{F}}
\providecommand{\cHom}{}\renewcommand{\cHom}{R\mathcal{H}\!om}
\providecommand{\cI}{}\renewcommand{\cI}{\mathcal{I}}
\providecommand{\cL}{}\renewcommand{\cL}{\mathcal L}
\providecommand{\cM}{}\renewcommand{\cM}{\mathcal M}
\providecommand{\cN}{}\renewcommand{\cN}{\mathcal N}
\providecommand{\cO}{}\renewcommand{\cO}{\mathcal O}
\providecommand{\cT}{}\renewcommand{\cT}{\mathcal T}
\providecommand{\checked}{}\renewcommand{\checked}{\color{red}\aftergroup\restorecolor}
\providecommand{\coeff}{}\renewcommand{\coeff}{\operatorname{coeff}}
\providecommand{\coker}{}\renewcommand{\coker}{\operatorname{coker}}
\providecommand{\colim}{}\renewcommand{\colim}{\operatorname{colim}}
\providecommand{\constructedby}{}\renewcommand{\constructedby}[1]{\par\smallskip\noindent\textcolor{purple!60!black}{\textsf{Constructed by:} }
  \footnotesize\sloppy\nolinkurl{#1}\par}
\providecommand{\deepresult}{}\renewcommand{\deepresult}[1]{\par\smallskip\noindent\textcolor{orange!45!black}{
  \textsf{Blueprint role: named deep theorem node; see the cluster proof architecture.}}
  \par\noindent\textcolor{orange!45!black}{\textsf{Mathematical source:} }#1\par}
\providecommand{\dt}{}\renewcommand{\dt}{\mathbin{\otimes^{\mathbf L}}}
\providecommand{\dw}{}\renewcommand{\dw}{\operatorname{dw}}
\providecommand{\et}{}\renewcommand{\et}[1]{{#1}_{\mathrm{ét}}}
\providecommand{\ev}{}\renewcommand{\ev}{\operatorname{ev}}
\providecommand{\fib}{}\renewcommand{\fib}{\operatorname{fib}}
\providecommand{\formalizationnote}{}\renewcommand{\formalizationnote}[1]{\begin{quote}\small\raggedright\textbf{Formalization note.} #1\end{quote}}
\providecommand{\fpqc}{}\renewcommand{\fpqc}{X_{\mathrm{fpqc}}}
\providecommand{\gammacont}{}\renewcommand{\gammacont}{\Gamma_{\mathrm{cont}}}
\providecommand{\gm}{}\renewcommand{\gm}{\mathbf G_m}
\providecommand{\hocolim}{}\renewcommand{\hocolim}{\operatorname{hocolim}}
\providecommand{\holim}{}\renewcommand{\holim}{\operatorname{holim}}
\providecommand{\hrB}{}\renewcommand{\hrB}{\@ifnextchar]{\hrEnd}{\hrStep}}
\providecommand{\hrEnd}{}\renewcommand{\hrEnd}{\@firstoftwo{}}
\providecommand{\hrStep}{}\renewcommand{\hrStep}{\@firstoftwo{\hrB}}
\providecommand{\id}{}\renewcommand{\id}{\operatorname{id}}
\providecommand{\kunFl}{}\renewcommand{\kunFl}{\mathbf F_\ell}
\providecommand{\kunZn}{}\renewcommand{\kunZn}{\mathbf Z/n\mathbf Z}
\providecommand{\kuncd}{}\renewcommand{\kuncd}{\operatorname{cd}}
\providecommand{\kundim}{}\renewcommand{\kundim}{\dim}
\providecommand{\kunotimes}{}\renewcommand{\kunotimes}[1]{\otimes^{\mathbf L}_{#1}}
\providecommand{\kunsupp}{}\renewcommand{\kunsupp}{\operatorname{Supp}}
\providecommand{\lotimes}{}\renewcommand{\lotimes}{\otimes^{\mathrm L}}
\providecommand{\m}{}\renewcommand{\m}{{\mathfrak{m}}}
\providecommand{\n}{}\renewcommand{\n}{{\mathfrak{n}}}
\providecommand{\nagataGap}{}\renewcommand{\nagataGap}[1]{\par\noindent\textbf{Formalization gap.} #1\par}
\providecommand{\openobligation}{}\renewcommand{\openobligation}[1]{\par\smallskip\noindent\textcolor{red!65!black}{\textsf{Open obligation:} }
  \footnotesize\sloppy\nolinkurl{#1}\par}
\providecommand{\p}{}\renewcommand{\p}{{\mathfrak{p}}}
\providecommand{\pr}{}\renewcommand{\pr}{\operatorname{pr}}
\providecommand{\preshv}{}\renewcommand{\preshv}[1]{\mathrm{PreShv}(#1)}
\providecommand{\pro}{}\renewcommand{\pro}[1]{\mathrm{Pro}(#1)}
\providecommand{\proet}{}\renewcommand{\proet}{X_{\mathrm{proét}}}
\providecommand{\q}{}\renewcommand{\q}{{\mathfrak{q}}}
\providecommand{\rad}{}\renewcommand{\rad}{\operatorname{rad}}
\providecommand{\rank}{}\renewcommand{\rank}{\operatorname{rank}}
\providecommand{\red}{}\renewcommand{\red}{\color{red}}
\providecommand{\restorecolor}{}\renewcommand{\restorecolor}{\color{black}}
\providecommand{\retiredobligation}{}\renewcommand{\retiredobligation}[1]{\par\smallskip\noindent\textcolor{red!75!black}{\textsf{Retired invalid route:} }
  \footnotesize\sloppy\nolinkurl{#1}\par}
\providecommand{\rk}{}\renewcommand{\rk}{\operatorname{rk}}
\providecommand{\sep}{}\renewcommand{\sep}{\textnormal{sep}}
\providecommand{\shv}{}\renewcommand{\shv}[1]{\mathrm{Shv}(#1)}
\providecommand{\spec}{}\renewcommand{\spec}[1]{\mathrm{Spec}(#1)}
\providecommand{\stacks}{}\renewcommand{\stacks}[1]{\href{https://stacks.math.columbia.edu/tag/#1}{\textnormal{Stacks, Tag~#1}}}
\providecommand{\stacksproject}{}\renewcommand{\stacksproject}[1]{{\cite[\href{https://stacks.math.columbia.edu/tag/#1}{Tag #1}]{stacks-project}}}
\providecommand{\statusnote}{}\renewcommand{\statusnote}[1]{\begin{quote}\small\raggedright\textbf{Status.} #1\end{quote}}
\providecommand{\tildeideal}{}\renewcommand{\tildeideal}[2]{\widetilde{(#1, #2)}}
\providecommand{\tildering}{}\renewcommand{\tildering}[3]{{#1}_{\widetilde{(#2, #3)}}}
\providecommand{\tr}{}\renewcommand{\tr}{\operatorname{tr}}
\providecommand{\wt}{}\renewcommand{\wt}{\operatorname{wt}}

\begin{document}
\pagestyle{plain}
\begin{center}
{\LARGE\bfseries Weil Conjectures\par}
\vspace{0.7em}
{\large Chapter 2: The pro-étale site\par}
\vspace{0.5em}
{\small\color{bpmuted}Snapshot 2026-09-26, revision 8e4cb88e2a\par}
\end{center}
\vspace{0.5em}
{\small\tableofcontents}
\bigskip\hrule\bigskip

% object: chap-The-pro-etale-site (blueprint/chap-The-pro-etale-site/_section.md)
\section*{\texorpdfstring{2\quad The pro-étale site}{2 The pro-étale site}}
\phantomsection\label{bp1}%
\addcontentsline{toc}{section}{\texorpdfstring{2\quad The pro-étale site}{2 The pro-étale site}}
% text: prose-02-01
\phantomsection\label{bp2}%
On the étale site the constant sheaf \(\mathbf Z_\ell\) forgets that \(\mathbf Z_\ell\) is profinite. The pro-étale site \(X_{\mathrm{pro}\acute{e}t}\) of Bhatt and Scholze [\hyperref[bp191]{BS15}] repairs this. Its objects are the weakly étale \(X\)-schemes, and its coverings are fpqc coverings (Definition \hyperref[bp12]{2.4}).

We need three properties of this site.

\begin{itemize}
\item (a) \emph{Limits.} For a tower \(\cdots\to F_2\to F_1\) of surjections of sheaves, cohomology commutes with the limit: \(R\Gamma(X,\varprojlim_nF_n)\simeq R\varprojlim_nR\Gamma(X,F_n)\).
\item (b) \emph{Topology.} The sheaf \(\mathbf Z_\ell\) remembers the topology of \(\mathbf Z_\ell\).
\item (c) \emph{Comparison.} Sheaves coming from the étale site keep their étale cohomology, so that the results of SGA 4 still apply to them.
\end{itemize}

Property (b) can be seen on a point. Over a separably closed field the topos of \(X_{\mathrm{pro}\acute{e}t}\) is, up to a set-theoretic cutoff, that of condensed sets (Proposition \hyperref[bp151]{2.50}, proved in Subsection \hyperref[bp149]{2.12.1} below), and there the sections of \(\mathbf Z_\ell\) over a profinite set \(S\) are all continuous maps \(S\to\mathbf Z_\ell\), not only the locally constant ones.

\begin{center}\small
\begin{tabularx}{\linewidth}{>{\raggedright\arraybackslash}X>{\raggedright\arraybackslash}X>{\raggedright\arraybackslash}X}
\toprule
 & on \(X_{\acute{e}t}\) & on \(X_{\mathrm{pro}\acute{e}t}\) \\
\midrule
objects & étale \(X\)-schemes & weakly étale \(X\)-schemes \\
local objects & strictly henselian local rings & w-contractible rings \\
topos of a geometric point & sets & condensed sets \\
the sheaf \(\mathbf Z_\ell\) on a point & \(\mathbf Z_\ell\) with discrete topology & continuous maps to \(\mathbf Z_\ell\) \\
\bottomrule
\end{tabularx}
\end{center}

Property (a) fails on \(X_{\acute{e}t}\), for the following reason. A surjection of sheaves is surjective on sections only after refining the cover. A countable product of surjections needs one refinement for each factor, and these infinitely many refinements need not have a common refinement among étale covers. So a countable product of surjections need not be surjective, \(R^1\varprojlim_nF_n\) need not vanish, and \(\varprojlim_nF_n\) differs from \(R\varprojlim_nF_n\). For \(F_n=\mathbf Z/\ell^n\) on \(\operatorname{Spec}\mathbf F_q\) this is the lost \(H^1\). On \(X_{\mathrm{pro}\acute{e}t}\) there are enough objects on which no refinement is needed: the spectra of w-contractible rings, roughly families of strictly henselian local rings parametrized by an extremally disconnected profinite set. Every cover of such a spectrum splits, so every surjection of sheaves is surjective on its sections there, and limits of sheaves can be computed there as limits of sets and modules. However, building these objects takes two theorems from outside algebraic geometry: Gleason's theorem on extremally disconnected spaces and Olivier's theorem on weakly étale maps. Once we have them, property (a) follows in a few lines (Proposition \hyperref[bp170]{2.55}). In particular, over \(\mathbf F_q\) it gives back \(H^1=\mathbf Z_\ell\) (Example \hyperref[bp176]{2.57}).

For property (c) we compare the two sites through the pullback functor \(\nu_X^*\) from sheaves on \(X_{\acute{e}t}\) to sheaves on \(X_{\mathrm{pro}\acute{e}t}\). Its image consists of the sheaves that we call \emph{classical}: a sheaf on \(X_{\mathrm{pro}\acute{e}t}\) is classical if its sections over a limit \(U=\varprojlim_iU_i\) of affine étale \(X\)-schemes are the colimit of its sections over the \(U_i\) (Definition \hyperref[bp89]{2.30}). In other words, a classical sheaf is determined by its values on étale \(X\)-schemes, as an étale sheaf is. The comparison has four parts.

\begin{itemize}
\item \(\nu_X^*\) is an equivalence from étale sheaves onto classical sheaves (Theorem \hyperref[bp93]{2.32}).
\item It preserves the cohomology of bounded below complexes (Theorem \hyperref[bp101]{2.35}).
\item It commutes with direct images along quasi-compact quasi-separated morphisms (Theorem \hyperref[bp111]{2.38}).
\item For a finite coefficient ring, a constructible sheaf on \(X_{\mathrm{pro}\acute{e}t}\) is classical, and \(\nu_X^*\) is an equivalence from \(D^b_c(X_{\acute{e}t},\Lambda)\) onto \(\Dbc(X,\Lambda)\) (Theorem \hyperref[bp127]{2.43}). So with finite coefficients the two sites have the same constructible complexes.
\end{itemize}

Later, \(f^!\) is built by Brown representability on the unbounded derived category, so we also need the comparison for unbounded complexes. With finite coefficients it holds as well, because affine schemes étale over a scheme of finite type over \(k\) have uniformly bounded cohomological dimension. Each cohomology group of an unbounded complex therefore depends only on a bounded below truncation, where the comparison is already known (Proposition \hyperref[bp144]{2.49}). The sheaf \(\mathbf Z_\ell\), however, is not classical (Example \hyperref[bp95]{2.33}). This is what the computation of \(H^1(\operatorname{Spec}\mathbf F_q,\mathbf Z_\ell)\) in Section \bpout{1.2} needs: \(H^1\) comes out as \(\mathbf Z_\ell\) only when sections see the topology of \(\mathbf Z_\ell\). So the comparison applies to the \(\mathbf Z/\ell^n\) but not to \(\mathbf Z_\ell\) itself, and we reach \(\mathbf Z_\ell\) from them by derived completion (Section \hyperref[bp177]{2.14}).

Throughout this chapter, schemes are arbitrary unless a result states a finiteness hypothesis, because objects of \(X_{\mathrm{pro}\acute{e}t}\) are in general neither Noetherian nor of finite type, even when \(X\) is a variety. The standing assumption that schemes are of finite type over \(k\) resumes in Chapter \bpout{3}.

% object: sec-Weakly-etale-morphisms (blueprint/chap-The-pro-etale-site/sec-Weakly-etale-morphisms/_section.md), content hidden
\subsection*{\texorpdfstring{2.1\quad Weakly étale morphisms}{2.1 Weakly étale morphisms}}
\phantomsection\label{bp3}%
\addcontentsline{toc}{subsection}{\texorpdfstring{2.1\quad Weakly étale morphisms}{2.1 Weakly étale morphisms}}
% text: prose-02-02
\phantomsection\label{bp4}%
The objects we want are inverse limits of étale \(X\)-schemes. Being such a limit is a condition on a presentation, which is awkward to work with, so we take as objects the weakly étale \(X\)-schemes of the following definition. This class is defined by flatness conditions. For this reason it is stable under base change and composition, and membership in it is local for the Zariski topology (Lemma \hyperref[bp10]{2.3}). We lose little by the change: after a faithfully flat refinement every weakly étale scheme becomes a limit of étale schemes (Theorem \hyperref[bp40]{2.14}).

Typical examples are the schemes that the étale site can only reach through colimits. The strict henselization \(\operatorname{Spec}\mathcal O^{\mathrm{sh}}_{X,\bar x}\) is the limit of the affine étale neighbourhoods of \(\bar x\), and on \(X_{\acute{e}t}\) it appears only through the stalk at \(\bar x\). On \(X_{\mathrm{pro}\acute{e}t}\) it is an object. Similarly \(\operatorname{Spec}\overline{\mathbf F}_q=\varprojlim_m\operatorname{Spec}\mathbf F_{q^m}\) is an object over \(\operatorname{Spec}\mathbf F_q\), and for a profinite set \(S=\varprojlim S_i\) the scheme \(X\times S=\varprojlim_i\coprod_{S_i}X\) is an object over \(X\). Through the last family the site sees profinite sets, and hence the topology of \(\mathbf Z_\ell\).

% object: definition-ind-etale-pro-etale-weakly (blueprint/chap-The-pro-etale-site/sec-Weakly-etale-morphisms/definition-ind-etale-pro-etale-weakly/_section.md)
\phantomsection\label{bp5}%
\begin{bpstatement}{Definition~2.1}{ (Ind-étale, pro-étale and weakly étale morphisms)}
A ring map \(R\to B\) is \emph{ind-étale} if \(B=\varinjlim_\alpha B_\alpha\) for a filtered system of étale \(R\)-algebras \(B_\alpha\). A \emph{pro-étale affine over \(X\)} is a scheme \(U=\varprojlim_iU_i\) that is the limit of a cofiltered system of affine schemes \(U_i\) étale over \(X\). If \(X=\operatorname{Spec}(R)\) is affine, the pro-étale affines over \(X\) are the schemes

\[
\operatorname{Spec}(B)=\varprojlim_\alpha
\operatorname{Spec}(B_\alpha)\longrightarrow X
\]

with \(R\to B\) ind-étale.

A scheme morphism \(u:U\to X\) is \emph{weakly étale} if \(u\) and its diagonal \(\Delta_u:U\to U\times_XU\) are flat.
\end{bpstatement}

% text: prose-02-03
\phantomsection\label{bp6}%
The three classes of ring maps are related as follows.

% object: proposition-etale-implications (blueprint/chap-The-pro-etale-site/sec-Weakly-etale-morphisms/proposition-etale-implications/_section.md)
\phantomsection\label{bp7}%
\begin{bpstatementit}{Proposition~2.2}{ (Étale, ind-étale and weakly étale ring maps)}
Let \(R\to B\) be a ring map.

\begin{enumerate}
\item If \(R\to B\) is étale, it is ind-étale; if it is ind-étale, it is weakly étale.
\item If \(R\to B\) is weakly étale and of finite presentation, it is étale.
\end{enumerate}
\end{bpstatementit}

% proof: proof/proposition-etale-implications
\begin{proof}[{Proof}]\phantomsection\label{bp8}%
The first implication in (1) is immediate from the definition, since an étale algebra is the colimit of the constant system on itself. For the second, a filtered colimit of étale algebras is flat, and its multiplication map is the filtered colimit of the multiplication maps \(B_\alpha\otimes_RB_\alpha\to B_\alpha\), which are flat; see [Stacks, \href{https://stacks.math.columbia.edu/tag/092A}{092A}].

For (2) we show that \(R\to B\) is unramified, which together with flatness and finite presentation means that it is étale. The multiplication map \(B\otimes_RB\to B\) is surjective, and it is flat by assumption. It is of finite presentation because \(B\) is of finite presentation over \(R\). Hence the diagonal of \(\operatorname{Spec}(B)\to\operatorname{Spec}(R)\) is a closed immersion that is flat and of finite presentation. Its ideal \(I\) is then finitely generated with \(I=I^2\), hence generated by an idempotent, so the diagonal is also an open immersion. A morphism locally of finite type whose diagonal is an open immersion is unramified, and this proves (2).
\end{proof}

% text: prose-02-04
\phantomsection\label{bp9}%
For instance, \(\mathbf F_q\to\overline{\mathbf F}_q\) is ind-étale, as the colimit of the finite fields \(\mathbf F_{q^m}\), but it is not étale, because \(\overline{\mathbf F}_q\) is not of finite type over \(\mathbf F_q\).

In (2) the flatness of the diagonal takes the place of unramifiedness; flatness of \(R\to B\) alone would not suffice. The converse of the second implication in (1) fails in general: a weakly étale map need not be a filtered colimit of étale maps. It does become ind-étale after a faithfully flat refinement (Theorem \hyperref[bp40]{2.14} below, and [Stacks, \href{https://stacks.math.columbia.edu/tag/097Z}{097Z}]). Geometrically, a weakly étale scheme over \(X\) is covered, in the pro-étale topology, by a limit of étale \(X\)-schemes.

% object: lemma-stability (blueprint/chap-The-pro-etale-site/sec-Weakly-etale-morphisms/lemma-stability/_section.md)
\phantomsection\label{bp10}%
\begin{bpstatementit}{Lemma~2.3}{ (Stability)}
Weakly étale morphisms are stable under base change and composition, and the property is local on source and target for the Zariski topology. Every \(X\)-morphism between two weakly étale \(X\)-schemes is weakly étale. Ind-étale ring maps are stable under base change and composition. On affine charts, a morphism \(\operatorname{Spec}(B)\to\operatorname{Spec}(R)\) is weakly étale if and only if both

\[
R\longrightarrow B,
\qquad B\otimes_RB\longrightarrow B
\]

are flat.
\end{bpstatementit}

% proof: proof/lemma-stability
\begin{proof}[{Proof}]\phantomsection\label{bp11}%
Since weakly étale means that a morphism and its diagonal are flat, each assertion about weakly étale morphisms comes down to writing the relevant diagonal in terms of diagonals and base changes that we already know to be flat.

Base change preserves flatness of a morphism and its diagonal. For a composite \(V\to U\to X\), the morphism \(V\to X\) is flat as a composite of flat morphisms. Its diagonal \(\Delta_{V/X}\) is \(\Delta_{V/U}\) followed by \(V\times_UV\to V\times_XV\), and the second map is a base change of \(\Delta_{U/X}\). Both maps are flat, and hence so is \(\Delta_{V/X}\).

Let \(h:U\to V\) be an \(X\)-morphism between weakly étale \(X\)-schemes. We factor \(h\) as the graph \(U\to U\times_XV\) followed by the projection \(U\times_XV\to V\). The graph is a base change of the diagonal of \(V/X\), which is flat, and the projection is a base change of \(U\to X\), which is also flat. Their composite \(h\) is therefore flat. Similarly, the diagonal \(\Delta_h:U\to
U\times_VU\) is the base change of \(\Delta_{U/X}\) along the monomorphism \(U\times_VU\to U\times_XU\), so it is flat as well. Finally, Zariski locality follows from the fact that flatness is Zariski local.

On affine charts, the diagonal of \(\operatorname{Spec}(B)\to\operatorname{Spec}(R)\) is the closed immersion \(\operatorname{Spec}(B)\to\operatorname{Spec}(B\otimes_RB)\) given by the multiplication map, so the morphism is weakly étale exactly when \(R\to B\) and \(B\otimes_RB\to B\) are flat. By Zariski locality on source and target, this describes weakly étale morphisms on any affine charts.

For ind-étale ring maps, base change is clear because tensor products commute with filtered colimits. For composition, let \(R\to B\) and \(B\to C\) be ind-étale, with \(B=\varinjlim_\alpha B_\alpha\). Each étale \(B\)-algebra \(C'\) occurring in a presentation of \(C\) is of finite presentation over \(B\), hence is the base change of an étale algebra \(C'_\alpha\) over some \(B_\alpha\). Since \(B_\alpha\) is étale over \(R\), so is \(C'_\alpha\), and \(C'\) is the filtered colimit of the étale \(R\)-algebras \(C'_\alpha\otimes_{B_\alpha}B_\beta\) for \(\beta\geq\alpha\). Thus \(C\) is a filtered colimit of ind-étale \(R\)-algebras, and such a colimit is again ind-étale [Stacks, \href{https://stacks.math.columbia.edu/tag/097L}{097L}].
\end{proof}

% object: definition-small-pro-etale-site (blueprint/chap-The-pro-etale-site/sec-Weakly-etale-morphisms/definition-small-pro-etale-site/_section.md)
\phantomsection\label{bp12}%
\begin{bpstatement}{Definition~2.4}{ (The small pro-étale site)}
The objects of \(X_{\mathrm{pro}\acute{e}t}\) are weakly étale \(X\)-schemes, with all \(X\)-morphisms as arrows. A family \(\{U_i\to U\}_{i\in I}\) is a covering if it is an fpqc covering. Explicitly, for every affine open \(W\subset U\) there must be finitely many indices \(i_1,\ldots,i_m\) and affine opens \(V_j\subset U_{i_j}\times_UW\) with

\[
W=\bigcup_{j=1}^m\operatorname{image}(V_j\to W).
\]

The indices may depend on \(W\), and \(I\) may be infinite. Equivalently one may use finitely many quasi-compact opens above each \(W\). See [Stacks, \href{https://stacks.math.columbia.edu/tag/0989}{0989}, \href{https://stacks.math.columbia.edu/tag/098A}{098A}].
\end{bpstatement}

% text: prose-02-05
\phantomsection\label{bp13}%
The finite refinement condition is part of the topology, and a jointly surjective family that does not satisfy it is not a covering. We record two consequences of the definition.

% object: proposition-pro-etale-coverings (blueprint/chap-The-pro-etale-site/sec-Weakly-etale-morphisms/proposition-pro-etale-coverings/_section.md)
\phantomsection\label{bp14}%
\begin{bpstatementit}{Proposition~2.5}{ (Coverings in the pro-étale site)}
Let \(U\) be an object of \(X_{\mathrm{pro}\acute{e}t}\).

\begin{enumerate}
\item If \(U\) is affine, every covering of \(U\) can be refined by a finite affine covering, whose disjoint union is affine and faithfully flat over \(U\).
\item Every étale covering and every Zariski covering of \(U\) is a pro-étale covering.
\end{enumerate}
\end{bpstatementit}

% proof: proof/proposition-pro-etale-coverings
\begin{proof}[{Proof}]\phantomsection\label{bp15}%
Part (1) is the covering condition of Definition \hyperref[bp12]{2.4} with \(W=U\). For (2), recall that flat morphisms locally of finite presentation are open. By quasi-compactness we then find the required finite refinements; see [Stacks, \href{https://stacks.math.columbia.edu/tag/098B}{098B}].
\end{proof}

% text: prose-02-06
\phantomsection\label{bp16}%
By (2) the pro-étale topology is finer than the étale topology. By (1), over an affine object a covering may always be replaced by a single faithfully flat map. For example, the single map \(\operatorname{Spec}\overline{\mathbf F}_q\to\operatorname{Spec}\mathbf F_q\) is a pro-étale covering, although it is not an étale one. These coverings define a topos with the expected functoriality in \(X\).

% object: proposition-topos-functoriality (blueprint/chap-The-pro-etale-site/sec-Weakly-etale-morphisms/proposition-topos-functoriality/_section.md)
\phantomsection\label{bp17}%
\begin{bpstatementit}{Proposition~2.6}{ (The topos and its functoriality)}
The above coverings define a subcanonical site. Its category of sheaves \(\Sh(X_{\mathrm{pro}\acute{e}t})\) is a Grothendieck topos. Every scheme morphism \(f:X\to Y\) induces a geometric morphism

\[
f_{\mathrm{pro}\acute{e}t}:\Sh(X_{\mathrm{pro}\acute{e}t})\longrightarrow\Sh(Y_{\mathrm{pro}\acute{e}t}),
\qquad f^{-1}\dashv f_*.
\]

Here \(f^{-1}:\Sh(Y_{\mathrm{pro}\acute{e}t})\to\Sh(X_{\mathrm{pro}\acute{e}t})\) preserves finite limits and all colimits, and is exact on abelian sheaves. For a second morphism \(g:Y\to Z\) there is a canonical isomorphism \((g\circ f)^{-1}\simeq f^{-1}\circ g^{-1}\).
\end{bpstatementit}

% proof: proof/proposition-topos-functoriality
\begin{proof}[{Proof}]\phantomsection\label{bp18}%
There are two things to check: that the coverings form a subcanonical topology, and that base change of objects is a morphism of sites.

Identity covers, composites of covers, and base changes of covers are again covers. For the underlying morphisms this is Lemma \hyperref[bp10]{2.3}, and for the finite refinement condition it is part of the fpqc covering axioms. Representable presheaves satisfy descent for these covers by fpqc descent for morphisms of schemes, so they are sheaves, and the topology is subcanonical.

For \(f:X\to Y\), the functor \(V\mapsto V\times_YX\) sends weakly étale \(Y\)-schemes to weakly étale \(X\)-schemes by Lemma \hyperref[bp10]{2.3}. It preserves fibre products and the final object, hence finite limits, and it carries covering families to covering families because fpqc coverings are stable under base change. Such a functor defines a geometric morphism, whose inverse image \(f^{-1}\) is left exact and, as a left adjoint, preserves all colimits. Exactness on abelian sheaves follows, and compatibility with composition holds because \((V\times_ZY)\times_YX\simeq V\times_ZX\) canonically. This is the small-site version of the constructions in [Stacks, \href{https://stacks.math.columbia.edu/tag/0988}{0988}].
\end{proof}


% object: sec-w-local-and-w-contractible-rings (blueprint/chap-The-pro-etale-site/sec-w-local-and-w-contractible-rings/_section.md), content hidden
\subsection*{\texorpdfstring{2.2\quad w-local and w-contractible rings}{2.2 w-local and w-contractible rings}}
\phantomsection\label{bp19}%
\addcontentsline{toc}{subsection}{\texorpdfstring{2.2\quad w-local and w-contractible rings}{2.2 w-local and w-contractible rings}}
% text: prose-02-07
\phantomsection\label{bp20}%
To compute limits of sheaves by taking sections, we want objects over which every cover splits, so that sheafification changes nothing there. On the étale site the spectrum of a strictly henselian local ring, such as the étale local ring of a variety at a geometric point, has this property, but such spectra do not cover \(X\). Instead, the rings we need are families of strictly henselian local rings over a profinite set of components. Moreover, we require this set to be extremally disconnected, which means that the closure of every open subset is open, because under this condition a cover splits. Namely, over each component of \(\operatorname{Spec}(R)\) we choose a component of the cover lying above it, and extremal disconnectedness lets us make this choice continuously. The model case is the ring of locally constant functions from a profinite set to a separably closed field (Example \hyperref[bp26]{2.9}), which is a family of copies of one field indexed by the profinite set. Below we define these rings, the w-contractible rings, by a hierarchy of local conditions, and we show that every ring has a faithfully flat ind-étale cover by one of them (Theorem \hyperref[bp37]{2.13}).

Recall that spectral spaces are quasi-compact and sober, with a basis of quasi-compact opens closed under finite intersections, and that the space of connected components of a spectral space is profinite.

% object: definition-w-local-w-strictly-local (blueprint/chap-The-pro-etale-site/sec-w-local-and-w-contractible-rings/definition-w-local-w-strictly-local/_section.md)
\phantomsection\label{bp21}%
\begin{bpstatement}{Definition~2.7}{ (w-local, w-strictly local and w-contractible rings)}
A topological space is \emph{extremally disconnected} if the closure of every open subset is open. A spectral space \(T\) is \emph{w-local} if its set \(T^c\) of closed points is closed and every point specializes to exactly one closed point. A ring \(R\) is w-local if \(\operatorname{Spec}(R)\) is w-local. It is \emph{w-strictly local} if it is w-local and \(R_{\mathfrak m}\) is strictly henselian for every maximal ideal \(\mathfrak m\). It is \emph{w-contractible} if it is w-strictly local and the profinite space \(\pi_0(\operatorname{Spec}(R))\) is extremally disconnected.
\end{bpstatement}

% text: prose-02-08
\phantomsection\label{bp22}%
Closedness of \(T^c\) does not follow from uniqueness of closed specializations, which is why the definition requires it separately. The basic properties of these notions are as follows.

% object: proposition-w-local-components (blueprint/chap-The-pro-etale-site/sec-w-local-and-w-contractible-rings/proposition-w-local-components/_section.md)
\phantomsection\label{bp23}%
\begin{bpstatementit}{Proposition~2.8}{ (Connected components of w-local spaces and rings)}
\begin{enumerate}
\item If \(T\) is a w-local spectral space, the natural map \(T^c\to\pi_0(T)\) is a homeomorphism.
\item If \(R\) is w-local, the component of \(\operatorname{Spec}(R)\) containing a closed point \(\mathfrak m\) is the image of \(\operatorname{Spec}(R_{\mathfrak m})\). In particular, if \(R\) is w-strictly local, every component is the spectrum of a strictly henselian local ring.
\item A strictly henselian local ring is w-contractible.
\end{enumerate}
\end{bpstatementit}

% proof: proof/proposition-w-local-components
\begin{proof}[{Proof}]\phantomsection\label{bp24}%
For (1), each component consists of the generalizations of its unique closed point, and \(T^c\) is compact. Part (2) follows, and (3) holds because the spectrum of a strictly henselian local ring has one closed point and one component. See [Stacks, \href{https://stacks.math.columbia.edu/tag/096A}{096A}, \href{https://stacks.math.columbia.edu/tag/0968}{0968}, \href{https://stacks.math.columbia.edu/tag/0969}{0969}].
\end{proof}

% text: prose-02-09
\phantomsection\label{bp25}%
So a w-strictly local ring is a family of strictly henselian local rings indexed by the profinite set of its components, and a w-contractible ring is such a family indexed by an extremally disconnected profinite set.

% object: example-locally-constant-functions-w-local (blueprint/chap-The-pro-etale-site/sec-w-local-and-w-contractible-rings/example-locally-constant-functions-w-local/_section.md)
\phantomsection\label{bp26}%
\begin{bpstatement}{Example~2.9}{ (Locally constant functions)}
Let \(k'\) be a separably closed field, let \(S=\varprojlim_iS_i\) be a profinite set written as a limit of finite sets, and let \(R=LC(S,k')\) be the ring of locally constant functions \(S\to k'\). Then \(R=\varinjlim_ik'^{\,S_i}\), and \(\operatorname{Spec}(k'^{\,S_i})\) is the finite discrete set \(S_i\), so \(\operatorname{Spec}(R)=\varprojlim_iS_i=S\): the point \(s\in S\) is the maximal ideal of functions vanishing at \(s\). Every point is closed and is its own component, so \(R\) is w-local with \(\pi_0(\operatorname{Spec}(R))=S\). The local ring at \(s\) is the colimit of the local rings \(k'\) of the \(k'^{\,S_i}\) at the image of \(s\), that is, \(k'\), which is strictly henselian. Thus \(R\) is w-strictly local: it is the family of copies of \(\operatorname{Spec}k'\) indexed by \(S\). It is w-contractible exactly when \(S\) is extremally disconnected. In this example the definitions are statements about \(S\) alone.
\end{bpstatement}

% text: prose-02-10
\phantomsection\label{bp27}%
The constructions below change the set of components while keeping the local rings, and for this we use the following class of maps. A ring map \(R\to A\) is a \emph{local isomorphism} if \(\operatorname{Spec}(A)\) is covered by principal opens \(D(g)\), each mapping isomorphically onto an open subscheme of \(\operatorname{Spec}(R)\) [Stacks, \href{https://stacks.math.columbia.edu/tag/096E}{096E}], and it is \emph{ind-Zariski} if it is a filtered colimit of local isomorphisms [Stacks, \href{https://stacks.math.columbia.edu/tag/096M}{096M}]. For instance, \(R\to R\times R\) is a local isomorphism. Taking the colimit over the finite quotients \(S_i\) of a profinite set \(S\), we see that \(R\to\varinjlim_iR^{S_i}=LC(S,R)\) is ind-Zariski. Its spectrum is \(\operatorname{Spec}(R)\times S\). A localization \(R\to R_{\mathfrak p}\) at a prime is ind-Zariski as well, as the colimit of the \(R_f\) with \(f\notin\mathfrak p\). The property we use is that an ind-Zariski map does not change local rings: each local ring of \(A\) is a local ring of \(R\).

We use extremally disconnected spaces through their projectivity. In Proposition \hyperref[bp34]{2.12} we use it to split a surjection of component spaces, and in Theorem \hyperref[bp37]{2.13} we cover a profinite component space by an extremally disconnected one (Lemma \hyperref[bp31]{2.11}).

% object: proposition-gleason-s-theorem (blueprint/chap-The-pro-etale-site/sec-w-local-and-w-contractible-rings/proposition-gleason-s-theorem/_section.md)
\phantomsection\label{bp28}%
\begin{bpstatementit}{Proposition~2.10}{ (Gleason's theorem)}
A compact Hausdorff space \(S\) is extremally disconnected if and only if it is projective with respect to continuous surjections of compact Hausdorff spaces: for every continuous surjection \(T_1\to T_2\) of compact Hausdorff spaces, every continuous map \(S\to T_2\) lifts to a continuous map \(S\to T_1\).
\end{bpstatementit}

% proof: proof/proposition-gleason-s-theorem
\begin{proof}[{Proof}]\phantomsection\label{bp29}%
This is the projectivity theorem for compact Hausdorff spaces; see [Stacks, \href{https://stacks.math.columbia.edu/tag/08YN}{08YN}]. Here we sketch the implication from extremal disconnectedness to projectivity, which is the direction used later. To lift a map \(S\to T_2\) along a surjection \(T_1\to T_2\), one pulls the surjection back to \(S\) and chooses a minimal closed subspace of the fibre product that still surjects onto \(S\). Then extremal disconnectedness forces this irreducible surjection to be a homeomorphism, and its inverse gives the lift.
\end{proof}

% text: prose-02-11
\phantomsection\label{bp30}%
To use the proposition we need enough extremally disconnected spaces. They come from Stone–Čech compactifications of discrete sets.

% object: lemma-extremally-disconnected-covers (blueprint/chap-The-pro-etale-site/sec-w-local-and-w-contractible-rings/lemma-extremally-disconnected-covers/_section.md)
\phantomsection\label{bp31}%
\begin{bpstatementit}{Lemma~2.11}{ (Extremally disconnected covers)}
Every compact Hausdorff space \(T\), in particular every profinite space, admits a continuous surjection \(S\to T\) with \(S\) compact Hausdorff and extremally disconnected. Every extremally disconnected compact Hausdorff space is profinite.
\end{bpstatementit}

% proof: proof/lemma-extremally-disconnected-covers
\begin{proof}[{Proof}]\phantomsection\label{bp32}%
Write \(T^{\delta}\) for the underlying set of \(T\) with the discrete topology and take \(S=\beta(T^{\delta})\), its Stone–Čech compactification. The Stone–Čech compactification of a discrete space is extremally disconnected, and the identity \(T^{\delta}\to T\) extends uniquely to a continuous surjection \(\beta(T^{\delta})\to T\); see [Stacks, \href{https://stacks.math.columbia.edu/tag/090D}{090D}]. An extremally disconnected compact Hausdorff space is totally disconnected, since the closure of an open set is clopen, and hence it is profinite.
\end{proof}

% text: prose-02-12
\phantomsection\label{bp33}%
A finite discrete space is extremally disconnected, while the Cantor set \(\{0,1\}^{\mathbf N}\) is not. To see this, let \(V_n\) be the clopen set of sequences beginning with \(0^n1\), and let \(U\) be the union of the \(V_n\) for odd \(n\). The closure of \(U\) is \(U\cup\{000\cdots\}\), and it is not open, because every neighbourhood of \(000\cdots\) contains some \(V_n\) with \(n\) even, which is disjoint from the closure. More generally, infinite extremally disconnected compact spaces have no nontrivial convergent sequences, and therefore they are never metrizable. In practice they arise as Stone–Čech compactifications.

Over a w-contractible ring an ind-étale cover splits at each closed point by strict henselianity, and Gleason's theorem lets us choose these local splittings continuously in the point.

% object: proposition-splitting-ind-etale-covers (blueprint/chap-The-pro-etale-site/sec-w-local-and-w-contractible-rings/proposition-splitting-ind-etale-covers/_section.md)
\phantomsection\label{bp34}%
\begin{bpstatementit}{Proposition~2.12}{ (Splitting ind-étale covers)}
If \(R\) is w-contractible, every faithfully flat ind-étale ring map \(R\to B\) has an \(R\)-algebra retraction \(B\to R\).
\end{bpstatementit}

% proof: proof/proposition-splitting-ind-etale-covers
\begin{proof}[{Proof (Proof outline)}]\phantomsection\label{bp35}%
We first replace \(B\) by an ind-Zariski \(B\)-algebra, so that its spectrum is w-local and its closed points lie exactly over the closed points of \(\operatorname{Spec}(R)\); such a refinement does not change the local rings. Let \(\mathfrak q\subset B\) lie over a closed point \(\mathfrak m\) of \(\operatorname{Spec}(R)\). Since \(R_{\mathfrak m}\) is strictly henselian, the ind-étale local homomorphism \(R_{\mathfrak m}\to B_{\mathfrak q}\) is an isomorphism. Hence every component of \(\operatorname{Spec}(B)\) maps isomorphically onto the component of \(\operatorname{Spec}(R)\) below it, and a retraction amounts to a continuous choice of one component above each component of \(\operatorname{Spec}(R)\). Such a choice exists because, by Proposition \hyperref[bp28]{2.10}, the surjection \(\pi_0(\operatorname{Spec}(B))\to\pi_0(\operatorname{Spec}(R))\) of profinite component spaces has a section. Given such a section, we restrict \(B\) to the corresponding closed set of components by a filtered colimit of localizations at idempotents. The result is an \(R\)-algebra isomorphic to \(R\), and this gives the retraction. For details see [Stacks, \href{https://stacks.math.columbia.edu/tag/097D}{097D}, \href{https://stacks.math.columbia.edu/tag/09AZ}{09AZ}, \href{https://stacks.math.columbia.edu/tag/0982}{0982}].
\end{proof}

% text: prose-02-13
\phantomsection\label{bp36}%
The proposition would be of little use if w-contractible rings were rare. The next theorem shows that every ring has a faithfully flat ind-étale cover by one, so that, after a single refinement, every affine object of \(X_{\mathrm{pro}\acute{e}t}\) is replaced by one over which covers split.

% object: theorem-existence-contractible-covers (blueprint/chap-The-pro-etale-site/sec-w-local-and-w-contractible-rings/theorem-existence-contractible-covers/_section.md)
\phantomsection\label{bp37}%
\begin{bpstatementit}{Theorem~2.13}{ (Existence of contractible covers)}
Every ring \(R\) admits a faithfully flat ind-étale algebra \(D\) that is w-contractible.
\end{bpstatementit}

% proof: proof/theorem-existence-contractible-covers
\begin{proof}[{Proof}]\phantomsection\label{bp38}%
The construction is that of [\hyperref[bp191]{BS15}, Theorem 2.4.9] and [Stacks, \href{https://stacks.math.columbia.edu/tag/0983}{0983}]. We first make the local rings strictly henselian, and then replace the profinite space of components by an extremally disconnected one surjecting onto it, which does not change the local rings.

In the first stage we construct a faithfully flat ind-étale \(R\)-algebra \(C\) that is w-strictly local. To do so we first replace \(R\) by an ind-Zariski algebra whose spectrum is w-local, and then take a filtered colimit of étale neighbourhoods of the closed points, which makes the local rings at the closed points strictly henselian. This is the construction of simultaneous strict localizations in [Stacks, \href{https://stacks.math.columbia.edu/tag/097X}{097X}]. Each step is a filtered colimit of étale algebras, so \(C\) is ind-étale over \(R\) by Lemma \hyperref[bp10]{2.3}.

For the second stage, choose a surjection \(S\to\pi_0(\operatorname{Spec}(C))\) with \(S\) extremally disconnected using Lemma \hyperref[bp31]{2.11}, and realize this map by an ind-Zariski \(C\)-algebra \(D\) such that

\[
|\operatorname{Spec}(D)|
\simeq |\operatorname{Spec}(C)|
\times_{\pi_0(\operatorname{Spec}(C))}S
\]

as topological spaces. Concretely, write \(S\) as a limit of finite quotients, use finite products of \(C\), restrict to the relevant closed sets of components by idempotent localizations, and take a filtered colimit; see [Stacks, \href{https://stacks.math.columbia.edu/tag/097D}{097D}]. This construction identifies local rings, is surjective when \(S\) is, and preserves w-locality. Hence \(C\to D\) is faithfully flat and the local rings at closed points of \(D\) remain strictly henselian. The space of components of \(\operatorname{Spec}(D)\) is \(S\), so \(D\) is w-contractible, and the composite \(R\to C\to D\) is ind-étale and faithfully flat.
\end{proof}

% text: prose-02-14
\phantomsection\label{bp39}%
To compute on \(X_{\mathrm{pro}\acute{e}t}\) we want the limits of étale affines to form a basis (Proposition \hyperref[bp44]{2.15}), so that étale sheaves can be evaluated on every object by a colimit. Objects of \(X_{\mathrm{pro}\acute{e}t}\) are only weakly étale, while the covers above are ind-étale, and the next theorem bridges the two.

% object: theorem-pro-etale-refinements-weakly-etale (blueprint/chap-The-pro-etale-site/sec-w-local-and-w-contractible-rings/theorem-pro-etale-refinements-weakly-etale/_section.md)
\phantomsection\label{bp40}%
\begin{bpstatementit}{Theorem~2.14}{ (Pro-étale refinements of weakly étale maps)}
If \(R\to B\) is weakly étale, there is a faithfully flat ind-étale map \(B\to B'\) such that \(R\to B'\) is ind-étale.
\end{bpstatementit}

% proof: proof/theorem-pro-etale-refinements-weakly-etale
\begin{proof}[{Proof}]\phantomsection\label{bp41}%
Weakly étale and ind-étale maps differ only in how local rings are glued. Namely, after strict localization of the base a weakly étale map identifies local rings, that is, for every prime \(\mathfrak q\) of \(B\) lying over \(\mathfrak p\subset R\) the induced map of local rings \(R_{\mathfrak p}\to B_{\mathfrak q}\) is an isomorphism. Such a map becomes ind-Zariski after refinement.

Applying the construction of simultaneous strict localizations to the pair \(R\to B\) we obtain compatible ind-étale refinements \(R\to C\) and \(B\to D\), with both spectra w-local and with closed points of \(D\) mapping to closed points of \(C\). The map \(C\to D\) is still weakly étale by Lemma \hyperref[bp10]{2.3}, and the local rings of \(C\) at closed points are strictly henselian. Olivier's rigidity theorem says that a weakly étale local map out of a strictly henselian local ring is an isomorphism; see [Stacks, \href{https://stacks.math.columbia.edu/tag/092Z}{092Z}]. Applying it at the closed points of \(D\), which lie over closed points of \(C\), we get isomorphisms of local rings there. Every point of the w-local \(\operatorname{Spec}(D)\) is a generalization of a closed point, and its local ring is a localization of the local ring at that closed point. Hence \(C\to D\) identifies all local rings. For a ring map that identifies local rings there is a faithfully flat ind-Zariski refinement of the target over which the map becomes ind-Zariski [Stacks, \href{https://stacks.math.columbia.edu/tag/097G}{097G}]. Applied to \(C\to D\), this gives a faithfully flat ind-Zariski map \(D\to D'\) such that \(C\to D'\) is ind-Zariski, and we take \(B'=D'\).

The algebra \(B'\) has the two required properties. The map \(B\to B'\) is the composite \(B\to D\to D'\) of two faithfully flat ind-étale maps, and \(R\to B'\) is the composite \(R\to C\to D'\) of an ind-étale and an ind-Zariski map. Both composites are ind-étale by Lemma \hyperref[bp10]{2.3}. For details see [Stacks, \href{https://stacks.math.columbia.edu/tag/097Z}{097Z}].
\end{proof}


% object: sec-Affine-bases (blueprint/chap-The-pro-etale-site/sec-Affine-bases/_section.md), content hidden
\subsection*{\texorpdfstring{2.3\quad Affine bases}{2.3 Affine bases}}
\phantomsection\label{bp42}%
\addcontentsline{toc}{subsection}{\texorpdfstring{2.3\quad Affine bases}{2.3 Affine bases}}
% text: prose-02-15
\phantomsection\label{bp43}%
To compute with sheaves on \(X_{\mathrm{pro}\acute{e}t}\) it suffices to evaluate on two kinds of affines. Every object of \(X_{\mathrm{pro}\acute{e}t}\) is covered by pro-étale affines (Theorem \hyperref[bp40]{2.14}), and these are in turn covered by w-contractible affines (Theorem \hyperref[bp37]{2.13}), over which taking sections is exact.

The two kinds of affines are used for different purposes. On a pro-étale affine \(U=\varprojlim_iU_i\), a sheaf pulled back from \(X_{\acute{e}t}\) can be evaluated by a colimit, and for this reason the comparison theorems of Section \hyperref[bp87]{2.7} are proved on such affines. On a w-contractible affine, by contrast, every epimorphism of sheaves is surjective on sections, so the homological results of the next three sections are proved there.

In Lemma \hyperref[bp47]{2.16} we turn the splitting of ring maps in Proposition \hyperref[bp34]{2.12} into a statement about sheaves. The idea is this. A finite cover of a w-contractible affine has a section. Such a section amounts to a decomposition of the affine into finitely many clopen pieces, each mapping into one member of the cover. Local lifts along an epimorphism can then be glued over these pieces without any further refinement.

% object: proposition-pro-etale-affines-form-basis (blueprint/chap-The-pro-etale-site/sec-Affine-bases/proposition-pro-etale-affines-form-basis/_section.md)
\phantomsection\label{bp44}%
\begin{bpstatementit}{Proposition~2.15}{ (Pro-étale affines form a basis)}
Every object of \(X_{\mathrm{pro}\acute{e}t}\) has a pro-étale covering by pro-étale affines over \(X\). These form a full subsite whose category of sheaves, of sets or of abelian groups, is equivalent to \(\Sh(X_{\mathrm{pro}\acute{e}t})\) by restriction.
\end{bpstatementit}

% proof: proof/proposition-pro-etale-affines-form-basis
\begin{proof}[{Proof}]\phantomsection\label{bp45}%
The first assertion is the refinement theorem applied chart by chart, and the second is then formal.

Given an object \(U\) of \(X_{\mathrm{pro}\acute{e}t}\), cover \(X\) by affine opens and cover their inverse images in \(U\) by affine opens. These Zariski coverings are pro-étale coverings, so it suffices to cover each chart. On such a chart the map is a weakly étale ring map \(R\to B\). By Theorem \hyperref[bp40]{2.14} there is a single faithfully flat cover \(\operatorname{Spec}(B')\to\operatorname{Spec}(B)\) such that \(\operatorname{Spec}(B')\) is a limit of affine schemes étale over \(\operatorname{Spec}(R)\), hence over \(X\). The equivalence of categories of sheaves follows from the comparison lemma for sites; see [\hyperref[bp191]{BS15}, Lemma 4.2.4].
\end{proof}

% text: prose-02-16
\phantomsection\label{bp46}%
Over \(X=\operatorname{Spec}\mathbf F_q\), for example, the weakly étale \(X\)-scheme \(\operatorname{Spec}\overline{\mathbf F}_q\) is already a pro-étale affine, and \(X\) itself is covered by it. Applying Theorem \hyperref[bp37]{2.13} to the charts of such a basis, we pass from pro-étale affines to w-contractible ones. The next two results show that the spectra of w-contractible rings have the splitting property for sheaves and that they cover every object.

% object: lemma-w-contractible-affines-projective (blueprint/chap-The-pro-etale-site/sec-Affine-bases/lemma-w-contractible-affines-projective/_section.md)
\phantomsection\label{bp47}%
\begin{bpstatementit}{Lemma~2.16}{ (Sections of epimorphisms onto w-contractible affines)}
Let \(U=\operatorname{Spec}(D)\) be an object of \(X_{\mathrm{pro}\acute{e}t}\) with \(D\) w-contractible. Then every epimorphism of sheaves \(F\to h_U\) has a section. Equivalently, \(h_U\) is projective in \(\Sh(X_{\mathrm{pro}\acute{e}t})\). More generally, for an object \(W\) of \(X_{\mathrm{pro}\acute{e}t}\) the following are equivalent: every epimorphism of sheaves of sets \(F\to h_W\) admits a section; \(h_W\) is projective in \(\Sh(X_{\mathrm{pro}\acute{e}t})\); every epimorphism \(F\to G\) induces a surjection \(F(W)\to G(W)\).
\end{bpstatementit}

% proof: proof/lemma-w-contractible-affines-projective
\begin{proof}[{Proof}]\phantomsection\label{bp48}%
For the equivalence of the three conditions on an object \(W\): by the Yoneda lemma, \(h_W\) is projective exactly when every epimorphism \(F\to G\) induces a surjection \(F(W)=\Hom(h_W,F)\to\Hom(h_W,G)=G(W)\). If every epimorphism onto \(h_W\) splits, pull an epimorphism \(F\to G\) back along a section \(h_W\to G\); the resulting epimorphism onto \(h_W\) splits, so the section lifts to \(F\). Conversely, if every epimorphism is surjective on sections over \(W\), then \(\mathrm{id}_{h_W}\) lifts along any epimorphism onto \(h_W\).

To see that covers of \(U\) split, note that by Theorem \hyperref[bp40]{2.14} a faithfully flat weakly étale \(D\)-algebra can be refined to a faithfully flat ind-étale \(D\)-algebra. This refinement splits by Proposition \hyperref[bp34]{2.12}, and hence so does the original map. Now, given a cover of the affine \(U\), we refine it by a finite affine cover \(\{V_j\to U\}\). Its disjoint union \(\coprod_jV_j\) is a single faithfully flat weakly étale affine cover, which therefore has a section \(s:U\to\coprod_jV_j\). The sets \(s^{-1}(V_j)\) then form a finite clopen partition of \(U\), and on \(s^{-1}(V_j)\) the inclusion into \(U\) factors through \(V_j\).

If \(F\to h_U\) is an epimorphism, then by the definition of an epimorphism of sheaves \(\mathrm{id}_U\) lifts after passing to a covering of \(U\). We refine this covering by the clopen partition just constructed, so that on each piece \(s^{-1}(V_j)\) we have a lift. Since the pieces are disjoint and cover \(U\), the sheaf condition glues these lifts with no compatibility to check, and the result is a global lift of \(\mathrm{id}_U\), that is, a section. The same argument, after pulling back an arbitrary epimorphism along a map from \(h_U\), proves projectivity. See [Stacks, \href{https://stacks.math.columbia.edu/tag/098F}{098F}, \href{https://stacks.math.columbia.edu/tag/0990}{0990}].
\end{proof}

% object: theorem-enough-affine-weakly-contractible-objects (blueprint/chap-The-pro-etale-site/sec-Affine-bases/theorem-enough-affine-weakly-contractible-objects/_section.md)
\phantomsection\label{bp49}%
\begin{bpstatementit}{Theorem~2.17}{ (Enough affine weakly contractible objects)}
Every object of \(X_{\mathrm{pro}\acute{e}t}\) has a covering by affines \(U=\operatorname{Spec}(D)\) with \(D\) w-contractible.
\end{bpstatementit}

% proof: proof/theorem-enough-affine-weakly-contractible-objects
\begin{proof}[{Proof}]\phantomsection\label{bp50}%
Cover the given object by pro-étale affines \(\operatorname{Spec}(B)\) (Proposition \hyperref[bp44]{2.15}), and for each of them choose, by Theorem \hyperref[bp37]{2.13}, a faithfully flat ind-étale \(B\)-algebra \(D\) that is w-contractible. Then \(\operatorname{Spec}(D)\to\operatorname{Spec}(B)\) is weakly étale by Proposition \hyperref[bp7]{2.2}, so \(\operatorname{Spec}(D)\) is an object of \(X_{\mathrm{pro}\acute{e}t}\) by Lemma \hyperref[bp10]{2.3}, and as a faithfully flat map of affines it is a covering. Composing these coverings gives the result.
\end{proof}


% object: sec-Weakly-contractible-objects (blueprint/chap-The-pro-etale-site/sec-Weakly-contractible-objects/_section.md), content hidden
\subsection*{\texorpdfstring{2.4\quad Weakly contractible objects}{2.4 Weakly contractible objects}}
\phantomsection\label{bp51}%
\addcontentsline{toc}{subsection}{\texorpdfstring{2.4\quad Weakly contractible objects}{2.4 Weakly contractible objects}}
% text: prose-02-17
\phantomsection\label{bp52}%
On a topological space we check exactness at points, and on the étale site at geometric points. On \(X_{\mathrm{pro}\acute{e}t}\), instead, we check exactness, and compute limits and colimits, by evaluating on weakly contractible objects. Over these objects we can compute in the ordinary derived categories of sheaves of modules by taking sections. Until Section \hyperref[bp87]{2.7}, \(\mathcal R\) denotes a sheaf of rings on \(X_{\mathrm{pro}\acute{e}t}\), for instance the pullback of a condensed ring (Section \hyperref[bp147]{2.12}).

The main result is Corollary \hyperref[bp57]{2.20}. On a weakly contractible affine \(W\), taking sections is exact and commutes with all limits and colimits of module sheaves. Moreover, these functors, for all such \(W\), are jointly conservative.

Stalk functors behave worse. Indeed, a stalk functor is exact and commutes with colimits, but not with infinite products. Moreover, products of sheaves on \(X_{\acute{e}t}\) need not be exact either, because a section of a product of surjections lifts in each factor only after its own refinement of the cover, and infinitely many refinements cannot in general be made at once. Evaluation at \(W\) avoids both problems: it commutes with products because limits of sheaves are computed on sections, and it is exact because covers of \(W\) split.

In the simplest case, \(X\) a geometric point, the site is that of profinite sets (Proposition \hyperref[bp151]{2.50}, proved in Subsection \hyperref[bp149]{2.12.1} below). The weakly contractible affines are then the extremally disconnected ones, among them \(\beta(T^\delta)\) for a set \(T\). Indeed, a surjection of profinite sets onto such an \(S\) has a section by Gleason's theorem.

% object: definition-weak-contractibility (blueprint/chap-The-pro-etale-site/sec-Weakly-contractible-objects/definition-weak-contractibility/_section.md)
\phantomsection\label{bp53}%
\begin{bpstatement}{Definition~2.18}{ (Weak contractibility)}
An object \(W\) of \(X_{\mathrm{pro}\acute{e}t}\) is \emph{weakly contractible} if every epimorphism of sheaves of sets \(F\to h_W\) admits a section, where \(h_W\) is the representable sheaf associated to \(W\). By Lemma \hyperref[bp47]{2.16}, equivalently \(h_W\) is projective in \(\Sh(X_{\mathrm{pro}\acute{e}t})\), or every epimorphism \(F\to G\) induces a surjection \(F(W)\to G(W)\).
\end{bpstatement}

% text: prose-02-18
\phantomsection\label{bp54}%
This is the categorical definition of [\hyperref[bp191]{BS15}, §3.2].

% object: example-weak-contractibility-point (blueprint/chap-The-pro-etale-site/sec-Weakly-contractible-objects/example-weak-contractibility-point/_section.md)
\phantomsection\label{bp55}%
\begin{bpstatement}{Example~2.19}{ (Weak contractibility over a point)}
Let \(X\) be a geometric point, so that the affine objects of \(X_{\mathrm{pro}\acute{e}t}\) are the profinite sets (Proposition \hyperref[bp151]{2.50} below). A continuous surjection \(S\to T\) of profinite sets is a covering, so \(h_S\to h_T\) is an epimorphism, and by the Yoneda lemma a section of it is a continuous map \(T\to S\) that is a section of the surjection. If \(T\) is finite, such a section exists: choose a preimage of each point; it is continuous because \(T\) is discrete. If \(T=\{1,2,3,\ldots\}\cup\{\infty\}\) is a convergent sequence, it need not exist. Let \(S=T_0\amalg T_1\), where \(T_0\) is the convergent sequence \(\{2,4,6,\ldots\}\cup\{\infty_0\}\) and \(T_1\) is \(\{1,3,5,\ldots\}\cup\{\infty_1\}\), and map \(S\to T\) by the identity on integers and \(\infty_0,\infty_1\mapsto\infty\). A section must send each integer \(n\) to its only preimage \(n\), and then the even integers converge to \(\infty_0\) while the odd ones converge to \(\infty_1\), so no value at \(\infty\) makes the section continuous. Hence \(T\) is not weakly contractible. In agreement with Gleason's theorem, \(T\) is not extremally disconnected: the set of odd integers is open, and its closure, which contains \(\infty\), is not open, because every neighbourhood of \(\infty\) contains even integers.
\end{bpstatement}

% text: prose-02-19
\phantomsection\label{bp56}%
Note that the splitting condition concerns a covering family as a whole. If \(W\) is disconnected, a lift of \(\mathrm{id}_{h_W}\) along a covering family may use different members of the family on different clopen pieces of \(W\).

By Lemma \hyperref[bp47]{2.16}, the spectrum of a w-contractible ring is weakly contractible. By Theorem \hyperref[bp49]{2.17}, such affines cover every object of \(X_{\mathrm{pro}\acute{e}t}\). For abelian sheaves we use these two facts in the following form.

% object: corollary-evaluation-weakly-contractible-objects (blueprint/chap-The-pro-etale-site/sec-Weakly-contractible-objects/corollary-evaluation-weakly-contractible-objects/_section.md)
\phantomsection\label{bp57}%
\begin{bpstatementit}{Corollary~2.20}{ (Evaluation on weakly contractible objects)}
Let \(W\) be a weakly contractible object of \(X_{\mathrm{pro}\acute{e}t}\), for instance \(W=\operatorname{Spec}(D)\) with \(D\) w-contractible.

\begin{enumerate}
\item The functor \(F\mapsto F(W)\) from abelian sheaves on \(X_{\mathrm{pro}\acute{e}t}\) to abelian groups is exact. Hence \(H^q(W,F)=0\) for \(q>0\), and \(R\Gamma(W,K)\simeq K(W)\) for every complex \(K\), with the right side computed term by term.
\item If \(W\) is affine, evaluation at \(W\) preserves all limits and all colimits of \(\mathcal R\)-modules.
\item Exactness of complexes and isomorphisms of \(\mathcal R\)-module sheaves can be tested by evaluation on all affine weakly contractible objects of \(X_{\mathrm{pro}\acute{e}t}\).
\end{enumerate}
\end{bpstatementit}

% proof: proof/corollary-evaluation-weakly-contractible-objects
\begin{proof}[{Proof}]\phantomsection\label{bp58}%
All three parts follow from the fact that every epimorphism onto a sheaf is surjective on sections over \(W\), together with the fact that such \(W\) cover every object. Exactness and conservativity follow directly, and for colimits we need in addition that \(W\) is quasi-compact.

(1) Taking sections is left exact for any object of a site, since it preserves kernels. By weak contractibility every epimorphism of sheaves is surjective on sections over \(W\), so the sections functor is also right exact. An exact functor carries acyclic complexes to acyclic complexes, and so it agrees with its derived functor, also for unbounded complexes. In particular \(H^q(W,F)=0\) for \(q>0\).

(2) Limits of sheaves are computed on sections, so only colimits need an argument. We first consider a direct sum \(\bigoplus_iM_i\). The direct sum sheaf is the sheafification of the presheaf direct sum, so a section over \(W\) is, locally on \(W\), a finite sum of sections of the \(M_i\). Since \(W\) is affine, hence quasi-compact as an object of the site, a finite covering suffices, and a single finite set of indices \(i\) serves for all members of this covering. By the sheaf condition, the section is then a section of the corresponding finite subsum, and finite sums commute with evaluation. It follows that evaluation preserves arbitrary direct sums [Stacks, \href{https://stacks.math.columbia.edu/tag/0935}{0935}]. An arbitrary colimit of modules is the cokernel of a map between direct sums, and evaluation preserves cokernels by (1). Hence it preserves all colimits.

(3) Suppose that a sequence \(F'\to F\to F''\) with composite zero becomes exact on all affine weakly contractible \(W\). A local section of the kernel of \(F\to F''\) is, after restriction to a covering by such \(W\), the image of a section of \(F'\), because these objects cover every object. Hence it lies in the image sheaf, and the sequence is exact. Applied to kernels and cokernels of a map, this criterion shows that a map which is an isomorphism on all such \(W\) is an isomorphism, and applied in each degree it gives the statement for complexes; see [Stacks, \href{https://stacks.math.columbia.edu/tag/0945}{0945}].
\end{proof}

% text: prose-02-20
\phantomsection\label{bp59}%
The quasi-compactness used here is quasi-compactness as an object of the site: coverings admit finite refinements. Affine objects of \(X_{\mathrm{pro}\acute{e}t}\) have this property. The affine weakly contractible objects are quasi-compact, quasi-separated and projective, and they cover every object.

Cohomology on other objects need not vanish. For instance, a fibre product of two weakly contractible affines need not be weakly contractible, and for this reason the Čech complex of a cover by them does not in general compute cohomology.


% object: sec-Repleteness-and-limits (blueprint/chap-The-pro-etale-site/sec-Repleteness-and-limits/_section.md), content hidden
\subsection*{\texorpdfstring{2.5\quad Repleteness and limits}{2.5 Repleteness and limits}}
\phantomsection\label{bp60}%
\addcontentsline{toc}{subsection}{\texorpdfstring{2.5\quad Repleteness and limits}{2.5 Repleteness and limits}}
% text: prose-02-21
\phantomsection\label{bp61}%
We need inverse limits to pass from torsion to adic coefficients, for example to form \(\varprojlim_n\mathbf Z/\ell^n\), and the first use of weakly contractible objects is to show that such limits are well behaved. In a general topos a limit of epimorphisms need not be an epimorphism, because the local lifts at each stage may require ever finer covers. On \(X_{\mathrm{pro}\acute{e}t}\), however, this does not happen, since surjectivity and exactness can be checked on sections over weakly contractible objects, and there no refinement is needed. As a result, the limit of a tower of epimorphisms maps epimorphically onto each term (Proposition \hyperref[bp63]{2.22}). For the same reason products are exact (Proposition \hyperref[bp66]{2.23}).

Exactness of products means that derived limits can be computed as for abelian groups. Namely, the derived limit of a tower is the fibre of \(1-T\) on the product. Its cohomology is therefore given by the Milnor sequence of Proposition \hyperref[bp69]{2.24}, and \(R^p\varprojlim\) vanishes for \(p\geq2\). For a tower of epimorphisms, such as \((\mathbf Z/\ell^n)_X\), moreover \(R^1\varprojlim\) vanishes, and the derived limit is the ordinary limit. In the next section we also deduce left completeness from the Milnor sequence. With these facts we compute the cohomology of \(\mathbf Z_\ell\) as continuous cohomology (Proposition \hyperref[bp170]{2.55}).

% object: definition-repleteness (blueprint/chap-The-pro-etale-site/sec-Repleteness-and-limits/definition-repleteness/_section.md)
\phantomsection\label{bp62}%
\begin{bpstatement}{Definition~2.21}{ (Repleteness)}
A topos \(\mathcal T\) is \emph{replete} if, for every inverse sequence \(\cdots\to F_2\to F_1\to F_0\) of sheaves of sets with epimorphic transition maps, the projection \(\varprojlim_nF_n\to F_m\) is an epimorphism for every \(m\).
\end{bpstatement}

% object: proposition-repleteness-pro-etale-topos (blueprint/chap-The-pro-etale-site/sec-Repleteness-and-limits/proposition-repleteness-pro-etale-topos/_section.md)
\phantomsection\label{bp63}%
\begin{bpstatementit}{Proposition~2.22}{ (Repleteness of the pro-étale topos)}
For every scheme \(X\), the topos \(\Sh(X_{\mathrm{pro}\acute{e}t})\) is replete.
\end{bpstatementit}

% proof: proof/proposition-repleteness-pro-etale-topos
\begin{proof}[{Proof}]\phantomsection\label{bp64}%
Evaluate the sequence on a weakly contractible \(W\). Since every epimorphism of sheaves is surjective on sections over \(W\), the transition maps \(F_{n+1}(W)\to F_n(W)\) are surjective. Given a section at stage \(m\), we choose lifts successively at stages \(m+1,m+2,\ldots\) (by countable choice), while the earlier stages are determined by the transition maps. This gives a compatible section of \(\varprojlim_nF_n(W)\). Since limits are computed on sections, the projection is surjective on \(W\).

An arbitrary \(U\) is covered by affine weakly contractible objects, so every section of \(F_m\) over \(U\) lifts to \(\varprojlim_nF_n\) on each member of such a cover. Hence the projection is an epimorphism of sheaves; see [\hyperref[bp191]{BS15}, §§3.1–3.2].
\end{proof}

% text: prose-02-22
\phantomsection\label{bp65}%
For sheaves of modules the same lifting argument, applied coordinatewise rather than along a tower, shows that products are exact. This lets us compute derived limits by termwise products of complexes.

% object: proposition-exact-products (blueprint/chap-The-pro-etale-site/sec-Repleteness-and-limits/proposition-exact-products/_section.md)
\phantomsection\label{bp66}%
\begin{bpstatementit}{Proposition~2.23}{ (Exact products)}
Arbitrary set-indexed products are exact in \(\operatorname{Mod}(X_{\mathrm{pro}\acute{e}t},\mathcal R)\). Products in \(D(X_{\mathrm{pro}\acute{e}t},\mathcal R)\) are computed by taking termwise products of representative complexes, and their cohomology sheaves are the products of the cohomology sheaves.
\end{bpstatementit}

% proof: proof/proposition-exact-products
\begin{proof}[{Proof}]\phantomsection\label{bp67}%
By Corollary corollary-evaluation-weakly-contractible-objects, exactness of sheaves can be checked on affine weakly contractible \(W\), and by part (1) evaluation at such \(W\) is exact. Since limits are computed on sections, a product of short exact sequences therefore evaluates on \(W\) to a product of short exact sequences of \(\mathcal R(W)\)-modules, which is exact because by the axiom of choice we may lift a tuple coordinate by coordinate. This proves the first assertion.

The derived product is the termwise product of K-injective representatives (complexes \(I\) such that maps from acyclic complexes to \(I\) are null-homotopic, the unbounded analogue of bounded-below injective complexes). Since products are exact, a product of quasi-isomorphisms is a quasi-isomorphism: its cone is the product of the cones, which is acyclic. We may therefore replace the K-injective representatives by any quasi-isomorphic ones. Hence termwise products compute the product in \(D(X_{\mathrm{pro}\acute{e}t},\mathcal R)\), and they commute with cohomology because products are exact; see [Stacks, \href{https://stacks.math.columbia.edu/tag/0947}{0947}, \href{https://stacks.math.columbia.edu/tag/08TB}{08TB}].
\end{proof}

% text: prose-02-23
\phantomsection\label{bp68}%
Derived limits of countable towers have an explicit model. For a countable tower of modules, write \(t_n:M_{n+1}\to M_n\) and let \(T:\prod_nM_n\to\prod_nM_n\) be \(T((m_n)_n)=(t_n(m_{n+1}))_n\). For a tower \((K_n)\) in \(D(X_{\mathrm{pro}\acute{e}t},\mathcal R)\), represented by a tower of complexes \(K_{n+1}\to K_n\), the same formula defines \(T\) termwise, and the derived inverse limit is the homotopy fibre of \(1-T\)

\[
R\varprojlim_n K_n
=
\fib\Bigl(
\prod_n K_n\xrightarrow{\,1-T\,}\prod_n K_n
\Bigr)=\Cone(1-T)[-1].
\tag{2.1}
\]

The cone is formed on the level of complexes, so that the construction is functorial in the tower. By Proposition \hyperref[bp66]{2.23} the termwise products are the derived products, and so, by [Stacks, \href{https://stacks.math.columbia.edu/tag/0940}{0940}], this cone agrees with the derived inverse limit.

% object: proposition-milnor-sequence (blueprint/chap-The-pro-etale-site/sec-Repleteness-and-limits/proposition-milnor-sequence/_section.md)
\phantomsection\label{bp69}%
\begin{bpstatementit}{Proposition~2.24}{ (The Milnor sequence)}
For a countable tower of \(\mathcal R\)-modules, \(R^p\varprojlim=0\) for \(p>1\), and

\[
R^1\varprojlim_n M_n
=\coker\left(\prod_n M_n\xrightarrow{\,1-T\,}\prod_n M_n\right).
\]

For a tower \((K_n)\) of complexes there is a short exact sequence of sheaves

\[
0\longrightarrow R^1\varprojlim_n H^{q-1}(K_n)
\longrightarrow H^q(R\varprojlim_nK_n)
\longrightarrow \varprojlim_n H^q(K_n)
\longrightarrow 0.
\tag{2.2}
\]

If the transitions of a tower of modules are epimorphisms, its derived limit is its ordinary limit in degree zero.
\end{bpstatementit}

% proof: proof/proposition-milnor-sequence
\begin{proof}[{Proof}]\phantomsection\label{bp70}%
Everything follows from the triangle (2.1). Indeed, for a tower of modules, \(R\varprojlim_nM_n\) is the fibre of a map between two sheaves in degree \(0\), so it has cohomology only in degrees \(0\) and \(1\), namely the kernel and the cokernel of \(1-T\). Here the kernel is \(\varprojlim_nM_n\), and the cokernel is \(R^1\varprojlim_nM_n\), which proves the first two assertions. Similarly, for a tower of complexes, the long exact cohomology sequence of (2.1) reads

\[
\prod_nH^{q-1}(K_n)\xrightarrow{1-T}\prod_nH^{q-1}(K_n)
\to H^q(R\varprojlim_nK_n)\to\prod_nH^q(K_n)\xrightarrow{1-T}\prod_nH^q(K_n),
\]

where we used that by Proposition \hyperref[bp66]{2.23} the cohomology of each product is the product of the cohomology sheaves. The cokernel on the left is \(R^1\varprojlim_nH^{q-1}(K_n)\) and the kernel on the right is \(\varprojlim_nH^q(K_n)\), which gives (2.2).

Suppose now that the transition maps of a tower of modules are epimorphisms. We show that \(1-T\) is an epimorphism; then \(R^1\varprojlim_nM_n=0\) and the derived limit is \(\varprojlim_nM_n\). By Corollary \hyperref[bp57]{2.20} it suffices to check surjectivity on sections over every weakly contractible \(W\), where the maps \(t_n\) are surjective. Given \((a_n)\), choose \(m_0\) arbitrarily and then, recursively, \(m_{n+1}\) with \(t_n(m_{n+1})=m_n-a_n\). Then \((1-T)(m_n)_n=(a_n)_n\). See [Stacks, \href{https://stacks.math.columbia.edu/tag/0947}{0947}].
\end{proof}

% object: example-nonzero-rlim-one (blueprint/chap-The-pro-etale-site/sec-Repleteness-and-limits/example-nonzero-rlim-one/_section.md)
\phantomsection\label{bp71}%
\begin{bpstatement}{Example~2.25}{ (A nonzero $R^1\varprojlim$)}
Without surjective transition maps the derived limit can be nonzero in degree one. To see this, work over a geometric point and evaluate on the final object, which is weakly contractible, so that we may compute with abelian groups. Let \(M_n=\mathbf Z\) for all \(n\), with transition maps multiplication by \(\ell\). Then the maps \(M_n\to\mathbf Z\), \(m\mapsto\ell^nm\), form a map to the constant tower \(\mathbf Z\) whose cokernel is the tower \((\mathbf Z/\ell^n)\). The constant tower has limit \(\mathbf Z\) and no \(R^1\varprojlim\), and \(\varprojlim_nM_n=0\), because an element of \(M_0\) in the image of every \(M_n\) is divisible by every power of \(\ell\). The long exact sequence of \(R\varprojlim\) therefore reads

\[
0\longrightarrow\mathbf Z\longrightarrow\mathbf Z_\ell\longrightarrow R^1\varprojlim_nM_n\longrightarrow0,
\]

so the sections of \(R^1\varprojlim_nM_n\) over the point form \(\mathbf Z_\ell/\mathbf Z\), an uncountable group. If the \(M_n\) are the groups \(H^{q-1}\) of a tower of complexes, this group appears in \(H^q\) of the derived limit by (2.2).
\end{bpstatement}

% text: prose-02-24
\phantomsection\label{bp72}%
More generally, the \(R^1\varprojlim\) term vanishes if, on every affine weakly contractible \(W\), the tower of section modules satisfies the Mittag–Leffler condition. To see this, note that evaluation on \(W\) preserves cokernels (Corollary \hyperref[bp57]{2.20}). So the cokernel of \(1-T\) evaluated on \(W\) is the cokernel of \(1-T\) on the tower \((M_n(W))\), which vanishes by the Mittag–Leffler lemma [Stacks, \href{https://stacks.math.columbia.edu/tag/07KV}{07KV}]. As evaluation on these \(W\) detects vanishing, the \(R^1\varprojlim\) term vanishes.


% object: sec-Left-completeness-and-hypercovers (blueprint/chap-The-pro-etale-site/sec-Left-completeness-and-hypercovers/_section.md), content hidden
\subsection*{\texorpdfstring{2.6\quad Left completeness and hypercovers}{2.6 Left completeness and hypercovers}}
\phantomsection\label{bp73}%
\addcontentsline{toc}{subsection}{\texorpdfstring{2.6\quad Left completeness and hypercovers}{2.6 Left completeness and hypercovers}}
% text: prose-02-25
\phantomsection\label{bp74}%
An unbounded complex on \(X_{\mathrm{pro}\acute{e}t}\) is determined by its bounded-below truncations, and its cohomology is computed from theirs by a derived limit. On the étale site of a scheme of infinite cohomological dimension this can fail, whereas on \(X_{\mathrm{pro}\acute{e}t}\) it holds because products are exact, so that the Milnor sequence applies.

Left completeness (Proposition \hyperref[bp75]{2.26}) lets us treat an unbounded complex as the derived limit of bounded below ones, to which injective resolutions and the comparison theorems apply. In this way we prove the unbounded comparison of Proposition \hyperref[bp144]{2.49}.

For bounded below complexes we also give an explicit model of cohomology. On a quasi-compact quasi-separated object \(U\), Proposition \hyperref[bp81]{2.28} provides a hypercover all of whose terms \(W_r\) are affine weakly contractible. Since these terms have no higher cohomology, the hypercover spectral sequence collapses, and \(R\Gamma(U,K)\) is the total complex of the sections \(K^b(W_r)\). This model has a consequence we use later. Namely, in each total degree only finitely many terms occur, and evaluation on each \(W_r\) commutes with filtered colimits. Hence cohomology commutes with filtered colimits of uniformly bounded below complexes (Corollary \hyperref[bp84]{2.29}). We use this to show that cohomology and direct images commute with inverting \(\ell\) and with extension of scalars to \(\overline{\mathbf Q}_\ell\), the colimit over the finite extensions \(E'\) of \(E\).

% object: proposition-left-completeness (blueprint/chap-The-pro-etale-site/sec-Left-completeness-and-hypercovers/proposition-left-completeness/_section.md)
\phantomsection\label{bp75}%
\begin{bpstatementit}{Proposition~2.26}{ (Left completeness)}
The standard \(t\)-structure on \(D(X_{\mathrm{pro}\acute{e}t},\mathcal R)\) is left complete: for every complex \(K\) the natural map

\[
K\longrightarrow R\varprojlim_{n\geq0}\tau_{\geq-n}K
\]

is an isomorphism.
\end{bpstatementit}

% proof: proof/proposition-left-completeness
\begin{proof}[{Proof}]\phantomsection\label{bp76}%
We compute the cohomology sheaves of the derived limit with the Milnor sequence and find that they are those of \(K\).

Represent \(K\) by a complex and form its canonical truncation tower, with maps \(K\to\tau_{\geq-n}K\) and \(\tau_{\geq-(n+1)}K\to\tau_{\geq-n}K\). For fixed \(q\), the sheaves \(H^q(\tau_{\geq-n}K)\) and \(H^{q-1}(\tau_{\geq-n}K)\) are eventually constant, with identity transition maps, equal to the corresponding cohomology sheaves of \(K\), so the \(R^1\) term in (2.2) vanishes and the limit term is \(H^q(K)\). As the natural map induces this identification in every degree, it is an isomorphism; see Stacks, [0947](https://stacks.math.columbia.edu/tag/0947).
\end{proof}

% text: prose-02-26
\phantomsection\label{bp77}%
To compute the cohomology of an unbounded complex from its truncations, we need derived functors on the whole unbounded derived category that commute with derived inverse limits.

% object: proposition-unbounded-functors-products (blueprint/chap-The-pro-etale-site/sec-Left-completeness-and-hypercovers/proposition-unbounded-functors-products/_section.md)
\phantomsection\label{bp78}%
\begin{bpstatementit}{Proposition~2.27}{ (Unbounded derived functors preserve products and limits)}
Let \(U\) be an object of \(X_{\mathrm{pro}\acute{e}t}\) and \(f:X\to Y\) a morphism of schemes. The functor \(R\Gamma(U,-)\), and the functor \(Rf_*\) from \(D(X_{\mathrm{pro}\acute{e}t},\mathcal R)\) to the derived category of sheaves of \(f_*\mathcal R\)-modules on \(Y_{\mathrm{pro}\acute{e}t}\), are defined on the unbounded derived categories. They preserve products and derived inverse limits.
\end{bpstatementit}

% proof: proof/proposition-unbounded-functors-products
\begin{proof}[{Proof}]\phantomsection\label{bp79}%
K-injective resolutions of arbitrary complexes exist in the Grothendieck category of \(\mathcal R\)-modules [Stacks, \href{https://stacks.math.columbia.edu/tag/079P}{079P}], which defines both functors. Both are right adjoints on derived categories: \(R\Gamma(U,-)\) is right adjoint to the functor sending a complex of abelian groups \(M\) to \(\mathcal R[h_U]\otimes^{\mathbf L}_{\mathbf Z}M\), where \(\mathcal R[h_U]\) is the free \(\mathcal R\)-module on \(h_U\), and \(Rf_*\) is right adjoint to the derived pullback of modules. Right adjoints preserve products. Since (2.1) is built from products and a cone, an exact functor preserving products also preserves derived inverse limits [Stacks, \href{https://stacks.math.columbia.edu/tag/0940}{0940}].
\end{proof}

% text: prose-02-27
\phantomsection\label{bp80}%
For bounded-below complexes, weakly contractible objects provide explicit cochain complexes computing cohomology.

% object: proposition-weakly-contractible-hypercovers (blueprint/chap-The-pro-etale-site/sec-Left-completeness-and-hypercovers/proposition-weakly-contractible-hypercovers/_section.md)
\phantomsection\label{bp81}%
\begin{bpstatementit}{Proposition~2.28}{ (Weakly contractible hypercovers)}
Let \(U\) be a quasi-compact, quasi-separated object of \(X_{\mathrm{pro}\acute{e}t}\). There is a hypercover \(W_\bullet\to U\) whose term in every simplicial degree is one affine weakly contractible object. If \(K\in D^+(X_{\mathrm{pro}\acute{e}t},\mathcal R)\) is represented by a complex \(K^\bullet\) vanishing below degree \(a\), then

\[
R\Gamma(U,K)
\simeq \operatorname{Tot}\bigl(K^b(W_r)\bigr)_{r\geq0,\ b\geq a}.
\tag{2.3}
\]

The cosimplicial differential is the alternating sum of pullbacks along face maps, combined with \((-1)^r\) times the differential of \(K^\bullet\). In total degree \(q\) this complex has the finite sum \(\bigoplus_{0\leq r\leq q-a}K^{q-r}(W_r)\).
\end{bpstatementit}

% proof: proof/proposition-weakly-contractible-hypercovers
\begin{proof}[{Proof}]\phantomsection\label{bp82}%
We build a hypercover whose terms have no higher cohomology, so that the hypercover spectral sequence collapses to the complex of sections. By quasi-compactness each term can be taken to be a single affine.

We start with a weakly contractible cover of \(U\), and at each stage we cover the matching object by weakly contractible affines, as for fibre products above. The matching objects are finite limits of quasi-compact, quasi-separated schemes, so they are again quasi-compact and quasi-separated and are covered by finitely many affine weakly contractible objects. Taking finite disjoint unions, we get a single affine term in each degree; see [Stacks, \href{https://stacks.math.columbia.edu/tag/09A1}{09A1}].

For a sheaf \(M\), the hypercover spectral sequence has \(E_1^{r,s}=H^s(W_r,M)\), and the rows with \(s>0\) vanish by Corollary corollary-evaluation-weakly-contractible-objects. The remaining row is the cosimplicial complex of sections. For bounded-below \(K^\bullet\), choose a bounded-below injective resolution \(K^\bullet\to I^\bullet\). Since evaluation on \(W_r\) is exact, \(K^\bullet(W_r)\to I^\bullet(W_r)\) is a quasi-isomorphism. After a shift, both double complexes lie in the first quadrant, so their total complexes are quasi-isomorphic, and hypercover descent for \(I^\bullet\) identifies the second with \(R\Gamma(U,K)\). See [Stacks, \href{https://stacks.math.columbia.edu/tag/09A2}{09A2}].
\end{proof}

% text: prose-02-28
\phantomsection\label{bp83}%
For a single sheaf this is the cochain complex \(M(W_0)\to M(W_1)\to M(W_2)\to\cdots\). It may be nonzero in every degree, so the hypercover gives no bound on cohomological dimension. For unbounded \(K\), by left completeness (Proposition \hyperref[bp75]{2.26}) and since \(R\Gamma(U,-)\) commutes with derived inverse limits, we get \(R\Gamma(U,K)\simeq
R\varprojlim_n R\Gamma(U,\tau_{\geq-n}K)\), and the hypercover complex computes each term.

% object: corollary-continuity-cohomology (blueprint/chap-The-pro-etale-site/sec-Left-completeness-and-hypercovers/corollary-continuity-cohomology/_section.md)
\phantomsection\label{bp84}%
\begin{bpstatementit}{Corollary~2.29}{ (Continuity of cohomology)}
Suppose \(U\) is quasi-compact and quasi-separated and \((K_i^\bullet)\) is a filtered diagram of complexes of \(\mathcal R\)-modules vanishing below one fixed degree \(a\). Then for every \(q\),

\[
\varinjlim_i H^q(U,K_i^\bullet)
\simeq H^q(U,\varinjlim_iK_i^\bullet).
\]

If \(f:X\to Y\) is quasi-compact and quasi-separated, the analogous map \(\varinjlim_iR^qf_*K_i^\bullet\to
R^qf_*(\varinjlim_iK_i^\bullet)\) is an isomorphism.
\end{bpstatementit}

% proof: proof/corollary-continuity-cohomology
\begin{proof}[{Proof}]\phantomsection\label{bp85}%
We compute both sides with one hypercover by weakly contractible affines, chosen independently of \(i\). Because of the common lower bound \(a\), each total degree \(q\) of (2.3) is a finite sum, of the terms \(K_i^{q-r}(W_r)\) with \(0\leq r\leq q-a\), and the range of \(r\) is the same for all \(i\).

By Corollary \hyperref[bp57]{2.20}, sections on each \(W_r\) commute with filtered colimits, hence so does the finite sum in each total degree of (2.3), and the hypercover complex of \(\varinjlim_iK_i^\bullet\) is the filtered colimit of the hypercover complexes of the \(K_i^\bullet\). The first assertion follows because filtered colimits of modules are exact and so commute with cohomology.

The second assertion reduces to the first. Recall that \(R^qf_*K\) is the sheaf associated to the presheaf \(V\mapsto H^q(X\times_YV,K)\), and that the affine objects \(V\) of \(Y_{\mathrm{pro}\acute{e}t}\) form a basis. For such \(V\) the object \(X\times_YV\) is quasi-compact and quasi-separated, so the first assertion applies to it, and sheafification commutes with colimits.
\end{proof}

% text: prose-02-29
\phantomsection\label{bp86}%
In particular, \(H^q(U,-)\) commutes with direct sums of sheaves of modules. Quasi-compactness of \(U\) is needed already for \(H^0\). Let \(X\) be a geometric point and \(U=\coprod_{n\in\mathbf N}X\), and consider the direct sum \(\bigoplus_{i\in\mathbf N}\mathbf Z\) of copies of the constant sheaf. Its sections over \(U\) are \(\prod_n\bigoplus_i\mathbf Z\), which contains the element whose \(n\)th component is the \(n\)th basis vector, whereas \(\bigoplus_i H^0(U,\mathbf Z)=\bigoplus_i\prod_n\mathbf Z\) consists of the elements with only finitely many nonzero indices \(i\).

For direct sums of unbounded complexes the common lower bound is replaced by uniform bounds on cohomological dimension (Proposition \hyperref[bp144]{2.49}).


% object: sec-Classical-sheaves (blueprint/chap-The-pro-etale-site/sec-Classical-sheaves/_section.md), content hidden
\subsection*{\texorpdfstring{2.7\quad Classical sheaves}{2.7 Classical sheaves}}
\phantomsection\label{bp87}%
\addcontentsline{toc}{subsection}{\texorpdfstring{2.7\quad Classical sheaves}{2.7 Classical sheaves}}
% text: prose-02-30
\phantomsection\label{bp88}%
From here until Section \hyperref[bp168]{2.13}, \(\Lambda\) denotes a discrete commutative ring and \(\Lambda_X\) its constant sheaf on \(X_{\mathrm{pro}\acute{e}t}\). To use the torsion arguments of étale cohomology, we must know which pro-étale sheaves come from étale sheaves. The main result answers this: étale sheaves embed fully faithfully, and the essential image consists of the sheaves whose sections on a limit \(\varprojlim U_i\) are the colimit of the sections on the \(U_i\). In particular, the reductions \(\mathbf Z/\ell^n\) lie in this image, while the \(\ell\)-adic coefficient sheaves themselves, such as \(\mathbf Z_{\ell,X}\), do not (Example \hyperref[bp95]{2.33}).

The colimit condition is forced on us by the way étale sheaves extend to limits. If \(F\) is an étale sheaf and \(\bar x\) a geometric point, the only reasonable value of the pullback of \(F\) on the strict henselization \(\operatorname{Spec}\mathcal O^{\mathrm{sh}}_{X,\bar x}=\varprojlim(U,\bar u)\) is \(\varinjlim F(U)=F_{\bar x}\). In this way stalks of étale sheaves become sections of pro-étale sheaves. Proposition \hyperref[bp90]{2.31} defines the pullback \(\nu_X^*\) by this rule, and Theorem \hyperref[bp93]{2.32} characterizes its image by it. The part that is not formal is cohomology, because pro-étale covers are not étale covers. In the next section we show that they create no new cohomology classes for classical sheaves.

% object: definition-classical-pro-etale-sheaves (blueprint/chap-The-pro-etale-site/sec-Classical-sheaves/definition-classical-pro-etale-sheaves/_section.md)
\phantomsection\label{bp89}%
\begin{bpstatement}{Definition~2.30}{ (Classical pro-étale sheaves)}
A sheaf \(F\) on \(X_{\mathrm{pro}\acute{e}t}\) is \emph{classical} if, for every pro-étale affine presentation \(U=\varprojlim_i U_i\) with \(U_i\) affine and étale over \(X\), its restriction maps induce a bijection

\[
\varinjlim_i F(U_i)\xrightarrow{\ \sim\ }F(U).
\tag{2.4}
\]

A complex is classical if all its cohomology sheaves are classical.
\end{bpstatement}

% object: proposition-comparison-morphism-inverse-image (blueprint/chap-The-pro-etale-site/sec-Classical-sheaves/proposition-comparison-morphism-inverse-image/_section.md)
\phantomsection\label{bp90}%
\begin{bpstatementit}{Proposition~2.31}{ (The comparison morphism and its inverse image)}
Inclusion of étale objects defines a morphism of topoi

\[
\nu_X:X_{\mathrm{pro}\acute{e}t}\longrightarrow X_{\acute{e}t}.
\]

Its direct image \(\nu_{X,*}:\Sh(X_{\mathrm{pro}\acute{e}t})\to\Sh(X_{\acute{e}t})\) is restriction to the étale subsite. Its inverse image \(\nu_X^*:\Sh(X_{\acute{e}t})\to\Sh(X_{\mathrm{pro}\acute{e}t})\) is exact and has the description

\[
(\nu_X^*F)(h:V\to X)
=\Gamma(V_{\acute{e}t},h_{\acute{e}t}^*F),
\tag{2.5}
\]

where \(h_{\acute{e}t}^*:\Sh(X_{\acute{e}t})\to\Sh(V_{\acute{e}t})\) is the étale inverse image along the weakly étale map \(h\). Both functors carry sheaves of \(\Lambda\)-modules to sheaves of \(\Lambda\)-modules, so they restrict to functors between \(\operatorname{Mod}(X_{\acute{e}t},\Lambda)\) and \(\operatorname{Mod}(X_{\mathrm{pro}\acute{e}t},\Lambda)\). In particular, for \(U=\varprojlim_iU_i\) as above, we have

\[
(\nu_X^*F)(U)=\varinjlim_i F(U_i).
\tag{2.6}
\]
\end{bpstatementit}

% proof: proof/proposition-comparison-morphism-inverse-image
\begin{proof}[{Proof}]\phantomsection\label{bp91}%
The morphism \(\nu_X\) is formal. For the formula, we write down the right-hand side of (2.5), check that it is a sheaf, and check that it has the universal property of the left adjoint of \(\nu_{X,*}\).

Since an étale morphism is weakly étale and an étale covering is a pro-étale covering, the inclusion is a continuous functor preserving finite limits, and it defines \(\nu_X\). For (2.5), we define the right-hand side for every weakly étale \(V\). Since pro-étale coverings are fpqc coverings, fpqc descent for étale sheaves shows that this rule is already a sheaf, so no sheafification is needed. If we pull back the restriction of a pro-étale sheaf \(G\) to \(X_{\acute{e}t}\), we get the counit \(\nu_X^*\nu_{X,*}G\to G\). On an étale object \(V\) the rule returns \(F(V)\), and this is the unit \(F\to\nu_{X,*}\nu_X^*F\). The two adjunction identities follow by restriction and pullback, so the rule is \(\nu_X^*\) [Stacks, \href{https://stacks.math.columbia.edu/tag/0GLZ}{0GLZ}].

For an affine limit \(U\), formula (2.6) follows from the fact that sections of an étale sheaf commute with affine inverse limits [Stacks, \href{https://stacks.math.columbia.edu/tag/099S}{099S}].
\end{proof}

% text: prose-02-31
\phantomsection\label{bp92}%
For \(X=\operatorname{Spec}\mathbf F_q\) and \(U=\operatorname{Spec}\overline{\mathbf F}_q\), formula (2.6) says that \((\nu_X^*F)(U)\) is the colimit of the \(F(\operatorname{Spec}\mathbf F_{q^m})\), that is, the underlying set of the Galois module corresponding to \(F\). Formula (2.6) also shows at once that every sheaf \(\nu_X^*F\) is classical. The next theorem says that there are no other classical sheaves.

% object: theorem-classical-sheaves-essential-image (blueprint/chap-The-pro-etale-site/sec-Classical-sheaves/theorem-classical-sheaves-essential-image/_section.md)
\phantomsection\label{bp93}%
\begin{bpstatementit}{Theorem~2.32}{ (Classical sheaves as the essential image)}
The functor \(\nu_X^*\) is fully faithful. Its essential image is exactly the classical sheaves of Definition \hyperref[bp89]{2.30}. For modules it is an exact equivalence between \(\operatorname{Mod}(X_{\acute{e}t},\Lambda)\) and the full subcategory of classical \(\Lambda_X\)-module sheaves on \(X_{\mathrm{pro}\acute{e}t}\).
\end{bpstatementit}

% proof: proof/theorem-classical-sheaves-essential-image
\begin{proof}[{Proof}]\phantomsection\label{bp94}%
A left adjoint is fully faithful exactly when its unit is an isomorphism, and the essential image is then the set of objects on which the counit is an isomorphism. We check both on pro-étale affines, where (2.6) computes \(\nu_X^*\).

Evaluating (2.5) on an étale object \(V\) gives \(F(V)\), so \(F\xrightarrow{\sim}\nu_{X,*}\nu_X^*F\), which proves full faithfulness. Moreover, every \(\nu_X^*F\) is classical by (2.6). Conversely, let \(G\) be classical. On a pro-étale affine the counit \(\nu_X^*\nu_{X,*}G\to G\) is

\[
(\nu_X^*\nu_{X,*}G)(U)
=\varinjlim_iG(U_i)\longrightarrow G(U),
\]

which is an isomorphism by (2.4). Since pro-étale affines form a covering basis (Proposition \hyperref[bp44]{2.15}), the counit is an isomorphism of sheaves.
\end{proof}

% object: example-coefficient-sheaf-outside-classical-part (blueprint/chap-The-pro-etale-site/sec-Classical-sheaves/example-coefficient-sheaf-outside-classical-part/_section.md)
\phantomsection\label{bp95}%
\begin{bpstatement}{Example~2.33}{ (A coefficient sheaf outside the classical part)}
Let \(X\) be a geometric point and let \(S=\mathbf Z_\ell\) be regarded as a profinite object of \(X_{\mathrm{pro}\acute{e}t}\). The completed coefficient sheaf \(\mathbf Z_{\ell,X}\) has sections \(\operatorname{Cont}(S,\mathbf Z_\ell)\) on \(S\), among them the identity function. The pullback \(\nu_X^*\) from \(X_{\acute{e}t}\) of the constant sheaf with value the \emph{discrete} abelian group \(\mathbf Z_\ell\) has only locally constant functions as sections, and these have finite image on the compact space \(S\). Hence the identity function does not come from any finite stage of the profinite presentation of \(S\), and \(\mathbf Z_{\ell,X}\) is not classical. Its reductions \(\mathbf Z/\ell^n\), on the other hand, are classical. So étale arguments apply at each finite level, while the coefficient sheaf itself is defined once on the pro-étale site.
\end{bpstatement}


% object: sec-Comparison-with-the-etale-site (blueprint/chap-The-pro-etale-site/sec-Comparison-with-the-etale-site/_section.md), content hidden
\subsection*{\texorpdfstring{2.8\quad Comparison with the étale site}{2.8 Comparison with the étale site}}
\phantomsection\label{bp96}%
\addcontentsline{toc}{subsection}{\texorpdfstring{2.8\quad Comparison with the étale site}{2.8 Comparison with the étale site}}
% text: prose-02-32
\phantomsection\label{bp97}%
For a bounded-below complex \(K\) on \(X_{\acute{e}t}\), the pro-étale cohomology of \(\nu_X^*K\) is the étale cohomology of \(K\) (Theorem \hyperref[bp101]{2.35}). Bounds on cohomological dimension are needed only for unbounded complexes (Section \hyperref[bp133]{2.11}). With this theorem, statements of SGA 4 about the cohomology of torsion sheaves, such as \(\mathbf Z/\ell^n\) on a variety, can be quoted on the pro-étale site.

The theorem rests on two facts. A pro-étale cover of an affine can be refined to a limit of étale covers, and sections of a classical sheaf on a limit are the colimit of sections. It follows that the Čech complex of a pro-étale cover is a filtered colimit of étale Čech complexes. Because of this, an injective étale sheaf stays acyclic after pullback (Lemma \hyperref[bp98]{2.34}), and resolving by injectives gives the theorem. For non-classical coefficients the comparison fails, as the computation of \(H^1(\operatorname{Spec}\mathbf F_q,\mathbf Z_\ell)\) in Section \bpout{1.2} shows.

It follows that the cohomology of a pro-étale affine is the colimit of the cohomology of its étale approximations (Corollary \hyperref[bp104]{2.36}). Moreover, classical complexes in \(D^+\) are exactly the pullbacks of étale ones (Corollary \hyperref[bp107]{2.37}), and this is what we need for proper base change on the pro-étale site.

% object: lemma-acyclicity-pulled-back-injectives (blueprint/chap-The-pro-etale-site/sec-Comparison-with-the-etale-site/lemma-acyclicity-pulled-back-injectives/_section.md)
\phantomsection\label{bp98}%
\begin{bpstatementit}{Lemma~2.34}{ (Acyclicity of pulled-back injectives)}
If \(I\) is an injective sheaf of \(\Lambda\)-modules on \(X_{\acute{e}t}\), then \(R^q\nu_{X,*}\nu_X^*I=0\) for \(q>0\).
\end{bpstatementit}

% proof: proof/lemma-acyclicity-pulled-back-injectives
\begin{proof}[{Proof}]\phantomsection\label{bp99}%
The sheaf \(R^q\nu_{X,*}\nu_X^*I\) is associated to \(V\mapsto H^q(V_{\mathrm{pro}\acute{e}t},\nu_X^*I)\), so it suffices to show that positive cohomology vanishes on affine étale objects \(V\). The restriction of \(I\) to \(V_{\acute{e}t}\) is injective, since restriction has an exact left adjoint, namely extension by zero. Hence we may assume \(X\) affine. By Theorem \hyperref[bp40]{2.14}, every pro-étale covering of an affine object is refined by a finite covering by affines that are ind-étale over it. By the Čech acyclicity criterion for a basis, it therefore suffices to show that the Čech complexes of \(\nu_X^*I\) for such coverings have no positive cohomology.

As in [Stacks, \href{https://stacks.math.columbia.edu/tag/099U}{099U}], we claim that the Čech complex of such a covering is a filtered colimit of Čech complexes of étale coverings. Write each covering map as an inverse limit of affine étale maps over the object being covered. By (2.6), the Čech complex of the covering is the filtered colimit of the Čech complexes of these étale approximations. Each finitely presented approximation descends to a common étale stage over the base, becomes a covering after increasing that stage, and its nerve descends by taking fibre products. Each augmented Čech complex at a finite stage is obtained by applying \(\Hom(-,I)\) to the exact augmented complex of free sheaves of \(\Lambda\)-modules on the nerve, so it has no positive cohomology because \(I\) is injective. Filtered colimits are exact, so the same holds for the pro-étale covering, and the criterion implies the required vanishing.
\end{proof}

% text: prose-02-33
\phantomsection\label{bp100}%
Since injective sheaves compute derived functors, the lemma suffices to compare derived direct images, and we obtain the comparison of cohomology by resolving a bounded-below complex by injectives.

% object: theorem-bounded-below-cohomology-comparison (blueprint/chap-The-pro-etale-site/sec-Comparison-with-the-etale-site/theorem-bounded-below-cohomology-comparison/_section.md)
\phantomsection\label{bp101}%
\begin{bpstatementit}{Theorem~2.35}{ (Bounded-below cohomology comparison)}
For \(K\in D^+(X_{\acute{e}t},\Lambda)\) the unit of the adjunction between \(\nu_X^*:D^+(X_{\acute{e}t},\Lambda)\to D^+(X_{\mathrm{pro}\acute{e}t},\Lambda)\) and \(R\nu_{X,*}\) is an isomorphism

\[
K\xrightarrow{\ \sim\ }R\nu_{X,*}\nu_X^*K.
\tag{2.7}
\]

Consequently \(\nu_X^*:D^+(X_{\acute{e}t},\Lambda)\to D^+(X_{\mathrm{pro}\acute{e}t},\Lambda)\) is fully faithful, and there is a canonical isomorphism

\[
R\Gamma(X_{\mathrm{pro}\acute{e}t},\nu_X^*K)
\simeq R\Gamma(X_{\acute{e}t},K).
\]
\end{bpstatementit}

% proof: proof/theorem-bounded-below-cohomology-comparison
\begin{proof}[{Proof}]\phantomsection\label{bp102}%
We compute \(R\nu_{X,*}\nu_X^*K\) with an injective resolution of \(K\) on the étale site. By the lemma its pullback is still acyclic for \(\nu_{X,*}\), and by full faithfulness on sheaves the result is the resolution itself.

Choose a bounded-below injective resolution \(K\to I^\bullet\). Since \(\nu_X^*\) is exact, \(\nu_X^*I^\bullet\) represents \(\nu_X^*K\), and by Lemma \hyperref[bp98]{2.34} its terms are \(\nu_{X,*}\)-acyclic, so \(R\nu_{X,*}\nu_X^*K\) is computed by the complex \(\nu_{X,*}\nu_X^*I^\bullet\). By Theorem \hyperref[bp93]{2.32} the unit \(I^q\to\nu_{X,*}\nu_X^*I^q\) is an isomorphism for every \(q\), so this complex is \(I^\bullet\), which proves (2.7). Since \(\Gamma(X_{\mathrm{pro}\acute{e}t},-)=\Gamma(X_{\acute{e}t},-)\circ\nu_{X,*}\), the statement on global sections follows by composing the derived functors; see [Stacks, \href{https://stacks.math.columbia.edu/tag/099V}{099V}, \href{https://stacks.math.columbia.edu/tag/099W}{099W}].
\end{proof}

% text: prose-02-34
\phantomsection\label{bp103}%
Applied to a pro-étale affine, the theorem computes its cohomology from the étale schemes of which it is the limit.

% object: corollary-cohomology-pro-etale-affine (blueprint/chap-The-pro-etale-site/sec-Comparison-with-the-etale-site/corollary-cohomology-pro-etale-affine/_section.md)
\phantomsection\label{bp104}%
\begin{bpstatementit}{Corollary~2.36}{ (Cohomology of a pro-étale affine)}
Let \(K\in D^+(X_{\acute{e}t},\Lambda)\), and let \(U=\varprojlim_iU_i\) be a pro-étale affine with \(U_i\) affine and étale over \(X\). Then there is an isomorphism

\[
R\Gamma(U_{\mathrm{pro}\acute{e}t},(\nu_X^*K)|_U)
\simeq \varinjlim_i R\Gamma((U_i)_{\acute{e}t},K|_{U_i}).
\tag{2.8}
\]
\end{bpstatementit}

% proof: proof/corollary-cohomology-pro-etale-affine
\begin{proof}[{Proof}]\phantomsection\label{bp105}%
By the theorem applied to \(U\), the left side is \(R\Gamma(U_{\acute{e}t},K|_U)\), so the assertion is about étale cohomology alone. For a sheaf, continuity of étale cohomology under affine limits [Stacks, \href{https://stacks.math.columbia.edu/tag/03Q6}{03Q6}] identifies this with the colimit in each degree. For a bounded-below complex, only finitely many terms of the hypercohomology spectral sequence contribute to each total degree, and filtered colimits are exact.
\end{proof}

% text: prose-02-35
\phantomsection\label{bp106}%
To apply statements proved on the étale site, such as proper base change, on \(X_{\mathrm{pro}\acute{e}t}\), we also need the essential image of \(\nu_X^*\) on \(D^+\). It consists of the classical complexes. Write \(D^+_{\mathrm{class}}(X_{\mathrm{pro}\acute{e}t},\Lambda)\) for the full subcategory of \(D^+(X_{\mathrm{pro}\acute{e}t},\Lambda)\) of classical complexes.

% object: corollary-bounded-below-classical-complexes (blueprint/chap-The-pro-etale-site/sec-Comparison-with-the-etale-site/corollary-bounded-below-classical-complexes/_section.md)
\phantomsection\label{bp107}%
\begin{bpstatementit}{Corollary~2.37}{ (Bounded-below classical complexes)}
Classical module sheaves form a weak Serre subcategory, closed under kernels, cokernels and extensions. There is an equivalence

\[
D^+_{\mathrm{class}}(X_{\mathrm{pro}\acute{e}t},\Lambda)
\simeq D^+(X_{\acute{e}t},\Lambda),
\]

with inverse functors \(R\nu_{X,*}\) and \(\nu_X^*\).
\end{bpstatementit}

% proof: proof/corollary-bounded-below-classical-complexes
\begin{proof}[{Proof}]\phantomsection\label{bp108}%
We need the first assertion so that “classical complex” defines a triangulated subcategory, and for extensions we use the vanishing of \(R^1\nu_{X,*}\) from the theorem. The equivalence then follows by comparing cohomology sheaves.

For kernels and cokernels the first assertion follows from exactness and full faithfulness of \(\nu_X^*\). Indeed, a map between classical sheaves is \(\nu_X^*\) of a map of étale sheaves, and \(\nu_X^*\) carries its kernel and cokernel to the kernel and cokernel. Given an extension \(0\to\nu_X^*F\to G\to\nu_X^*F'\to0\), we apply \(\nu_{X,*}\). Since \(R^1\nu_{X,*}\nu_X^*F=0\) by Theorem \hyperref[bp101]{2.35}, the result \(0\to F\to\nu_{X,*}G\to F'\to0\) is exact, and the counit maps its pullback to the original extension. This map is the identity on the outer terms, so by the five lemma \(G\simeq\nu_X^*\nu_{X,*}G\) is classical.

The functor \(\nu_X^*\) lands in classical complexes and is fully faithful on \(D^+\) by Theorem \hyperref[bp101]{2.35}, so the equivalence amounts to the counit \(\nu_X^*R\nu_{X,*}K\to K\) being an isomorphism for \(K\) in \(D^+_{\mathrm{class}}\). By (2.7), every classical sheaf \(\nu_X^*F\) is \(\nu_{X,*}\)-acyclic, so the hypercohomology spectral sequence for \(R\nu_{X,*}K\) degenerates and \(H^q(R\nu_{X,*}K)=\nu_{X,*}H^q(K)\). On \(H^q\) the counit is therefore the counit \(\nu_X^*\nu_{X,*}H^q(K)\to H^q(K)\) for a classical sheaf, which is an isomorphism by Theorem \hyperref[bp93]{2.32}; see [Stacks, \href{https://stacks.math.columbia.edu/tag/09B1}{09B1}].
\end{proof}


% object: sec-Direct-images (blueprint/chap-The-pro-etale-site/sec-Direct-images/_section.md), content hidden
\subsection*{\texorpdfstring{2.9\quad Direct images}{2.9 Direct images}}
\phantomsection\label{bp109}%
\addcontentsline{toc}{subsection}{\texorpdfstring{2.9\quad Direct images}{2.9 Direct images}}
% text: prose-02-36
\phantomsection\label{bp110}%
Proper base change is proved on the étale site (Chapter \bpout{4}), while \(Rf_!\) is defined by direct images on the pro-étale site. To transfer proper base change, and to compute these direct images, we compare direct images on the two sites.

Let \(g:X\to Y\) be a morphism of schemes. Base change of objects of the sites defines a commutative square of topoi

\[
\begin{CD}
X_{\mathrm{pro}\acute{e}t}@>{\nu_X}>>X_{\acute{e}t}\\
@V{g_{\mathrm{pro}\acute{e}t}}VV @VV{g_{\acute{e}t}}V\\
Y_{\mathrm{pro}\acute{e}t}@>{\nu_Y}>>Y_{\acute{e}t}.
\end{CD}
\]

It induces an isomorphism \(g_{\mathrm{pro}\acute{e}t}^*\nu_Y^*\simeq\nu_X^*g_{\acute{e}t}^*\) of functors \(\operatorname{Mod}(Y_{\acute{e}t},\Lambda)\to\operatorname{Mod}(X_{\mathrm{pro}\acute{e}t},\Lambda)\), compatible with composition. Since all four functors are exact, it is also an isomorphism of functors \(D(Y_{\acute{e}t},\Lambda)\to D(X_{\mathrm{pro}\acute{e}t},\Lambda)\). For direct images the comparison needs a hypothesis on \(g\). The comparison map \(\nu_Y^*Rg_{\acute{e}t,*}\to Rg_{\mathrm{pro}\acute{e}t,*}\nu_X^*\) of functors \(D^+(X_{\acute{e}t},\Lambda)\to D^+(Y_{\mathrm{pro}\acute{e}t},\Lambda)\) is defined as the adjoint of the composite of \(g_{\mathrm{pro}\acute{e}t}^*\nu_Y^*Rg_{\acute{e}t,*}\simeq\nu_X^*g_{\acute{e}t}^*Rg_{\acute{e}t,*}\) with \(\nu_X^*\) applied to the counit of the adjunction \(g_{\acute{e}t}^*\dashv Rg_{\acute{e}t,*}\).

Evaluated on a pro-étale affine \(V=\varprojlim_iV_i\) over \(Y\), the comparison map is the map from \(\varinjlim_iH^q(X\times_YV_i,F)\) to \(H^q(X\times_YV,F)\). So the theorem below is the continuity of étale cohomology under affine limits, and like it, it needs \(g\) to be quasi-compact. Without this the map fails to be surjective already for \(q=0\). For example, let \(Y\) be affine and connected, \(X=\coprod_{n\in\mathbf N}Y\) and \(F\) the constant sheaf \(\mathbf Z/2\). Then \(g_*F\) is the product of countably many copies of \(\mathbf Z/2\) on \(Y\). On \(V=Y\times S\) with \(S\) an infinite profinite set, the pro-étale side has as sections all sequences of locally constant functions \(S\to\mathbf Z/2\), while the classical side only has the sequences that factor through a common finite quotient of \(S\).

% object: theorem-comparison-direct-images (blueprint/chap-The-pro-etale-site/sec-Direct-images/theorem-comparison-direct-images/_section.md)
\phantomsection\label{bp111}%
\begin{bpstatementit}{Theorem~2.38}{ (Comparison of direct images)}
If \(g\) is quasi-compact and quasi-separated, the canonical natural transformation

\[
\nu_Y^*Rg_{\acute{e}t,*}K\longrightarrow
Rg_{\mathrm{pro}\acute{e}t,*}\nu_X^*K
\tag{2.9}
\]

is an isomorphism for \(K\in D^+(X_{\acute{e}t},\Lambda)\).
\end{bpstatementit}

% proof: proof/theorem-comparison-direct-images
\begin{proof}[{Proof}]\phantomsection\label{bp112}%
We work over an affine open of \(Y\) and evaluate on the covering basis \(V=\varprojlim_iV_i\), with \(V_i\) affine and étale over \(Y\). By (2.6) and Theorem \hyperref[bp101]{2.35}, for a sheaf \(F\) the map in degree \(q\) is the continuity map

\[
\varinjlim_i H^q_{\acute{e}t}(X\times_YV_i,F)
\longrightarrow
H^q_{\acute{e}t}(X\times_YV,F).
\]

Since \(g\) is quasi-compact and quasi-separated, so are its base changes, and the map is an isomorphism by continuity of étale cohomology under affine limits [Stacks, \href{https://stacks.math.columbia.edu/tag/03Q6}{03Q6}]. Using bounded-below resolutions, or the convergent hypercohomology spectral sequences, we pass from sheaves to \(K\); see [Stacks, \href{https://stacks.math.columbia.edu/tag/09A6}{09A6}].
\end{proof}

% text: prose-02-37
\phantomsection\label{bp113}%
Since every bounded-below classical complex is of the form \(\nu_X^*K\) (Corollary \hyperref[bp107]{2.37}), the theorem shows that, for \(g\) quasi-compact and quasi-separated, \(Rg_{\mathrm{pro}\acute{e}t,*}\) carries bounded-below classical complexes to classical complexes.

In later arguments we transport maps built from units, counits and composition from the étale to the pro-étale site, and for this we record how the comparison map behaves under these operations. Write \(\beta_g:g_{\mathrm{pro}\acute{e}t}^*\nu_Y^*\simeq\nu_X^*g_{\acute{e}t}^*\) for the isomorphism above and \(c_g\) for the map (2.9).

% object: lemma-comparison-direct-images-compatibilities (blueprint/chap-The-pro-etale-site/sec-Direct-images/lemma-comparison-direct-images-compatibilities/_section.md)
\phantomsection\label{bp114}%
\begin{bpstatementit}{Lemma~2.39}{ (Compatibilities of the comparison of direct images)}
Let \(g:X\to Y\) be a morphism of schemes.

\begin{enumerate}
\item For \(L\in D^+(Y_{\acute{e}t},\Lambda)\), the composite \(\nu_Y^*L\to\nu_Y^*Rg_{\acute{e}t,*}g_{\acute{e}t}^*L\xrightarrow{c_g}
  Rg_{\mathrm{pro}\acute{e}t,*}\nu_X^*g_{\acute{e}t}^*L\) of \(\nu_Y^*\) of the unit with \(c_g\) equals the unit \(\nu_Y^*L\to Rg_{\mathrm{pro}\acute{e}t,*}g_{\mathrm{pro}\acute{e}t}^*\nu_Y^*L\) followed by \(Rg_{\mathrm{pro}\acute{e}t,*}(\beta_g)\).
\item For \(K\in D^+(X_{\acute{e}t},\Lambda)\), the composite \(g_{\mathrm{pro}\acute{e}t}^*\nu_Y^*Rg_{\acute{e}t,*}K\xrightarrow{g^*c_g}
  g_{\mathrm{pro}\acute{e}t}^*Rg_{\mathrm{pro}\acute{e}t,*}\nu_X^*K\to\nu_X^*K\) with the counit equals \(\nu_X^*\) of the counit \(g_{\acute{e}t}^*Rg_{\acute{e}t,*}K\to K\) precomposed with \(\beta_g\).
\item For a second morphism \(h:Y\to Z\), under the identifications \(R(h\circ g)_*\simeq Rh_*Rg_*\) on both sites, \(c_{h\circ g}\) is the composite \[
   \nu_Z^*Rh_{\acute{e}t,*}Rg_{\acute{e}t,*}K\xrightarrow{c_h}
   Rh_{\mathrm{pro}\acute{e}t,*}\nu_Y^*Rg_{\acute{e}t,*}K\xrightarrow{Rh_{\mathrm{pro}\acute{e}t,*}(c_g)}
   Rh_{\mathrm{pro}\acute{e}t,*}Rg_{\mathrm{pro}\acute{e}t,*}\nu_X^*K .
   \]
\end{enumerate}
\end{bpstatementit}

% proof: proof/lemma-comparison-direct-images-compatibilities
\begin{proof}[{Proof}]\phantomsection\label{bp115}%
Part (2) is the definition of \(c_g\) as an adjoint. For (1), by (2) and the naturality of \(\beta_g\), the adjoint of the composite in (1) is \(\beta_g\) followed by \(\nu_X^*\) of the composite \(g^*_{\acute{e}t}L\to g^*_{\acute{e}t}Rg_{\acute{e}t,*}g^*_{\acute{e}t}L\to g^*_{\acute{e}t}L\), which is the identity by the triangle identity. Hence the adjoint is \(\beta_g\), whose adjoint is the map on the right of (1). For (3), the isomorphisms \(\beta\) are compatible with composition, \(\beta_{h\circ g}\) being \(\beta_g\) composed with \(g_{\mathrm{pro}\acute{e}t}^*(\beta_h)\), and the adjoint of a composite of maps defined as in (2) is the composite of their adjoints.
\end{proof}

% text: prose-02-38
\phantomsection\label{bp116}%
In the comparison for constructible sheaves below we argue by induction on strata, using the functors of localization triangles. The next proposition identifies these functors on the two sites.

% object: proposition-comparison-immersions (blueprint/chap-The-pro-etale-site/sec-Direct-images/proposition-comparison-immersions/_section.md)
\phantomsection\label{bp117}%
\begin{bpstatementit}{Proposition~2.40}{ (Comparison for immersions)}
\begin{enumerate}
\item For an open immersion \(j:U\hookrightarrow X\), there is a canonical isomorphism \[
   \nu_X^*j_{\acute{e}t,!}\simeq j_{\mathrm{pro}\acute{e}t,!}\nu_U^*\quad\text{of functors }\operatorname{Mod}(U_{\acute{e}t},\Lambda)\to\operatorname{Mod}(X_{\mathrm{pro}\acute{e}t},\Lambda),
   \tag{2.10}
   \] and hence of functors \(D(U_{\acute{e}t},\Lambda)\to D(X_{\mathrm{pro}\acute{e}t},\Lambda)\).
\item For a closed immersion \(i:Z\hookrightarrow X\), the map (2.9) is an isomorphism \(\nu_X^*Ri_{\acute{e}t,*}\simeq Ri_{\mathrm{pro}\acute{e}t,*}\nu_Z^*\) of functors \(D^+(Z_{\acute{e}t},\Lambda)\to D^+(X_{\mathrm{pro}\acute{e}t},\Lambda)\), and \(i_{\mathrm{pro}\acute{e}t}^*\nu_X^*\simeq\nu_Z^*i_{\acute{e}t}^*\). For \(K\in D^+(X_{\acute{e}t},\Lambda)\), these isomorphisms and (2.10) carry \(\nu_X^*\) of the localization triangle \(j_{\acute{e}t,!}j_{\acute{e}t}^*K\to K\to Ri_{\acute{e}t,*}i_{\acute{e}t}^*K\to\) to a distinguished triangle \(j_{\mathrm{pro}\acute{e}t,!}j_{\mathrm{pro}\acute{e}t}^*\nu_X^*K\to\nu_X^*K\to Ri_{\mathrm{pro}\acute{e}t,*}i_{\mathrm{pro}\acute{e}t}^*\nu_X^*K\to\) whose first two maps are the counit and the unit of the adjunctions \(j_!\dashv j^*\) and \(i^*\dashv Ri_*\) on \(X_{\mathrm{pro}\acute{e}t}\).
\item Let \(f:X\to Y\) be a separated morphism of finite type with a factorization \(f=\bar f\circ j\), where \(j\) is an open immersion and \(\bar f\) is proper. Then there is a canonical isomorphism \(\nu_Y^*R\bar f_{\acute{e}t,*}j_{\acute{e}t,!}\simeq R\bar f_{\mathrm{pro}\acute{e}t,*}j_{\mathrm{pro}\acute{e}t,!}\nu_X^*\) of functors \(D^+(X_{\acute{e}t},\Lambda)\to D^+(Y_{\mathrm{pro}\acute{e}t},\Lambda)\).
\end{enumerate}
\end{bpstatementit}

% proof: proof/proposition-comparison-immersions
\begin{proof}[{Proof}]\phantomsection\label{bp118}%
For (1), both functors are left adjoint to restriction followed by \(\nu_{U,*}\), and this also identifies the adjunction maps [Stacks, \href{https://stacks.math.columbia.edu/tag/09AD}{09AD}]. As the functors are exact, the identity passes to derived categories. In (2), the first identity is Theorem \hyperref[bp111]{2.38}, and the second is the comparison of inverse images. Since \(\nu_X^*\) is exact, it carries the localization triangle to a distinguished triangle. Its first map is carried to the counit on \(X_{\mathrm{pro}\acute{e}t}\) because the isomorphism of (1) identifies the adjunction maps, and its second map to the unit by Lemma lemma-comparison-direct-images-compatibilities, applied to \(g=i\) and \(L=K\). Part (3) follows from (1) and Theorem \hyperref[bp111]{2.38} for the proper, hence quasi-compact and quasi-separated, map \(\bar f\).
\end{proof}


% object: sec-Constructible-sheaves (blueprint/chap-The-pro-etale-site/sec-Constructible-sheaves/_section.md), content hidden
\subsection*{\texorpdfstring{2.10\quad Constructible sheaves}{2.10 Constructible sheaves}}
\phantomsection\label{bp119}%
\addcontentsline{toc}{subsection}{\texorpdfstring{2.10\quad Constructible sheaves}{2.10 Constructible sheaves}}
% text: prose-02-39
\phantomsection\label{bp120}%
The six operations act on constructible complexes, and their finiteness theorems are proved on the étale site. To use these theorems we show that constructible pro-étale sheaves are classical, so that they correspond to constructible étale sheaves, and that the comparison preserves finite Tor amplitude.

A sheaf of modules on \(X_{\mathrm{pro}\acute{e}t}\) is locally constant if it is trivialized by a pro-étale covering, and it is of finite presentation if its constant models are finitely presented \(\Lambda\)-modules. Building on this, for Noetherian \(\Lambda\), a sheaf is constructible if, locally on \(X\), there is a finite partition into constructible locally closed strata such that it is locally constant of finite type on each stratum. We write \(\Cons(X,\Lambda)\) for the category of constructible sheaves and \(\Dbc(X,\Lambda)\) for the bounded complexes with constructible cohomology sheaves. A complex \(K\) has \emph{Tor amplitude in \([a,b]\)} if \(K\otimes^{\mathbf L}_{\Lambda_X}N\) has cohomology only in \([a,b]\) for every sheaf \(N\) of \(\Lambda_X\)-modules, and \(D^b_{\mathrm{ctf}}(X,\Lambda)\subset\Dbc(X,\Lambda)\) consists of the complexes of finite Tor amplitude. For a discrete ring these notions agree with Definition \bpout{3.1}, which is stated for all coefficient rings in Section \bpout{3.2} below. These conditions refer only to pro-étale covers, so the main step is to descend torsors from pro-étale to étale covers.

A priori there could be more pro-étale local systems than étale ones, since the covers are larger. For finite groups, however, there are not. Indeed, a torsor under a finite group is described by a cocycle on one pro-étale affine cover. That cover is a limit of étale covers and the group is classical, so the cocycle is defined at a finite stage. For \(\mathbf Z_\ell\), by contrast, there are more local systems, and this is the reason for using the pro-étale site. For example, over \(\operatorname{Spec}\mathbf F_q\), a rank-one \(\mathbf Z_\ell\)-local system on which Frobenius acts by a unit \(u\in\mathbf Z_\ell^\times\) of infinite order is trivialized by a pro-étale cover but by no étale one. Our comparison therefore concerns only finite coefficients, and we study \(\ell\)-adic local systems through their reductions.

What we use later is the case of a finite ring \(A\). For such \(A\), Theorem \hyperref[bp127]{2.43} and Corollary \hyperref[bp130]{2.44} identify the category \(D^b_{\mathrm{ctf}}(X,A)\) on the pro-étale site with the classical one.

% object: lemma-torsors-under-classical-groups (blueprint/chap-The-pro-etale-site/sec-Constructible-sheaves/lemma-torsors-under-classical-groups/_section.md)
\phantomsection\label{bp121}%
\begin{bpstatementit}{Lemma~2.41}{ (Torsors under classical groups)}
Every pro-étale torsor under a classical sheaf of groups is classical, and hence has a trivializing étale covering.
\end{bpstatementit}

% proof: proof/lemma-torsors-under-classical-groups
\begin{proof}[{Proof}]\phantomsection\label{bp122}%
Write the group as \(\nu_X^*G_0\). We may work on an affine open of \(X\). Using Theorem \hyperref[bp40]{2.14} and finite affine refinements, we refine a trivializing covering to a single faithfully flat ind-étale affine \(U\to X\), and we write \(U=\varprojlim_iU_i\) with \(U_i\) affine étale over \(X\). A trivialization over \(U\) determines a transition element in \(\nu_X^*G_0(U\times_XU)\). By continuity for affine limits it descends to \(G_0(U_i\times_XU_i)\) for some \(i\), and the cocycle identity on the triple product holds after enlarging \(i\). Since \(U\to X\) is surjective, each \(U_i\to X\) is an étale covering. Gluing the trivial torsor along this cocycle, we get an étale torsor whose pullback is isomorphic to the given torsor. Hence the given torsor is classical by Theorem \hyperref[bp93]{2.32}, and it is trivialized by the étale covering \(U_i\to X\). See [Stacks, \href{https://stacks.math.columbia.edu/tag/099X}{099X}].
\end{proof}

% text: prose-02-40
\phantomsection\label{bp123}%
The hypothesis that the group is classical is essential. For instance, the group \(\mathbf Z_\ell^\times\) is not classical, and the rank-one \(\mathbf Z_\ell\)-local system on \(\operatorname{Spec}\mathbf F_q\) described above is a \(\mathbf Z_\ell^\times\)-torsor with no étale trivialization. On the other hand, for locally constant sheaves with finitely presented constant models over a discrete ring, the relevant group is the automorphism group of the model, which is discrete, and the lemma applies.

% object: proposition-locally-constant-sheaves-finite-presentation (blueprint/chap-The-pro-etale-site/sec-Constructible-sheaves/proposition-locally-constant-sheaves-finite-presentation/_section.md)
\phantomsection\label{bp124}%
\begin{bpstatementit}{Proposition~2.42}{ (Locally constant sheaves of finite presentation)}
Every locally constant pro-étale \(\Lambda_X\)-module sheaf of finite presentation is classical. In particular, it has a trivializing étale covering, and \(\nu_X^*\) identifies the étale local systems of finite presentation with these pro-étale ones.
\end{bpstatementit}

% proof: proof/proposition-locally-constant-sheaves-finite-presentation
\begin{proof}[{Proof}]\phantomsection\label{bp125}%
A locally constant sheaf with model \(M\) is the same as a torsor under the automorphism sheaf of the constant sheaf \(\underline M\), so the proposition follows from the lemma once we know that this automorphism sheaf is classical.

Decompose \(X\) into the open pieces on which the constant model has a fixed isomorphism type, and fix a finitely presented module \(M\). A presentation \(\Lambda^m\to\Lambda^n\to M\to0\) identifies the internal endomorphism sheaf of \(\underline M\) with the kernel of \(\underline M^{\,n}\to\underline M^{\,m}\). Since the constant sheaf functor is exact, this kernel is the constant sheaf on \(\operatorname{End}_\Lambda(M)\), and so the automorphism sheaf of \(\underline M\) is the constant discrete group \(\operatorname{Aut}_\Lambda(M)\). Hence \(\mathit{Isom}(\underline M,F)\) is a torsor under this classical group, and by Lemma \hyperref[bp121]{2.41} it has sections étale locally. Since \(F\) is obtained from this torsor by twisting \(\underline M\), it is the pullback of the corresponding twist of the constant étale sheaf, and so it is classical. This is the descent argument of [Stacks, \href{https://stacks.math.columbia.edu/tag/099Y}{099Y}].
\end{proof}

% text: prose-02-41
\phantomsection\label{bp126}%
To pass from a single stratum to arbitrary constructible sheaves, note that for a finite ring \(A\) the proposition applies to every finite local system of \(A\)-modules. A constructible sheaf is built from such local systems on the strata of a stratification, by extension by zero and extensions. Both operations preserve classical sheaves. For extension by zero this follows from the comparisons of Section \hyperref[bp109]{2.9}. For extensions it is Corollary \hyperref[bp107]{2.37}, which uses that \(R^1\nu_{X,*}\) vanishes on classical sheaves.

% object: theorem-comparison-constructible-sheaves (blueprint/chap-The-pro-etale-site/sec-Constructible-sheaves/theorem-comparison-constructible-sheaves/_section.md)
\phantomsection\label{bp127}%
\begin{bpstatementit}{Theorem~2.43}{ (Comparison for constructible sheaves)}
Let \(X\) be quasi-compact and quasi-separated. For a discrete Noetherian \(\Lambda\), every constructible pro-étale \(\Lambda_X\)-module sheaf is classical. The functor \(\nu_X^*\) gives an equivalence from constructible étale sheaves of \(\Lambda\)-modules to \(\Cons(X,\Lambda)\), and an equivalence \(D^b_c(X_{\acute{e}t},\Lambda)\to\Dbc(X,\Lambda)\) of bounded constructible derived categories.
\end{bpstatementit}

% proof: proof/theorem-comparison-constructible-sheaves
\begin{proof}[{Proof}]\phantomsection\label{bp128}%
The preceding proposition treats the case of a single stratum. For a finite stratification we argue locally, by induction on the number of strata. Once the sheaves are known to be classical, the equivalences follow from full faithfulness of \(\nu_X^*\) and its version for \(D^+\).

Given a point, we shrink to an affine open on which the stratum containing it is closed, with inclusion \(i\) and open complement \(j\). The exact sequence

\[
0\longrightarrow j_!j^*F\longrightarrow F
\longrightarrow i_*i^*F\longrightarrow0
\]

reduces the assertion to fewer strata, because on pro-étale sheaves of modules the functors \(j_!\) and \(i_*\) preserve classical sheaves, by (2.10) and Theorem \hyperref[bp111]{2.38}, and classical sheaves are closed under extensions (Corollary \hyperref[bp107]{2.37}). Since \(\nu_X^*\) is fully faithful, sheaves that are classical on the members of a Zariski covering are classical, so the local statement suffices. The derived statement follows from Corollary \hyperref[bp107]{2.37} and bounded truncation; see [Stacks, \href{https://stacks.math.columbia.edu/tag/09AK}{09AK}, \href{https://stacks.math.columbia.edu/tag/09AL}{09AL}].
\end{proof}

% text: prose-02-42
\phantomsection\label{bp129}%
The six operations act on complexes of finite Tor amplitude, so we also need to know that the equivalence just proved respects this condition. A priori the pro-étale condition is stronger, since it requires a Tor bound against all pro-étale modules, most of which are not classical.

% object: corollary-finite-tor-amplitude-under-comparison (blueprint/chap-The-pro-etale-site/sec-Constructible-sheaves/corollary-finite-tor-amplitude-under-comparison/_section.md)
\phantomsection\label{bp130}%
\begin{bpstatementit}{Corollary~2.44}{ (Finite Tor amplitude under comparison)}
Let \(X\) be quasi-compact and quasi-separated and let \(\Lambda\) be a discrete Noetherian ring. The equivalence \(\nu_X^*:D^b_c(X_{\acute{e}t},\Lambda)\to\Dbc(X,\Lambda)\) of Theorem \hyperref[bp127]{2.43} restricts to an equivalence

\[
D^b_{\mathrm{ctf}}(X_{\acute{e}t},\Lambda)\simeq D^b_{\mathrm{ctf}}(X,\Lambda).
\]
\end{bpstatementit}

% proof: proof/corollary-finite-tor-amplitude-under-comparison
\begin{proof}[{Proof}]\phantomsection\label{bp131}%
For one direction we represent an étale complex by a bounded flat complex and pull it back, and for the other we use the fact that \(\nu_X^*\) detects vanishing.

The functor \(\nu_X^*\) is exact (Proposition \hyperref[bp90]{2.31}), and since \(\nu_X^*\Lambda=\Lambda_X\) it is symmetric monoidal on sheaves of modules. The pullback of a flat sheaf of modules is flat [Stacks, \href{https://stacks.math.columbia.edu/tag/05VD}{05VD}]. An étale complex of Tor amplitude in \([a,b]\) is represented by a complex of flat sheaves in degrees \([a,b]\), so its pullback is a complex of flat pro-étale sheaves in the same degrees and has Tor amplitude in \([a,b]\) against every pro-étale module, not only against pulled-back ones. Conversely, let \(K\) be an étale complex. Its Tor amplitude can be read off from \(K\otimes^{\mathbf L}_\Lambda M\) for arbitrary étale sheaves of modules \(M\). The inverse image of this complex is \(\nu_X^*K\otimes^{\mathbf L}_{\Lambda_X}\nu_X^*M\), and since \(\nu_X^*\) is exact and fully faithful, it detects vanishing of cohomology. Hence a Tor bound on \(\nu_X^*K\) implies the same bound on \(K\).
\end{proof}

% text: prose-02-43
\phantomsection\label{bp132}%
In particular, the corollary applies to every finite commutative \(\ell\)-primary ring \(A\).


% object: sec-Unbounded-complexes (blueprint/chap-The-pro-etale-site/sec-Unbounded-complexes/_section.md), content hidden
\subsection*{\texorpdfstring{2.11\quad Unbounded complexes}{2.11 Unbounded complexes}}
\phantomsection\label{bp133}%
\addcontentsline{toc}{subsection}{\texorpdfstring{2.11\quad Unbounded complexes}{2.11 Unbounded complexes}}
% text: prose-02-44
\phantomsection\label{bp134}%
Proper base change needs the comparison with the étale site only for bounded-below complexes. Brown representability, however, takes place in the unbounded derived category, so we extend the comparison to unbounded complexes. The difficulty is that a morphism from a sheaf into an unbounded complex \(K\) involves in principle all the cohomology sheaves of \(K\) at once, and so does a cohomology sheaf of \(Rg_*K\). Both depend on only finitely many of them if Ext groups and cohomology of sheaves vanish above a fixed degree, uniformly in the sheaf. The smallest such degree is the cohomological dimension defined below.

Without such a bound, \(\Hom\) out of the constant sheaf need not commute with direct sums. For example, take \(\ell=2\) and \(S=\operatorname{Spec}\mathbf R\). Here étale cohomology is the cohomology of \(\operatorname{Gal}(\mathbf C/\mathbf R)=\mathbf Z/2\), and \(H^q(\mathbf Z/2,\mathbf F_2)=\mathbf F_2\) for every \(q\geq0\). Thus \(\operatorname{cd}_2(S)=\infty\), and \(\Hom(\mathbf F_2,-)\) does not commute with the direct sum \(\bigoplus_{n\geq0}\mathbf F_2[n]\) (see the opening of Chapter \bpout{5}).

Over a field of finite cohomological dimension, the étale cohomology of schemes of finite type is uniformly bounded. In that case every computation in a fixed degree can be carried out on a bounded-below truncation. We use the bounds in this way in Proposition \hyperref[bp144]{2.49} below and, later, in the proofs of compactness and of the preservation of direct sums by \(Rp_*\) (Lemmas \bpout{5.6} and \bpout{Lemma 5.17}).

% object: definition-cohomological-dimension (blueprint/chap-The-pro-etale-site/sec-Unbounded-complexes/definition-cohomological-dimension/_section.md)
\phantomsection\label{bp135}%
\begin{bpstatement}{Definition~2.45}{ (Cohomological dimension)}
For a qcqs scheme \(S\) and a qcqs morphism \(g:X\to Y\), the \emph{\(\ell\)-cohomological dimensions} are

\[
\begin{aligned}
\operatorname{cd}_\ell(S)&=\sup\{q: H^q(S_{\mathrm{pro}\acute{e}t},F)\neq0
\text{ for some classical $\ell$-power torsion sheaf $F$}\},\\
\operatorname{cd}_\ell(g)&=\sup\{q: R^qg_*F\neq0
\text{ for some classical $\ell$-power torsion sheaf $F$ on }X_{\mathrm{pro}\acute{e}t}\}.
\end{aligned}
\]

For a field \(k'\) we write \(\operatorname{cd}_\ell(k')=\operatorname{cd}_\ell(\operatorname{Spec}k')\).
\end{bpstatement}

% text: prose-02-45
\phantomsection\label{bp136}%
By Theorems \hyperref[bp101]{2.35} and \hyperref[bp111]{Theorem 2.38}, these are the usual \(\ell\)-cohomological dimensions computed on the étale sites. So the classical bounds apply. If \(k'\) is a field with \(\operatorname{cd}_\ell(k')<\infty\), then \(\operatorname{cd}_\ell(S)\) and \(\operatorname{cd}_\ell(g)\) are finite for all schemes \(S\) and morphisms \(g\) of finite type over \(k'\) [Stacks, \href{https://stacks.math.columbia.edu/tag/0F10}{0F10}]. For our field \(k\) the hypothesis holds by the following lemma.

% object: lemma-cohomological-dimension-base-field (blueprint/chap-The-pro-etale-site/sec-Unbounded-complexes/lemma-cohomological-dimension-base-field/_section.md)
\phantomsection\label{bp137}%
\begin{bpstatementit}{Lemma~2.46}{ (Cohomological dimension of the base field)}
If \(k'\) is an algebraically closed field, then \(\operatorname{cd}_\ell(k')=0\). If \(k'\) is a finite field, then \(\operatorname{cd}_\ell(k')\le1\).
\end{bpstatementit}

% proof: proof/lemma-cohomological-dimension-base-field
\begin{proof}[{Proof}]\phantomsection\label{bp138}%
Étale cohomology of \(\operatorname{Spec}k'\) is Galois cohomology [Stacks, \href{https://stacks.math.columbia.edu/tag/03QQ}{03QQ}]. For algebraically closed \(k'\) the Galois group is trivial. For finite \(k'\) it is \(\widehat{\mathbf Z}\), which has cohomological dimension one (see Serre, \emph{Galois cohomology}, II.§3).
\end{proof}

% text: prose-02-46
\phantomsection\label{bp139}%
Over \(\overline{\mathbf F}_q\), for example, a variety of dimension \(n\) has \(\ell\)-cohomological dimension at most \(2n\), and the curve \(\mathbf P^1\) attains this bound, since \(H^2(\mathbf P^1,\mathbf Z/\ell)\neq0\). Over \(\mathbf F_q\), however, one more degree is possible. Indeed, for \(\ell\neq p\) the Hochschild–Serre spectral sequence gives \(H^3(\mathbf P^1_{\mathbf F_q},\mu_\ell)\simeq
H^1(\operatorname{Spec}\mathbf F_q,H^2(\mathbf P^1_{\overline{\mathbf F}_q},\mu_\ell))
=H^1(\operatorname{Spec}\mathbf F_q,\mathbf Z/\ell)=\mathbf Z/\ell\).

To extend the comparison theorems to unbounded complexes with these bounds, we need two facts. First, the bounds hold uniformly over the affine objects on which we take sections (Lemma \hyperref[bp140]{2.47}). Second, on the pro-étale side a cohomology group of an unbounded complex in a fixed degree only sees a bounded-below truncation (Lemma \hyperref[bp142]{2.48}).

% object: lemma-uniform-bounds-pro-etale-affines (blueprint/chap-The-pro-etale-site/sec-Unbounded-complexes/lemma-uniform-bounds-pro-etale-affines/_section.md)
\phantomsection\label{bp140}%
\begin{bpstatementit}{Lemma~2.47}{ (Uniform bounds)}
Let \(X\) and \(Y\) be of finite type over \(k\), let \(A\) be a finite commutative \(\ell\)-primary ring, let \(g:X\to Y\) be of finite type, and put \(d=\dim(X)+\operatorname{cd}_\ell(k)\). There is an integer \(c\), depending only on \(X\), such that for every classical sheaf \(F\) of \(A\)-modules on \(X_{\mathrm{pro}\acute{e}t}\), every pro-étale affine \(U\) over \(X\) and every pro-étale affine \(V\) over \(Y\),

\[
H^q(U,F)=0\ \text{ for }q>d,
\qquad
H^q(X\times_YV,F)=0\ \text{ for }q>c.
\]

Moreover \(\operatorname{cd}_\ell(g)<\infty\).
\end{bpstatementit}

% proof: proof/lemma-uniform-bounds-pro-etale-affines
\begin{proof}[{Proof}]\phantomsection\label{bp141}%
We write both cohomologies as colimits of étale cohomology of schemes of finite type over \(k\) and bound these uniformly. Write \(U=\varprojlim_iU_i\) and \(V=\varprojlim_iV_i\) with \(U_i\) affine étale over \(X\) and \(V_i\) affine étale over \(Y\), and write \(F=\nu_X^*F_{\acute{e}t}\) (Theorem \hyperref[bp93]{2.32}). By Theorem \hyperref[bp101]{2.35}, applied to the schemes \(U\) and \(X\times_YV\), the cohomology of \(F\) on them is the étale cohomology of \(F_{\acute{e}t}\). As \(U=\varprojlim_iU_i\) and \(X\times_YV=\varprojlim_iX\times_YV_i\) are limits of qcqs schemes with affine transition maps, [Stacks, \href{https://stacks.math.columbia.edu/tag/03Q6}{03Q6}] identifies this with the colimit of the étale cohomology on the \(U_i\), respectively on the \(X\times_YV_i\). Each \(U_i\) is affine of finite type over \(k\) with \(\dim(U_i)\leq\dim(X)\), so by [Stacks, \href{https://stacks.math.columbia.edu/tag/0F0V}{0F0V}] we have \(\operatorname{cd}_\ell(U_i)\leq d\) for sheaves of \(A\)-modules. To bound cohomology on \(X\times_YV_i\) independently of \(i\), we choose a finite affine open cover of \(X\), with \(r\) members. Since \(X\) and \(Y\) are separated over \(k\), its intersections and their pullbacks to the affine \(V_i\) are affine; each pullback is étale over the corresponding intersection and has dimension at most \(\dim(X)\). The alternating Čech complex has length \(r-1\), so the Čech spectral sequence gives the bound \(c=d+r-1\). The last assertion is [Stacks, \href{https://stacks.math.columbia.edu/tag/0F10}{0F10}], which applies because \(g\) is of finite type over a field with \(\operatorname{cd}_\ell(k)<\infty\) (Lemma \hyperref[bp137]{2.46}).
\end{proof}

% object: lemma-truncation-pro-etale-sections (blueprint/chap-The-pro-etale-site/sec-Unbounded-complexes/lemma-truncation-pro-etale-sections/_section.md)
\phantomsection\label{bp142}%
\begin{bpstatementit}{Lemma~2.48}{ (Truncation on the pro-étale site)}
Let \(A\) be a finite commutative \(\ell\)-primary ring, let \(W\) be an object of \(X_{\mathrm{pro}\acute{e}t}\), and let \(e\) be an integer such that \(H^q(W,F)=0\) for \(q>e\) and every classical sheaf \(F\) of \(A\)-modules on \(X_{\mathrm{pro}\acute{e}t}\). Let \(L\in D(X_{\mathrm{pro}\acute{e}t},A)\) have classical cohomology sheaves, with \(H^q(L)=0\) for \(q>b\). Then \(H^jR\Gamma(W,L)=0\) for \(j>b+e+1\).
\end{bpstatementit}

% proof: proof/lemma-truncation-pro-etale-sections
\begin{proof}[{Proof}]\phantomsection\label{bp143}%
By left completeness (Proposition \hyperref[bp75]{2.26}), \(L\simeq R\varprojlim_n\tau_{\geq-n}L\), and derived sections preserve derived inverse limits, so \(R\Gamma(W,L)\simeq R\varprojlim_nR\Gamma(W,\tau_{\geq-n}L)\). Each \(\tau_{\geq-n}L\) is bounded with classical cohomology sheaves in degrees at most \(b\), so by the hypothesis and the hypercohomology spectral sequence, \(R\Gamma(W,\tau_{\geq-n}L)\) has no cohomology above \(b+e\). The Milnor sequence for this tower of complexes of abelian groups, proved as in Proposition \hyperref[bp69]{2.24}, expresses \(H^j\) of the derived limit through \(\varprojlim^1\) of \(H^{j-1}\) and \(\varprojlim\) of \(H^j\) of the terms. Both vanish for \(j>b+e+1\).
\end{proof}

% object: proposition-unbounded-comparison-finite-coefficients (blueprint/chap-The-pro-etale-site/sec-Unbounded-complexes/proposition-unbounded-comparison-finite-coefficients/_section.md)
\phantomsection\label{bp144}%
\begin{bpstatementit}{Proposition~2.49}{ (Unbounded comparison for finite coefficients)}
Let \(X\) and \(Y\) be of finite type over the base field \(k\), let \(A\) be a finite commutative \(\ell\)-primary ring, and let \(g:X\to Y\) be of finite type. Then \(\nu_X^*:D(X_{\acute{e}t},A)\to D(X_{\mathrm{pro}\acute{e}t},A)\) is fully faithful and preserves arbitrary direct sums. The map (2.9), now a natural transformation of functors \(D(X_{\acute{e}t},A)\to D(Y_{\mathrm{pro}\acute{e}t},A)\), is an isomorphism for every \(K\in D(X_{\acute{e}t},A)\).
\end{bpstatementit}

% proof: proof/proposition-unbounded-comparison-finite-coefficients
\begin{proof}[{Proof}]\phantomsection\label{bp145}%
By the two lemmas, both sides in a fixed degree depend only on a bounded-below truncation of \(K\), where Theorems \hyperref[bp101]{2.35} and \hyperref[bp111]{Theorem 2.38} apply. Let \(d\) and \(c\) be as in Lemma \hyperref[bp140]{2.47}, fix \(K\in
D(X_{\acute{e}t},A)\) and a degree \(m\), and consider the triangle \(\tau_{<N}K\to K\to\tau_{\geq N}K\) for an integer \(N\) chosen below.

Full faithfulness of \(\nu_X^*\) means that the unit \(K\to R\nu_{X,*}\nu_X^*K\) is an isomorphism, and it suffices to check this on the maps \(H^mR\Gamma(U,K)\to H^mR\Gamma(U,\nu_X^*K)\) for affine étale \(U\) over \(X\). We first reduce to bounded-below complexes. On the étale side, \(\Gamma(U,-)\) has cohomological dimension at most \(d\) on sheaves of \(A\)-modules, and by [Stacks, \href{https://stacks.math.columbia.edu/tag/07K7}{07K7}] the group \(H^mR\Gamma(U,K)\) depends only on \(\tau_{\geq N}K\) as soon as \(N\leq m-d\). On the pro-étale side, since \(\nu_X^*\) is exact it commutes with truncations, and Lemma \hyperref[bp142]{2.48}, applied with \(e=d\) to \(L=\nu_X^*\tau_{<N}K\) and \(b=N-1\), shows that \(H^m\) and \(H^{m+1}\) of \(R\Gamma(U,L)\) vanish as soon as \(N<m-d\). Hence for such \(N\) both sides are unchanged when \(K\) is replaced by \(\tau_{\geq N}K\), and for this bounded-below complex the map is an isomorphism by Theorem \hyperref[bp101]{2.35}.

For direct images, we evaluate (2.9) on a pro-étale affine \(V\) over \(Y\); these objects form a basis of \(Y_{\mathrm{pro}\acute{e}t}\). As before, we bound the contribution of \(\tau_{<N}K\) on both sides. On the étale side, the complex \(Rg_{\acute{e}t,*}\tau_{<N}K\) has cohomology sheaves in degrees at most \(N-1+\operatorname{cd}_\ell(g)\), so \(\nu_Y^*Rg_{\acute{e}t,*}\tau_{<N}K\) has classical cohomology in the same range, and Lemma \hyperref[bp142]{2.48} on \(Y\), with \(e=\dim(Y)+\operatorname{cd}_\ell(k)\) (Lemma \hyperref[bp140]{2.47} for \(Y\)), shows that its derived sections over \(V\) vanish above degree \(N+\operatorname{cd}_\ell(g)+\dim(Y)+\operatorname{cd}_\ell(k)\). On the other side, \(R\Gamma(V,Rg_{\mathrm{pro}\acute{e}t,*}\nu_X^*\tau_{<N}K)=R\Gamma(X\times_YV,\nu_X^*\tau_{<N}K)\), and the same lemma with \(W=X\times_YV\) and \(e=c\) gives vanishing above degree \(N+c\). Therefore, choosing \(N\) so small that \(H^m\) and \(H^{m+1}\) vanish on both sides, both sides of (2.9), evaluated on \(V\) in degree \(m\), are unchanged when \(K\) is replaced by \(\tau_{\geq N}K\), and for this bounded-below complex the map is an isomorphism by Theorem \hyperref[bp111]{2.38}.

The functor \(\nu_X^*\) preserves direct sums because it is an exact left adjoint.
\end{proof}

% text: prose-02-47
\phantomsection\label{bp146}%
Since by Theorem \hyperref[bp93]{2.32} the classical sheaves of \(A\)-modules are the essential image of \(\nu_X^*\) on sheaves, and \(\nu_X^*\) commutes with direct sums of sheaves, classical sheaves of \(A\)-modules are closed under arbitrary direct sums in \(\operatorname{Mod}(X_{\mathrm{pro}\acute{e}t},A)\).


% object: sec-Coefficient-rings (blueprint/chap-The-pro-etale-site/sec-Coefficient-rings/_section.md), content hidden
\subsection*{\texorpdfstring{2.12\quad Coefficient rings}{2.12 Coefficient rings}}
\phantomsection\label{bp147}%
\addcontentsline{toc}{subsection}{\texorpdfstring{2.12\quad Coefficient rings}{2.12 Coefficient rings}}
% text: prose-02-48
\phantomsection\label{bp148}%
We want pullback along morphisms and change of coefficients to be compatible without any bookkeeping. For this we define each coefficient ring once, on a point, and pull it back to every scheme along a single morphism of topoi. This gives \(A_X\), \(\mathcal O_{E,X}\), \(E_X\) and \(\QbarX\). Here \(A\) is a finite commutative ring killed by a power of \(\ell\), with the discrete topology, and \(A_X\) is the usual constant sheaf. For the point we take the site of profinite sets described below. By Proposition \hyperref[bp151]{2.50} it is the pro-étale site of a separably closed field. Unlike the étale site of the field, it keeps the topology of the coefficient rings.

% object: subsec-Condensed-rings (blueprint/chap-The-pro-etale-site/sec-Coefficient-rings/subsec-Condensed-rings/_section.md), content hidden
\subsubsection*{\texorpdfstring{2.12.1\quad Condensed rings}{2.12.1 Condensed rings}}
\phantomsection\label{bp149}%
\addcontentsline{toc}{subsubsection}{\texorpdfstring{2.12.1\quad Condensed rings}{2.12.1 Condensed rings}}
% text: prose-02-49
\phantomsection\label{bp150}%
The constant sheaf with value the abstract ring \(\mathbf Z_\ell\) forgets the topology of \(\mathbf Z_\ell\), and on \(\operatorname{Spec}\mathbf F_q\) it has \(H^1=0\) (Section \bpout{1.2}). We therefore need a notion of coefficient ring on a geometric point that retains the topology. The pro-étale site of a separably closed field is the site of profinite sets (Proposition \hyperref[bp151]{2.50} below). So a coefficient ring must at least be a sheaf of rings on profinite sets, that is, a condensed ring. The continuous maps from profinite sets into \(\mathbf Z_\ell\) form such a sheaf. With it, \(H^1\) of \(\mathbf F_q\) becomes \(\mathbf Z_\ell\) (Example \hyperref[bp176]{2.57}).

Let \(*_{\mathrm{pro}\acute{e}t}\) be the site of profinite sets, with finite jointly surjective families as covers. A \emph{condensed ring} is a sheaf of rings on \(*_{\mathrm{pro}\acute{e}t}\). A topological ring \(V\) defines a condensed ring \(\underline V\), whose sections over \(S\) are the continuous maps \(S\to V\). This is a sheaf because a surjection of profinite sets is a quotient map.

% object: proposition-pro-etale-site-separably-closed (blueprint/chap-The-pro-etale-site/sec-Coefficient-rings/subsec-Condensed-rings/proposition-pro-etale-site-separably-closed/_section.md)
\phantomsection\label{bp151}%
\begin{bpstatementit}{Proposition~2.50}{ (The pro-étale site of a separably closed field)}
Let \(\Omega\) be a separably closed field. The category of pro-étale affines over \(\operatorname{Spec}\Omega\) is equivalent to the category of profinite sets, with finite jointly surjective families as covers. Consequently its topos is, up to a set-theoretic cutoff, the topos of condensed sets. A discrete ring \(A\) defines the classical constant sheaf, while a profinite or locally profinite coefficient ring defines its sheaf of continuous functions on profinite sets.
\end{bpstatementit}

% proof: proof/proposition-pro-etale-site-separably-closed
\begin{proof}[{Proof}]\phantomsection\label{bp152}%
Since a finitely presented étale \(\Omega\)-algebra is a finite product of copies of \(\Omega\), filtered colimits of such algebras correspond, under \(\operatorname{Spec}\), to inverse limits of finite sets. Under this correspondence, faithfully flat covering maps become surjections of profinite sets, and by quasi-compactness every cover can be refined by a finite one, so the sites are equivalent. The topoi then agree because the pro-étale affines form a basis and sheaves are determined by their values on them. Finally, the last assertion is the definition of the condensed ring \(\underline R\) of a topological ring \(R\) given above, whose sections over a profinite set \(S\) are the continuous maps \(S\to R\); for discrete \(R\) these are the locally constant functions, which are the sections of the classical constant sheaf.
\end{proof}

% text: prose-02-50
\phantomsection\label{bp153}%
The étale topos of \(\Omega\) is the topos of sets, whereas by the proposition its pro-étale topos is the topos of condensed sets. The final object of \(*_{\mathrm{pro}\acute{e}t}\) is weakly contractible, so taking global sections is exact. However, a condensed sheaf is not determined by its global sections. For classical sheaves geometric stalks still suffice, whereas a general sheaf of modules on the pro-étale site has to be evaluated on all weakly contractible affines (Corollary \hyperref[bp57]{2.20}). For example, the completed coefficient sheaf \(\mathbf Z_{\ell,X}\) of Example \hyperref[bp95]{2.33} is of this general kind: its sections on a profinite set \(S\) are all continuous functions \(S\to\mathbf Z_\ell\).

Fix an algebraic closure \(\overline{\mathbf Q}_\ell\) of \(\mathbf Q_\ell\). Throughout, \(E\) denotes a finite extension of \(\mathbf Q_\ell\) inside \(\overline{\mathbf Q}_\ell\), with ring of integers \(\mathcal O_E\) and uniformizer \(\pi_E\). For \(\mathcal O_E\) and \(E\) the \(\ell\)-adic topologies define condensed rings \(\underline{\mathcal O_E}\) and \(\underline E\). For \(\overline{\mathbf Q}_\ell\), however, the topology causes a problem. To see why, recall that classically a \(\overline{\mathbf Q}_\ell\)-sheaf is an \(E\)-sheaf for some finite \(E\) with scalars extended, and all arguments with lattices and reduction modulo \(\pi_E\) take place over that \(E\). They cannot be made over \(\overline{\mathbf Q}_\ell\) itself, because \(\overline{\mathbf Q}_\ell\) is not complete and its valuation ring is not a discrete valuation ring. But the sheaf of all continuous maps to \(\overline{\mathbf Q}_\ell\) has sections that are defined over no finite \(E\) (see below), so it would lose this reduction. To keep it available, we take instead the filtered colimit over the finite coefficient fields. Using the valuation topology on \(\overline{\mathbf Q}_\ell\), we put

\[
\underline{\overline{\mathbf Q}_\ell}(S)=
\{a:S\to\overline{\mathbf Q}_\ell\text{ continuous}:a(S)\subset E\text{ for some finite }E/\mathbf Q_\ell\}
=\varinjlim_E\underline E(S).
\]

This is already a sheaf, because a cover of a profinite set can be refined by a finite one, and finitely many finite extensions of \(\mathbf Q_\ell\) lie in a common one. Since filtered colimits of sheaves on \(*_{\mathrm{pro}\acute{e}t}\) are computed sectionwise, \(\underline{\overline{\mathbf Q}_\ell}=\varinjlim_E\underline E\) as condensed rings.

This condensed ring is smaller than \(\underline V\) for \(V=\overline{\mathbf Q}_\ell\) with the valuation topology. Indeed, algebraic numbers of unbounded degree can tend to \(0\). Such a sequence is a continuous map from a convergent sequence to \(\overline{\mathbf Q}_\ell\) with image in no finite extension. With our choice, Lemma \bpout{3.23} shows that every constructible \(\overline{\mathbf Q}_\ell\)-complex and every morphism between such complexes descends to a finite \(E\).


% object: subsec-Coefficient-sheaves (blueprint/chap-The-pro-etale-site/sec-Coefficient-rings/subsec-Coefficient-sheaves/_section.md), content hidden
\subsubsection*{\texorpdfstring{2.12.2\quad Coefficient sheaves}{2.12.2 Coefficient sheaves}}
\phantomsection\label{bp154}%
\addcontentsline{toc}{subsubsection}{\texorpdfstring{2.12.2\quad Coefficient sheaves}{2.12.2 Coefficient sheaves}}
% text: prose-02-51
\phantomsection\label{bp155}%
A profinite set \(S\) defines a pro-étale scheme \(\operatorname{Spec}LC(S,\mathbf Z)\) over \(\mathbf Z\), where \(LC(S,\mathbf Z)\) is the ring of locally constant functions. The functor \(S\mapsto\operatorname{Spec}LC(S,\mathbf Z)\) is continuous and preserves finite limits, as one checks on finite quotients of \(S\), so it induces a morphism of topoi \(\Sh((\operatorname{Spec}\mathbf Z)_{\mathrm{pro}\acute{e}t})\to
\Sh(*_{\mathrm{pro}\acute{e}t})\) [Stacks, \href{https://stacks.math.columbia.edu/tag/00X6}{00X6}]. Composing with the structure morphism of \(X\), we obtain

\[
p_X:\Sh(X_{\mathrm{pro}\acute{e}t})\longrightarrow\Sh(*_{\mathrm{pro}\acute{e}t}),
\]

and for a condensed ring \(\Lambda\) we write \(\Lambda_X=p_X^*\Lambda\). For a morphism \(f:X\to Y\) we have \(p_Y\circ f=p_X\), and hence a canonical isomorphism \(f^{-1}\Lambda_Y\simeq\Lambda_X\), compatible with identities and composition. Likewise, a morphism \(\Lambda\to\Lambda'\) of condensed rings pulls back to \(\Lambda_X\to\Lambda'_X\), which defines extension of scalars \(\Lambda'_X\otimes_{\Lambda_X}-\) from \(\Lambda_X\)-modules to \(\Lambda'_X\)-modules on \(X_{\mathrm{pro}\acute{e}t}\). Thus both pullback along morphisms and change of coefficients are obtained from \(p_X^*\).

The next lemma identifies these rings with the completed coefficient rings of classical \(\ell\)-adic cohomology.

% object: lemma-comparison-completed-coefficients (blueprint/chap-The-pro-etale-site/sec-Coefficient-rings/subsec-Coefficient-sheaves/lemma-comparison-completed-coefficients/_section.md)
\phantomsection\label{bp156}%
\begin{bpstatementit}{Lemma~2.51}{ (Comparison with completed coefficients)}
Write \(\mathcal O_{E,X}\), \(E_X\) and \(\QbarX\) for the pullbacks to \(X_{\mathrm{pro}\acute{e}t}\) of \(\underline{\mathcal O_E}\), \(\underline E\) and \(\underline{\overline{\mathbf Q}_\ell}\). There are canonical isomorphisms of sheaves of rings

\[
\mathcal O_{E,X}\simeq\varprojlim_n(\mathcal O_E/\pi_E^n)_X,
\qquad E_X\simeq\mathcal O_{E,X}[1/\ell],
\qquad \QbarX\simeq\varinjlim_E E_X,
\]

where \((\mathcal O_E/\pi_E^n)_X\) is the constant sheaf. They are compatible with inverse image along every morphism of schemes.
\end{bpstatementit}

% proof: proof/lemma-comparison-completed-coefficients
\begin{proof}[{Proof}]\phantomsection\label{bp157}%
The functor \(p_X^*\) commutes with colimits, so the second and third isomorphisms reduce to identities on the point. The first is an inverse limit, and since \(p_X^*\) need not commute with inverse limits, we use instead that \(\underline{\mathcal O_E}\) is representable.

On the point, \(\underline{\mathcal O_E}\) is represented by the profinite set \(\mathcal O_E=\varprojlim_n\mathcal O_E/\pi_E^n\), so its pullback to \(X\) is represented by the pro-étale \(X\)-scheme \(\varprojlim_n\coprod_{\mathcal O_E/\pi_E^n}X\). By the Yoneda lemma the sheaf this scheme represents is \(\varprojlim_n(\mathcal O_E/\pi_E^n)_X\), and the ring structures match.

A continuous map from a profinite set to \(E\) has compact image, and hence lands in \(\ell^{-m}\mathcal O_E\) for some \(m\). Thus \(\underline E=\underline{\mathcal O_E}[1/\ell]\) on the point. Since inverting \(\ell\) is a filtered colimit and \(p_X^*\) commutes with colimits, we obtain the second isomorphism. In the same way the third follows from \(\underline{\overline{\mathbf Q}_\ell}=\varinjlim_E\underline E\). All three isomorphisms are obtained by applying \(p_X^*\), and \(p_Y\circ f=p_X\) for every morphism \(f:X\to Y\), so they are compatible with pullback.
\end{proof}

% text: prose-02-52
\phantomsection\label{bp158}%
The first isomorphism describes \(\mathcal O_{E,X}\) as a limit of its finite quotients. To compute derived reductions modulo \(\pi_E^n\) we also need these quotients as cokernels of multiplication by \(\pi_E^n\), which the next corollary provides.

% object: corollary-finite-quotients-their-transition-maps (blueprint/chap-The-pro-etale-site/sec-Coefficient-rings/subsec-Coefficient-sheaves/corollary-finite-quotients-their-transition-maps/_section.md)
\phantomsection\label{bp159}%
\begin{bpstatementit}{Corollary~2.52}{ (Finite quotients and their transition maps)}
For every \(n\geq1\) there is a canonical exact sequence on \(X_{\mathrm{pro}\acute{e}t}\)

\[
0\longrightarrow\mathcal O_{E,X}
\xrightarrow{\pi_E^n}\mathcal O_{E,X}
\longrightarrow(\mathcal O_E/\pi_E^n)_X\longrightarrow0.
\]

Thus the quotient is represented by the two-term complex \([\mathcal O_{E,X}\xrightarrow{\pi_E^n}\mathcal O_{E,X}]\) in degrees \(-1,0\). Under this representation, reduction from level \(n+1\) to level \(n\) is the chain map \((\pi_E,\mathrm{id})\). These resolutions and maps commute with inverse image along every morphism of schemes.
\end{bpstatementit}

% proof: proof/corollary-finite-quotients-their-transition-maps
\begin{proof}[{Proof}]\phantomsection\label{bp160}%
It suffices to prove exactness on the point, since \(p_X^*\) is exact and inverse image along \(f:X\to Y\) satisfies \(p_Y\circ f=p_X\). Now, on a profinite set \(S\), multiplication by \(\pi_E^n\) is injective on continuous functions \(S\to\mathcal O_E\), and its image is the kernel of reduction. Moreover, reduction is surjective because a continuous function \(S\to\mathcal O_E/\pi_E^n\) takes finitely many values, and choosing a representative of each gives a continuous lift. Finally, the chain map \((\pi_E,\mathrm{id})\) commutes with the differentials because \(\pi_E^{n}\pi_E=\pi_E^{n+1}\), and its degree-zero component is the identity, so it induces the quotient projection.
\end{proof}

% text: prose-02-53
\phantomsection\label{bp161}%
Using the two-term resolution we see that, for every complex \(K\) of \(\mathcal O_{E,X}\)-modules, the derived reduction \(K\otimes^{\mathbf L}(\mathcal O_E/\pi_E^n)_X\) is the cone of \(\pi_E^n:K\to K\) (Lemma \hyperref[bp188]{2.62}).

We call the rings \(A,\mathbf Z_\ell,\mathcal O_E,\mathbf Q_\ell,E,\overline{\mathbf Q}_\ell\) the \emph{admissible} coefficient rings. Here \(E=\mathbf Q_\ell\) is allowed, so that \(\mathbf Z_\ell\) and \(\mathbf Q_\ell\) are the case \(E=\mathbf Q_\ell\) of \(\mathcal O_E\) and \(E\). For \(E=\mathbf Q_\ell\) we write \(\mathbf Z_{\ell,X}=\mathcal O_{\mathbf Q_\ell,X}\) and \(\mathbf Q_{\ell,X}=(\mathbf Q_\ell)_X\). By the comparison lemma above,

\[
\mathbf Z_{\ell,X}\simeq\varprojlim_n(\mathbf Z/\ell^n)_X,
\qquad
\mathbf Q_{\ell,X}\simeq\mathbf Z_{\ell,X}[1/\ell].
\]

For a condensed ring \(\Lambda\) we write \(\operatorname{Mod}(X_{\mathrm{pro}\acute{e}t},\Lambda)\) for the category of \(\Lambda_X\)-modules. Its derived category is \(D(X_{\mathrm{pro}\acute{e}t},\Lambda)\).

Duality and the trace map need, besides the constant sheaf, its twists by roots of unity. The reason is already visible on \(\mathbf P^1\) over an algebraically closed field \(k\), for \(n\) invertible in \(k\). The Kummer sequence \(0\to\mu_n\to\mathbf G_m\xrightarrow{n}\mathbf G_m\to0\) gives an isomorphism \(H^2(\mathbf P^1,\mu_n)\simeq\operatorname{Pic}(\mathbf P^1)/n=\mathbf Z/n\), given by the degree of a line bundle. This isomorphism involves no choice, and it sends the class of a point to \(1\). An identification of \(H^2(\mathbf P^1,\mathbf Z/n)\) with \(\mathbf Z/n\), on the other hand, requires an isomorphism \(\mu_n\simeq\mathbf Z/n\), that is, a choice of a primitive \(n\)th root of unity. Frobenius does not respect such a choice.

To see this, let \(k\) be an algebraic closure of \(\mathbf F_q\). The arithmetic Frobenius acts on \(\mu_{\ell^n}(k)\) by \(\zeta\mapsto\zeta^q\). The geometric Frobenius is the inverse of the arithmetic Frobenius in \(\operatorname{Gal}(k/\mathbf F_q)\). Hence it acts on \(\Lambda(1)\) by \(q^{-1}\), and on \(H^2(\mathbf P^1,\Lambda)\simeq\Lambda(-1)\) by \(q\). This agrees with \(\#\mathbf P^1(\mathbf F_q)=1+q\).

Over \(\mathbf C\) the twist is \(2\pi i\mathbf Z\). The exponential sequence identifies \(H^2(\mathbf P^1(\mathbf C),2\pi i\mathbf Z)\) with \(\mathbf Z\). An identification of \(H^2(\mathbf P^1(\mathbf C),\mathbf Z)\) with \(\mathbf Z\), by contrast, is a choice of orientation, and the complex orientation depends on the choice of \(i\).

With twists, the trace map of \(\mathbf P^1\) (Proposition \bpout{6.7}) and smooth purity \(f^!\simeq f^*(d)[2d]\) (Theorem \bpout{9.3}) are canonical and commute with Frobenius (Proposition \bpout{10.20}). This is where the factor \(q^d\) in Poincaré duality (c) comes from.

% object: definition-tate-twists (blueprint/chap-The-pro-etale-site/sec-Coefficient-rings/subsec-Coefficient-sheaves/definition-tate-twists/_section.md)
\phantomsection\label{bp162}%
\begin{bpstatement}{Definition~2.53}{ (Tate twists)}
Let \(X\) be a scheme on which \(\ell\) is invertible, and let \(\mu_{\ell^n,X}\) be the sheaf of \(\ell^n\)th roots of unity on \(X_{\mathrm{pro}\acute{e}t}\). Then it is locally isomorphic to \((\mathbf Z/\ell^n\mathbf Z)_X\). For a finite ring \(A\) with \(\ell^rA=0\) put \(A_X(1)=\mu_{\ell^r,X}\otimes_{\mathbf Z/\ell^r\mathbf Z}A_X\), which by the Kummer maps is independent of \(r\). Put \(\mathcal O_{E,X}(1)=\varprojlim_n(\mathcal O_E/\ell^n)_X(1)\), and obtain \(E_X(1)\) and \(\QbarX(1)\) by inverting \(\ell\) and passing to the colimit over \(E\); equivalently, \(\QbarX(1)=\QbarX\otimes_{\mathbf Z_{\ell,X}}\mathbf Z_{\ell,X}(1)\) with \(\mathbf Z_{\ell,X}(1)=\varprojlim_n\mu_{\ell^n,X}\). Each \(\Lambda_X(1)\) is locally free of rank one. For \(d\ge0\), \(\Lambda_X(d)\) is the \(d\)-fold tensor power, \(\Lambda_X(-d)\) its dual, and \(K(d)=K\otimes^{\mathbf L}_{\Lambda_X}\Lambda_X(d)\).
\end{bpstatement}


% object: subsec-Inverse-image (blueprint/chap-The-pro-etale-site/sec-Coefficient-rings/subsec-Inverse-image/_section.md), content hidden
\subsubsection*{\texorpdfstring{2.12.3\quad Inverse image}{2.12.3 Inverse image}}
\phantomsection\label{bp163}%
\addcontentsline{toc}{subsubsection}{\texorpdfstring{2.12.3\quad Inverse image}{2.12.3 Inverse image}}
% text: prose-02-54
\phantomsection\label{bp164}%
Every base change argument below compares a module over \(\Lambda_Y\) with its pullback to \(X\), and for this we want the pullback of a \(\Lambda_Y\)-module to be its inverse image as a sheaf, with no tensor product. For coherent sheaves this fails, since \(f^*M=\mathcal O_X\otimes_{f^{-1}\mathcal O_Y}f^{-1}M\) and \(f^{-1}\mathcal O_Y\ne\mathcal O_X\) in general. For our coefficient sheaves the two notions agree, because all coefficient sheaves are pulled back from the point. The same holds for their Tate twists, where \(\ell\) is invertible.

The second part of the next lemma says that extension of scalars to a finite extension \(E'/E\) is a finite direct sum. We use this, for instance, in the proof of Proposition \bpout{4.40}.

% object: lemma-inverse-image-coefficient-sheaves (blueprint/chap-The-pro-etale-site/sec-Coefficient-rings/subsec-Inverse-image/lemma-inverse-image-coefficient-sheaves/_section.md)
\phantomsection\label{bp165}%
\begin{bpstatementit}{Lemma~2.54}{ (Inverse image of the coefficient sheaves)}
Let \(f:X\to Y\) be a morphism of schemes.

\begin{enumerate}
\item For every admissible \(\Lambda\), the canonical map \[
   f^{-1}\Lambda_Y\longrightarrow\Lambda_X
   \] is an isomorphism of sheaves of rings. If \(\ell\) is invertible on \(Y\), the same holds for its Tate twists, including the twists of a finite ring \(A\).
\item If \(E\subset E'\) is a finite extension, then \(\mathcal O_{E',X}\) is a finite free \(\mathcal O_{E,X}\)-module and \(E'_X\) is a finite free \(E_X\)-module.
\end{enumerate}
\end{bpstatementit}

% proof: proof/lemma-inverse-image-coefficient-sheaves
\begin{proof}[{Proof}]\phantomsection\label{bp166}%
(1) For every condensed ring \(\Lambda\), the identity \(p_Y\circ f=p_X\) shows that

\[
f^{-1}\Lambda_Y=f^{-1}p_Y^*\Lambda\simeq p_X^*\Lambda=\Lambda_X,
\]

and this proves the assertion for the coefficient rings themselves.

For the twists, \(\ell\) is invertible on \(Y\), hence on \(X\), as Definition \hyperref[bp162]{2.53} requires. The torsion twists \(\mu_{\ell^n}\) are then represented by finite étale group schemes, and these commute with base change. Their inverse limit is represented by the inverse limit of these schemes, so that its inverse image is represented by the base change [Stacks, \href{https://stacks.math.columbia.edu/tag/09A5}{09A5}]. This proves the assertion for \(\mathcal O_E(1)\). For a finite \(A\) killed by \(\ell^n\), its twist is obtained from the twist at finite level by extending scalars from \(\mathbf Z/\ell^n\mathbf Z\), so inverse image commutes with it as well. The other twists are obtained by inverting \(\ell\), extending scalars, and taking tensor powers and duals.

(2) An \(\mathcal O_E\)-basis of \(\mathcal O_{E'}\) is an isomorphism of topological \(\mathcal O_E\)-modules \(\mathcal O_{E'}\simeq\mathcal O_E^{[E':E]}\), hence of condensed modules \(\underline{\mathcal O_{E'}}\simeq\underline{\mathcal O_E}^{[E':E]}\). Pull it back along \(p_X\), and invert \(\ell\) for the fields.
\end{proof}

% text: prose-02-55
\phantomsection\label{bp167}%
By (1), the pullback of modules over these sheaves, defined in Section \bpout{3.1} below, is the exact functor \(f^{-1}\).



% object: sec-Continuous-cohomology (blueprint/chap-The-pro-etale-site/sec-Continuous-cohomology/_section.md), content hidden
\subsection*{\texorpdfstring{2.13\quad Continuous cohomology}{2.13 Continuous cohomology}}
\phantomsection\label{bp168}%
\addcontentsline{toc}{subsection}{\texorpdfstring{2.13\quad Continuous cohomology}{2.13 Continuous cohomology}}
% text: prose-02-56
\phantomsection\label{bp169}%
The comparison theorems above concern discrete coefficients, whereas the integral and rational \(\ell\)-adic coefficient rings of Section \hyperref[bp147]{2.12} are defined directly on the pro-étale site, because the discrete constant sheaves with the same values have the wrong cohomology (Example \hyperref[bp176]{2.57}). The main result of this section, Proposition \hyperref[bp170]{2.55}, says that pro-étale cohomology with coefficients \(\mathbf Z_\ell\) is the classical continuous cohomology, that is, the derived limit of cohomology with coefficients \(\mathbf Z/\ell^n\).

The sheaf \(\mathbf Z_{\ell,X}\) is the inverse limit of the classical sheaves \((\mathbf Z/\ell^n)_X\), with surjective transition maps (Lemma \hyperref[bp156]{2.51}). For such a limit, the cohomology is the derived limit of the cohomology of the terms. When the terms are classical, this is continuous étale cohomology.

% object: proposition-continuous-cohomology-pro-etale-cohomology (blueprint/chap-The-pro-etale-site/sec-Continuous-cohomology/proposition-continuous-cohomology-pro-etale-cohomology/_section.md)
\phantomsection\label{bp170}%
\begin{bpstatementit}{Proposition~2.55}{ (Continuous cohomology as pro-étale cohomology)}
Let \(X\) be a scheme and \((F_n)\) an inverse system of abelian sheaves on \(X_{\mathrm{pro}\acute{e}t}\) with epimorphic transition maps. Put \(F=\varprojlim_nF_n\) in the pro-étale topos. Then, in the derived category of abelian groups,

\[
R\Gamma(X_{\mathrm{pro}\acute{e}t},F)
\simeq R\varprojlim_n R\Gamma(X_{\mathrm{pro}\acute{e}t},F_n).
\]
\end{bpstatementit}

% proof: proof/proposition-continuous-cohomology-pro-etale-cohomology
\begin{proof}[{Proof}]\phantomsection\label{bp171}%
Since the transition maps are epimorphisms, \(F\simeq R\varprojlim_nF_n\) by the last assertion of Proposition \hyperref[bp69]{2.24}. Moreover, \(R\Gamma(X_{\mathrm{pro}\acute{e}t},-)\) commutes with derived inverse limits since it is a right adjoint.
\end{proof}

% text: prose-02-57
\phantomsection\label{bp172}%
When the \(F_n\) come from the étale site, as the sheaves \((\mathbf Z/\ell^n)_X\) do, each term on the right is étale cohomology by the comparison theorem, and the proposition identifies pro-étale cohomology of the limit with continuous étale cohomology \(R\varprojlim_nR\Gamma(X_{\acute{e}t},F_n)\).

% object: corollary-continuous-etale-cohomology (blueprint/chap-The-pro-etale-site/sec-Continuous-cohomology/corollary-continuous-etale-cohomology/_section.md)
\phantomsection\label{bp173}%
\begin{bpstatementit}{Corollary~2.56}{ (Continuous étale cohomology)}
In Proposition \hyperref[bp170]{2.55}, suppose that \(F_n=\nu_X^*F_{n,\acute{e}t}\) for an inverse system \((F_{n,\acute{e}t})\) of abelian sheaves on \(X_{\acute{e}t}\) with epimorphic transition maps. Then \(R\Gamma(X_{\mathrm{pro}\acute{e}t},F)\) is the continuous étale cohomology of \((F_{n,\acute{e}t})\): the right derived functor of \((G_n)\mapsto\Gamma(X_{\acute{e}t},\varprojlim_nG_n)\) on the abelian category of inverse systems, evaluated at \((F_{n,\acute{e}t})\).
\end{bpstatementit}

% proof: proof/corollary-continuous-etale-cohomology
\begin{proof}[{Proof}]\phantomsection\label{bp174}%
By Theorem \hyperref[bp101]{2.35}, each \(R\Gamma(X_{\mathrm{pro}\acute{e}t},F_n)\) is \(R\Gamma(X_{\acute{e}t},F_{n,\acute{e}t})\). The derived functor of \((G_n)\mapsto\Gamma(X_{\acute{e}t},\varprojlim_nG_n)\) is the composite of the derived functors of sections and of inverse limit, so its value at \((F_{n,\acute{e}t})\) is \(R\varprojlim_nR\Gamma(X_{\acute{e}t},F_{n,\acute{e}t})\) [\hyperref[bp191]{BS15}, Proposition 5.6.2].
\end{proof}

% text: prose-02-58
\phantomsection\label{bp175}%
Applying the Milnor sequence (2.2) to the tower of complexes \(R\Gamma(X_{\mathrm{pro}\acute{e}t},F_n)\) and using the proposition, we obtain for every \(q\) a natural exact sequence

\[
0\longrightarrow R^1\varprojlim_nH^{q-1}(X_{\mathrm{pro}\acute{e}t},F_n)
\longrightarrow H^q(X_{\mathrm{pro}\acute{e}t},F)
\longrightarrow\varprojlim_nH^q(X_{\mathrm{pro}\acute{e}t},F_n)
\longrightarrow0.
\tag{2.11}
\]

The first term can be nonzero, because the transition maps between the cohomology groups need not be surjective even if those between the sheaves \(F_n\) are. It vanishes if the groups \(H^{q-1}(X_{\mathrm{pro}\acute{e}t},F_n)\) satisfy the Mittag–Leffler condition. This holds for instance if they are finite, since then descending chains of images stabilize at each level.

% object: example-first-cohomology-finite-field (blueprint/chap-The-pro-etale-site/sec-Continuous-cohomology/example-first-cohomology-finite-field/_section.md)
\phantomsection\label{bp176}%
\begin{bpstatement}{Example~2.57}{ (The first cohomology of a finite field)}
For \(X=\operatorname{Spec}\mathbf F_q\) and \(\ell\nmid q\),

\[
H^1(X_{\mathrm{pro}\acute{e}t},\mathbf Z_{\ell,X})\simeq\mathbf Z_\ell,
\qquad H^1(X_{\mathrm{pro}\acute{e}t},\mathbf Q_{\ell,X})\simeq\mathbf Q_\ell.
\]

At each finite level the first cohomology is \(\operatorname{Hom}_{\mathrm{cont}}(\widehat{\mathbf Z},
\mathbf Z/\ell^n)=\mathbf Z/\ell^n\), and the groups in degree zero are finite, so that their \(\varprojlim^1\) vanishes. Hence the sequence (2.11) for \(F_n=(\mathbf Z/\ell^n)_X\) proves the integral statement. Since \(X\) is qcqs, its pro-étale cohomology commutes with filtered colimits [Stacks, \href{https://stacks.math.columbia.edu/tag/09A3}{09A3}], in particular with inverting \(\ell\), and the rational statement follows. For the constant étale sheaf with value the discrete group \(\mathbf Q_\ell\), on the other hand, the first cohomology is \(\operatorname{Hom}_{\mathrm{cont}}(\widehat{\mathbf Z},\mathbf Q_\ell)=0\), because a continuous homomorphism from a profinite group to a discrete torsion-free group has finite, hence zero, image.
\end{bpstatement}


% object: sec-Extension-by-zero-and-completion (blueprint/chap-The-pro-etale-site/sec-Extension-by-zero-and-completion/_section.md), content hidden
\subsection*{\texorpdfstring{2.14\quad Extension by zero and completion}{2.14 Extension by zero and completion}}
\phantomsection\label{bp177}%
\addcontentsline{toc}{subsection}{\texorpdfstring{2.14\quad Extension by zero and completion}{2.14 Extension by zero and completion}}
% text: prose-02-59
\phantomsection\label{bp178}%
The integral coefficient sheaves are inverse limits. To restrict integral complexes to the strata of a stratification, and to extend them by zero, we therefore need the functors \(j_!\), \(j^*\), \(i^*\) and \(i_*\) to commute with limits. On the pro-étale site this holds, even for inverse image along a closed immersion, and from this we deduce the localization triangles and construct the local cohomology functor. With these tools we can argue with closed supports before the exceptional inverse image and proper base change are available.

That \(i^*\) commutes with limits is special to the pro-étale site. On \(X_{\acute{e}t}\) the stalk of \(i^*F\) at a point of \(Z\) is the stalk of \(F\). This is a colimit, and colimits do not commute with infinite products. On \(X_{\mathrm{pro}\acute{e}t}\), by contrast, the value of \(i^*M\) on a weakly contractible affine over \(Z\) is the value of \(M\) on one affine object of \(X_{\mathrm{pro}\acute{e}t}\), henselian along \(Z\) (Proposition \hyperref[bp179]{2.58}), and evaluation on one object preserves limits. For \(\mathcal O_{E,X}=\varprojlim(\mathcal O_E/\pi_E^n)_X\) this means that \(i^*\) applied to the limit is the limit of the restrictions \((\mathcal O_E/\pi_E^n)_Z\). The same holds for any integral complex given as a derived limit of its reductions.

A complex over \(\mathcal O_E\) is determined by its reductions modulo \(\pi_E^n\) when it is derived complete (Definition \hyperref[bp186]{2.61} below). Combined with the computation of the cohomology of these reductions on the étale site, this is how we reduce integral statements to finite coefficients throughout the notes.

% object: proposition-extension-zero-restriction-closed-subscheme (blueprint/chap-The-pro-etale-site/sec-Extension-by-zero-and-completion/proposition-extension-zero-restriction-closed-subscheme/_section.md)
\phantomsection\label{bp179}%
\begin{bpstatementit}{Proposition~2.58}{ (Extension by zero and restriction to a closed subscheme)}
Let \(X\) be quasi-compact and quasi-separated, \(j:U\hookrightarrow X\) a quasi-compact open immersion, and \(i:Z\hookrightarrow X\) its closed complement. For a sheaf of rings \(\mathcal R\) on \(X_{\mathrm{pro}\acute{e}t}\), we use its inverse images as the coefficient rings on \(U\) and \(Z\). Thus \(i^*\) goes from \(\mathcal R\)-module sheaves on \(X_{\mathrm{pro}\acute{e}t}\) to \(i^*\mathcal R\)-module sheaves on \(Z_{\mathrm{pro}\acute{e}t}\), \(i_*\) goes back, and \(j^*\), \(j_!\) relate \(\mathcal R\)-module sheaves on \(X_{\mathrm{pro}\acute{e}t}\) and \(j^*\mathcal R\)-module sheaves on \(U_{\mathrm{pro}\acute{e}t}\). Then \(i^*\) is exact and preserves all small limits and colimits; \(i_*\) is exact and fully faithful. The functors \(j^*\) and \(j_!\) are exact, and every \(\mathcal R\)-module sheaf \(M\) has the exact sequence

\[
0\longrightarrow j_!j^*M\longrightarrow M
\longrightarrow i_*i^*M\longrightarrow0.
\]

Moreover \(j_!\) preserves products and derived inverse limits of module complexes. Finally, for a morphism \(g:X'\to X\) with base changes \(j':U'\hookrightarrow X'\) and \(g':U'\to U\), the canonical map \(j'_!g'^*\to g^*j_!\) is an isomorphism.
\end{bpstatementit}

% proof: proof/proposition-extension-zero-restriction-closed-subscheme
\begin{proof}[{Proof}]\phantomsection\label{bp180}%
The special feature of the pro-étale site is that \(i^*\) preserves limits; we prove this first, deduce exactness of \(j_!\) from the localization sequence, obtain the statement about products and derived limits from a triangle expressing \(j_!\) through \(Rj_*\), \(i^*\) and \(i_*\), and check base change by adjunction.

\emph{Step 1: the functors \(i^*\) and \(i_*\).} The functor \(i^*\) preserves limits, because evaluating \(i^*M\) on a weakly contractible affine over \(Z\) is the same as evaluating \(M\) on an affine object of \(X_{\mathrm{pro}\acute{e}t}\), and evaluation commutes with limits. The assertion is local on \(X\). Let \(V\) be a weakly contractible object of the affine ind-étale basis of \(Z\). It determines an affine object \(v(V)\) in the corresponding basis of \(X\), henselian along its inverse image of \(Z\): the functor \(v\) is the fully faithful left adjoint to base change of these affine bases, and \(v(V)\times_XZ\simeq V\) [Stacks, \href{https://stacks.math.columbia.edu/tag/09A9}{09A9}]. The pullback of a sheaf, evaluated on \(V\), is its evaluation on \(v(V)\), because every covering of a weakly contractible object splits, so that sheafification does not change values there. Since evaluation commutes with limits and these objects detect isomorphisms, \(i^*\) preserves limits. Being a left adjoint, it also preserves colimits. Since it preserves both limits and colimits, it is exact. Lifting covers along \(v\) shows that \(i_*\) is exact and that \(i^*i_*=\mathrm{id}\) and \(j^*i_*=0\). Finally, the localization sequence is exact on the pro-étale site [Stacks, \href{https://stacks.math.columbia.edu/tag/09AH}{09AH}], that is, the adjunction map \(M\to i_*i^*M\) is surjective with kernel \(j_!j^*M\).

\emph{Step 2: exactness of \(j^*\) and \(j_!\).} The functor \(j^*\) is restriction to \(U_{\mathrm{pro}\acute{e}t}\), hence exact. The functor \(j_!\) is exact because its composites \(j^*j_!\) and \(i^*j_!\) are exact and the pair \((j^*,i^*)\) detects exactness. Here the first composite is \(j^*j_!=\mathrm{id}\), and the second is \(i^*j_!=0\). Indeed, in the localization sequence of \(j_!N\) the first map \(j_!j^*j_!N\to j_!N\) is an isomorphism since \(j^*j_!N=N\). Hence \(i_*i^*j_!N=0\), and \(i^*j_!N=0\) because \(i_*\) is fully faithful. The pair \((j^*,i^*)\) detects exactness, because both functors are exact, hence commute with taking homology of a complex of sheaves, while a sheaf \(M\) with \(j^*M=0\) and \(i^*M=0\) is zero by the localization sequence.

\emph{Step 3: products and derived inverse limits.} For \(K\in D(U_{\mathrm{pro}\acute{e}t},j^*\mathcal R)\), applying the localization sequence to a complex representing \(Rj_*K\) and using \(j^*Rj_*K\simeq K\), we obtain the triangle

\[
j_!K\longrightarrow Rj_*K\longrightarrow i_*i^*Rj_*K
\longrightarrow j_!K[1].
\]

The last two terms preserve products and derived inverse limits, because \(Rj_*\) and \(i_*\) are right adjoints and \(i^*\) is exact and preserves limits by Step 1; hence so does their fibre \(j_!\). As \(j_!\) is exact and products are exact, the corresponding statements about limits of sheaves of modules follow.

\emph{Step 4: base change.} The map \(j'_!g'^*\to g^*j_!\) is an isomorphism because its right adjoint is. Indeed, \(j'_!g'^*\) is left adjoint to \(g'_*j'^*\) and \(g^*j_!\) is left adjoint to \(j^*g_*\), and the right adjoint of the map is the base change map \(j^*g_*\to g'_*j'^*\), which is an isomorphism since both sides send a sheaf on \(X'_{\mathrm{pro}\acute{e}t}\) to its restriction to the objects lying over \(U'\). The same holds on derived categories, since the coefficient rings are inverse images, so that \(g^*\), \(g'^*\) and \(j'_!\) are exact. The argument with the functor \(v\) is [\hyperref[bp191]{BS15}, §6.1]; see also [Stacks, \href{https://stacks.math.columbia.edu/tag/09AA}{09AA}, \href{https://stacks.math.columbia.edu/tag/09AH}{09AH}, \href{https://stacks.math.columbia.edu/tag/09BM}{09BM}] and [Stacks, \href{https://stacks.math.columbia.edu/tag/09BP}{09BP}].
\end{proof}

% text: prose-02-60
\phantomsection\label{bp181}%
For arguments with closed supports we need a right adjoint of \(i_*\) before \(i^!\) is constructed by Brown representability. We build it from \(i^*\), \(j^*\) and \(Rj_*\) alone. We keep the notation of Proposition \hyperref[bp179]{2.58}.

% object: definition-local-cohomology (blueprint/chap-The-pro-etale-site/sec-Extension-by-zero-and-completion/definition-local-cohomology/_section.md)
\phantomsection\label{bp182}%
\begin{bpstatement}{Definition~2.59}{ (Local cohomology)}
The \emph{local cohomology functor} \(R\underline\Gamma_Z:D(X_{\mathrm{pro}\acute{e}t},\mathcal R)\to D(Z_{\mathrm{pro}\acute{e}t},i^*\mathcal R)\) is the right derived functor of the functor of sections supported on \(Z\),

\[
F\longmapsto i^*\ker(F\longrightarrow j_*j^*F),
\]

where \(F\to j_*j^*F\) is the unit of \(j^*\dashv j_*\).
\end{bpstatement}

% object: lemma-local-cohomology-right-adjoint (blueprint/chap-The-pro-etale-site/sec-Extension-by-zero-and-completion/lemma-local-cohomology-right-adjoint/_section.md)
\phantomsection\label{bp183}%
\begin{bpstatementit}{Lemma~2.60}{ (Local cohomology is right adjoint to $i_*$)}
The functor \(R\underline\Gamma_Z\) is right adjoint to the exact fully faithful functor \(i_*:D(Z_{\mathrm{pro}\acute{e}t},i^*\mathcal R)\to D(X_{\mathrm{pro}\acute{e}t},\mathcal R)\). It is computed by

\[
R\underline\Gamma_ZK\simeq i^*\fib(K\longrightarrow Rj_*j^*K),
\]

where \(K\to Rj_*j^*K\) is the unit of \(j^*\dashv Rj_*\).
\end{bpstatementit}

% proof: proof/lemma-local-cohomology-right-adjoint
\begin{proof}[{Proof}]\phantomsection\label{bp184}%
We first compute \(R\underline\Gamma_Z\). Represent \(K\) by a K-injective complex \(J\) with injective terms, so that \(Rj_*j^*K\) is represented by \(j_*j^*J\). The unit \(J\to j_*j^*J\) is surjective in each degree: for an injective \(\mathcal R\)-module \(I\) and an object \(V\), a section of \(j_*j^*I\) over \(V\) is a map \(j_!j^*\mathcal R_V\to I\), and it extends along the monomorphism \(j_!j^*\mathcal R_V\to\mathcal R_V\). Hence the kernel of \(J\to j_*j^*J\), taken termwise, represents the fibre of the unit, and applying the exact functor \(i^*\) gives \(R\underline\Gamma_ZK\simeq i^*\fib(K\to Rj_*j^*K)\).

Let \(F=\fib(K\to Rj_*j^*K)\). Since \(j^*i_*M=0\), we have \(\operatorname{Hom}(i_*M,Rj_*j^*K)=\operatorname{Hom}(j^*i_*M,j^*K)=0\) in all degrees, so maps \(i_*M\to K\) are the same as maps \(i_*M\to F\). The restriction \(j^*F\) vanishes, because \(j^*\) of the unit is an isomorphism, so the localization sequence gives \(F\simeq i_*i^*F\), and \(\operatorname{Hom}(i_*M,i_*i^*F)=\operatorname{Hom}(M,i^*F)\) since \(i_*\) is fully faithful.
\end{proof}

% text: prose-02-61
\phantomsection\label{bp185}%
Once \(i^!\) exists, it agrees with \(R\underline\Gamma_Z\) on the localizing category on which \(i^!\) is defined, since both are right adjoints of \(i_*\) there.

We treat integral coefficients through their reductions modulo powers of \(\pi_E\). This is possible for complexes that are recovered from their reductions, the derived complete ones. Derived completeness is the analogue for complexes of the condition \(M=\varprojlim_nM/\pi_E^nM\) on a module. For example, \(\mathcal O_{E,X}\) itself is derived complete. The sheaf \(E_X\) is not, since all its reductions vanish.

% object: definition-derived-adic-completeness (blueprint/chap-The-pro-etale-site/sec-Extension-by-zero-and-completion/definition-derived-adic-completeness/_section.md)
\phantomsection\label{bp186}%
\begin{bpstatement}{Definition~2.61}{ (Derived adic completeness)}
Let \(E/\mathbf Q_\ell\) be finite, with uniformizer \(\pi_E\), and let \(\mathcal O_{E,X}=p_X^*\underline{\mathcal O_E}\) and \((\mathcal O_E/\pi_E^n)_X\) be the integral and finite coefficient sheaves of Section \hyperref[bp147]{2.12}. A complex \(K\in D(X_{\mathrm{pro}\acute{e}t},\mathcal O_E)\) is \emph{derived complete} if its canonical map

\[
K\longrightarrow R\varprojlim_n
(K\otimes^{\mathbf L}_{\mathcal O_{E,X}}(\mathcal O_E/\pi_E^{n})_X)
\]

is an isomorphism.
\end{bpstatement}

% text: prose-02-62
\phantomsection\label{bp187}%
We write \(K_n=K\otimes^{\mathbf L}_{\mathcal O_{E,X}}(\mathcal O_E/\pi_E^{n})_X\) for the reductions of \(K\). The canonical transition maps \(K_{n+1}\to K_n\) induce isomorphisms \(K_{n+1}\otimes^{\mathbf L}_{(\mathcal O_E/\pi_E^{n+1})_X}(\mathcal O_E/\pi_E^{n})_X\simeq K_n\), by associativity of the derived tensor product. We call \((K_n)\), with these isomorphisms, the \emph{tower of reductions} of \(K\).

% object: lemma-reduction-derived-nakayama (blueprint/chap-The-pro-etale-site/sec-Extension-by-zero-and-completion/lemma-reduction-derived-nakayama/_section.md)
\phantomsection\label{bp188}%
\begin{bpstatementit}{Lemma~2.62}{ (Reduction and derived Nakayama)}
Let \(K\in D(X_{\mathrm{pro}\acute{e}t},\mathcal O_E)\). Then \(K_n\) is the cone of \(\pi_E^n:K\to K\). If \(K\) is derived complete and \(K_1=0\), then \(K=0\). Consequently a morphism between derived complete complexes is an isomorphism whenever the induced map on \(K_1\) is an isomorphism.
\end{bpstatementit}

% proof: proof/lemma-reduction-derived-nakayama
\begin{proof}[{Proof}]\phantomsection\label{bp189}%
The first assertion follows from the two-term resolution of \((\mathcal O_E/\pi_E^n)_X\) in Corollary \hyperref[bp159]{2.52}. Tensoring \(K\) with the short exact sequence \(0\to\pi_E^n\mathcal O_E/\pi_E^{n+1}\to\mathcal O_E/\pi_E^{n+1}\to
\mathcal O_E/\pi_E^n\to0\), whose first term is isomorphic to \(\mathcal O_E/\pi_E\), gives a triangle \(K_1\to K_{n+1}\to K_n\). By induction \(K_1=0\) implies \(K_n=0\) for all \(n\). Then the derived limit of the \(K_n\) vanishes, and \(K=0\) by completeness.

For the last assertion, let \(u:K\to L\) be a morphism of derived complete complexes and \(C\) its cone. Derived completeness is preserved by cones, since both sides of the completion map are exact functors of \(K\), so \(C\) is derived complete. Reduction modulo \(\pi_E\) is exact, so \(C_1\) is the cone of \(u_1:K_1\to L_1\), which vanishes when \(u_1\) is an isomorphism. Hence \(C=0\), and \(u\) is an isomorphism.
\end{proof}

% text: prose-02-63
\phantomsection\label{bp190}%
Over a point, for \(K=\mathcal O_E\) placed in degree \(0\), the reductions are \(K_n=\mathcal O_E/\pi_E^n\), and the lemma is the classical Nakayama lemma for the complete local ring \(\mathcal O_E\). The hypothesis of completeness cannot be dropped, since \(E\neq0\) but \(E/\pi_E=0\).

Let \(K\) be derived complete. Since \(R\Gamma(X_{\mathrm{pro}\acute{e}t},-)\) is a right adjoint, it commutes with the derived limit, and therefore

\[
R\Gamma(X_{\mathrm{pro}\acute{e}t},K)
\simeq R\varprojlim_nR\Gamma(X_{\mathrm{pro}\acute{e}t},K_n).
\]

When the \(K_n\) are classical bounded-below complexes, each term on the right can be computed on the étale site by Theorem \hyperref[bp101]{2.35}. In this way we compute the cohomology of \(K\) from its tower of reductions \((K_n)\).

In Section \bpout{3.2} we define constructible complexes over \(\mathcal O_E\), and we show that for Noetherian \(X\) they are exactly the derived complete complexes \(K\) whose reduction \(K_1\) is constructible (Proposition \bpout{3.11}). Derived complete complexes can equally be studied over the constant sheaf with value the discrete ring \(\mathcal O_E\) (Lemma lemma-integral-models).

The rational coefficient sheaves are obtained by inverting \(\ell\) and passing to the colimit over \(E\): \(E_X\simeq\mathcal O_{E,X}[1/\ell]\) and \(\QbarX\simeq\varinjlim_E E_X\) (Lemma \hyperref[bp156]{2.51}). Chapter \bpout{3} also compares constructible complexes over \(E_X\) and \(\QbarX\) with the classical adic categories. It does so by lattices and by descent to a finite extension.



\clearpage\phantomsection\hypertarget{bp-outside}{}%
\section*{Outside this export}
\addcontentsline{toc}{section}{Outside this export}
References marked \ensuremath{{}^\ast} are to parts of the blueprint this document leaves out:
\begin{description}
\item[Section 1.2] Weil cohomology, in 1 Introduction (\bpslug{sec-{}\allowbreak{}Weil-{}\allowbreak{}cohomology})
\item[Chapter 3] Constructible complexes, in Ⅰ The six operations (\bpslug{chap-{}\allowbreak{}Constructible-{}\allowbreak{}complexes})
\item[Section 3.1] Derived functors, in Ⅰ The six operations (\bpslug{sec-{}\allowbreak{}Derived-{}\allowbreak{}functors})
\item[Section 3.2] Constructibility, in Ⅰ The six operations (\bpslug{sec-{}\allowbreak{}Constructibility})
\item[Definition 3.1] Constructible sheaves, in Ⅰ The six operations (\bpslug{definition-{}\allowbreak{}constructible-{}\allowbreak{}sheaves})
\item[Proposition 3.11] The integral constructibility criterion, in Ⅰ The six operations (\bpslug{proposition-{}\allowbreak{}integral-{}\allowbreak{}constructibility-{}\allowbreak{}criterion})
\item[Lemma 3.23] Descent to a finite coefficient field, in Ⅰ The six operations (\bpslug{lemma-{}\allowbreak{}descent-{}\allowbreak{}finite-{}\allowbreak{}coefficient-{}\allowbreak{}field})
\item[Chapter 4] Proper base change, in Ⅰ The six operations (\bpslug{chap-{}\allowbreak{}Proper-{}\allowbreak{}base-{}\allowbreak{}change})
\item[Proposition 4.40] Proper direct image of constructible complexes, in Ⅰ The six operations (\bpslug{proposition-{}\allowbreak{}proper-{}\allowbreak{}direct-{}\allowbreak{}image-{}\allowbreak{}constructible-{}\allowbreak{}complexes})
\item[Chapter 5] Construction of $f^!$, in Ⅰ The six operations (\bpslug{chap-{}\allowbreak{}Construction-{}\allowbreak{}of-{}\allowbreak{}f})
\item[Lemma 5.6] Uniform Ext bound, in Ⅰ The six operations (\bpslug{lemma-{}\allowbreak{}uniform-{}\allowbreak{}ext-{}\allowbreak{}bound})
\item[Lemma 5.17] A uniform bound for proper direct image, in Ⅰ The six operations (\bpslug{lemma-{}\allowbreak{}uniform-{}\allowbreak{}bound-{}\allowbreak{}proper-{}\allowbreak{}direct-{}\allowbreak{}image})
\item[Proposition 6.7] The adjunction for the projective line, in Ⅰ The six operations (\bpslug{proposition-{}\allowbreak{}adjunction-{}\allowbreak{}projective-{}\allowbreak{}line})
\item[Theorem 9.3] Smooth purity for all admissible coefficients, in Ⅱ Duality and the trace formula (\bpslug{theorem-{}\allowbreak{}smooth-{}\allowbreak{}purity-{}\allowbreak{}all-{}\allowbreak{}admissible-{}\allowbreak{}coefficients})
\item[Proposition 10.20] Purity and trace commute with Frobenius, in Ⅱ Duality and the trace formula (\bpslug{proposition-{}\allowbreak{}purity-{}\allowbreak{}trace-{}\allowbreak{}commute-{}\allowbreak{}frobenius})
\end{description}
\clearpage\phantomsection%
\section*{References}
\addcontentsline{toc}{section}{References}
\bpbibitem{BBD82}{bp192}{A. Beilinson, J. Bernstein, and P. Deligne. \emph{Faisceaux pervers}, volume 100 of \emph{Astérisque}. 1982.}

\bpbibitem{BS15}{bp191}{B. Bhatt and P. Scholze. The pro-étale topology for schemes. \emph{Astérisque}, 369:99–201, 2015. \href{https://people.mpim-bonn.mpg.de/scholze/proetale.pdf}{Author's manuscript}.}

\bpbibitem{Con07}{bp193}{B. Conrad. Deligne's notes on Nagata compactifications. \emph{Journal of the Ramanujan Mathematical Society}, 22:205–257, 2007. \href{https://math.stanford.edu/\~conrad/papers/nagatafinal.pdf}{Author's manuscript}.}

\bpbibitem{Del74}{bp194}{P. Deligne. La conjecture de Weil. I. \emph{Publ. Math. IHÉS}, 43:273–307, 1974.}

\bpbibitem{Del80}{bp195}{P. Deligne. La conjecture de Weil. II. \emph{Publ. Math. IHÉS}, 52:137–252, 1980.}

\bpbibitem{Den22}{bp196}{C. Deninger. There is no “Weil-”cohomology theory with real coefficients for arithmetic curves. \emph{Ann. Sc. Norm. Super. Pisa Cl. Sci.}, 2022. To appear. arXiv:\href{https://arxiv.org/abs/2204.02714}{2204.02714}.}

\bpbibitem{Deu41}{bp197}{M. Deuring. Die Typen der Multiplikatorenringe elliptischer Funktionenkörper. \emph{Abh. Math. Sem. Univ. Hamburg}, 14:197–272, 1941.}

\bpbibitem{Eke90}{bp198}{T. Ekedahl. On the adic formalism. In \emph{The Grothendieck Festschrift, Vol. II}, volume 87 of \emph{Progress in Mathematics}, pages 197–218. Birkhäuser, 1990.}

\bpbibitem{Has36}{bp199}{H. Hasse. Zur Theorie der abstrakten elliptischen Funktionenkörper. I, II, III. \emph{J. Reine Angew. Math.}, 175:55–62, 69–88, 193–208, 1936.}

\bpbibitem{Hum72}{bp200}{J. E. Humphreys. \emph{Introduction to Lie Algebras and Representation Theory}, volume 9 of \emph{Graduate Texts in Mathematics}. Springer, 1972.}

\bpbibitem{Ill94}{bp201}{L. Illusie. Autour du théorème de monodromie locale. \emph{Astérisque}, 223:9–57, 1994. §4. \href{https://www.numdam.org/item/AST_1994__223__9_0/}{Original article}.}

\bpbibitem{Kat88}{bp202}{N. M. Katz. Travaux de Laumon. \emph{Astérisque}, 161–162:105–132, 1988. Séminaire Bourbaki, exp. 691.}

\bpbibitem{KW01}{bp203}{R. Kiehl and R. Weissauer. \emph{Weil conjectures, perverse sheaves and \(\ell\)-adic Fourier transform}, volume 42 of \emph{Ergebnisse der Mathematik}. Springer, 2001.}

\bpbibitem{Lau87}{bp204}{G. Laumon. Transformation de Fourier, constantes d'équations fonctionnelles et conjecture de Weil. \emph{Publ. Math. IHÉS}, 65:131–210, 1987.}

\bpbibitem{Mil86a}{bp205}{J. S. Milne. \href{https://www.jmilne.org/math/xnotes/AVs.pdf}{Abelian Varieties}. In \emph{Arithmetic Geometry}. 1986. Corrected online version.}

\bpbibitem{Mil86b}{bp206}{J. S. Milne. \href{https://www.jmilne.org/math/xnotes/JVs.pdf}{Jacobian Varieties}. In \emph{Arithmetic Geometry}. 1986. Corrected online version.}

\bpbibitem{Ran39}{bp207}{R. A. Rankin. Contributions to the theory of Ramanujan's function \(\tau(n)\) and similar arithmetical functions. II. \emph{Proc. Cambridge Philos. Soc.}, 35:357–372, 1939.}

\bpbibitem{SGA4}{bp208}{M. Artin, A. Grothendieck, and J.-L. Verdier, editors. \emph{Théorie des topos et cohomologie étale des schémas}, volumes 269, 270, 305 of \emph{Lecture Notes in Mathematics}. Springer, 1972–1973. SGA 4.}

\bpbibitem{SGA4½}{bp209}{P. Deligne. \emph{Cohomologie étale}, volume 569 of \emph{Lecture Notes in Mathematics}. Springer, 1977. SGA 4½; \emph{Théorèmes de finitude}, §3, and \emph{Rapport sur la formule des traces}, §3, and \emph{Sommes trigonométriques}, §1.}

\bpbibitem{SGA5}{bp210}{A. Grothendieck, editor. \emph{Cohomologie \(\ell\)-adique et fonctions \(L\)}, volume 589 of \emph{Lecture Notes in Mathematics}. Springer, 1977. SGA 5.}

\bpbibitem{SGA7II}{bp211}{P. Deligne and N. Katz, editors. \emph{Groupes de monodromie en géométrie algébrique}, volume 340 of \emph{Lecture Notes in Mathematics}. Springer, 1973. SGA 7 II; Exposé XIII, (2.1.7).}

\bpbibitem{Stacks}{bp212}{The Stacks Project Authors. \href{https://stacks.math.columbia.edu}{The Stacks Project}. Individual results are cited by their stable tags in the text.}

\bpbibitem{Var05}{bp213}{Y. Varshavsky. Lefschetz–Verdier trace formula and a generalization of a theorem of Fujiwara. 2005. Version 2. arXiv:\href{https://arxiv.org/abs/math/0505564}{math/0505564}. \href{https://arxiv.org/abs/math/0505564v2}{Author's preprint}.}

\bpbibitem{Wei48a}{bp214}{A. Weil. \emph{Sur les courbes algébriques et les variétés qui s'en déduisent}. Hermann, Paris, 1948.}

\bpbibitem{Wei48b}{bp215}{A. Weil. \emph{Variétés abéliennes et courbes algébriques}. Hermann, Paris, 1948.}

\bpbibitem{Wei49}{bp216}{A. Weil. Numbers of solutions of equations in finite fields. \emph{Bull. Amer. Math. Soc.}, 55:497–508, 1949.}

\end{document}
